\documentclass[11pt]{article}
\usepackage[T1]{fontenc}
\usepackage[utf8]{inputenc}
\usepackage{lmodern,microtype,amsmath,amssymb,amsthm,mathtools,booktabs,array}
\usepackage[a4paper,margin=28mm]{geometry}
\usepackage[colorlinks=true,linkcolor=blue,citecolor=blue,urlcolor=blue,breaklinks=true]{hyperref}
\usepackage{xurl}
\setlength{\emergencystretch}{2em}
\hypersetup{pdftitle={A counterexample and a dimension formula for Grothendieck's strong transcendence conjecture},pdfauthor={},pdfsubject={Motivic period comparisons and conjugate Hodge filtrations},pdfkeywords={period conjecture, motives, Hodge filtration, integral Betti comparison}}
\newtheorem{theorem}{Theorem}[section]
\newtheorem{proposition}[theorem]{Proposition}
\newtheorem{lemma}[theorem]{Lemma}
\newtheorem{corollary}[theorem]{Corollary}
\theoremstyle{definition}\newtheorem{definition}[theorem]{Definition}
\theoremstyle{remark}\newtheorem{remark}[theorem]{Remark}
\newcommand{\Q}{\mathbf Q}\newcommand{\Z}{\mathbf Z}
\newcommand{\R}{\mathbf R}\newcommand{\C}{\mathbf C}
\newcommand{\Qb}{\overline{\mathbf Q}}\newcommand{\Fp}{\mathbf F_p}
\newcommand{\Gm}{\mathbf G_m}\newcommand{\B}{\mathrm B}\newcommand{\dR}{\mathrm{dR}}
\DeclareMathOperator{\Aut}{Aut}\DeclareMathOperator{\Isom}{Isom}
\DeclareMathOperator{\Lie}{Lie}\DeclareMathOperator{\trdeg}{trdeg}
\DeclareMathOperator{\gr}{gr}\DeclareMathOperator{\GL}{GL}
\DeclareMathOperator{\SL}{SL}\DeclareMathOperator{\End}{End}
\DeclareMathOperator{\Gal}{Gal}\DeclareMathOperator{\im}{im}
\DeclareMathOperator{\Spec}{Spec}\DeclareMathOperator{\Res}{Res}
\newcommand{\workbenchurl}{https://github.com/KokunoYumeto/algebraic-cycles-and-motives-workbench}
\newcommand{\recordurl}{https://github.com/KokunoYumeto/mathematics-workbenches}
\title{A counterexample and a dimension formula for\newline Grothendieck's strong transcendence conjecture}
\author{}
\date{23 November 2026}
\begin{document}
\maketitle
\begin{abstract}
Grothendieck's Montpellier notes propose a strong transcendence conjecture for the comparison between Betti, de Rham and Hodge realizations. We show that five of its predicted degrees are too large for a non-CM elliptic curve defined over the rationals. For a pure motivated motive and its complex conjugate, we compute the six algebraic image dimensions of the comparison triangle. One normal subgroup of their joint motivic Galois group measures the missing freedom. These dimensions are unconditional upper bounds for the corresponding transcendence degrees, and become equalities under the period conjecture for the joint motive. An explicit elliptic analysis describes the splitting parameter, its CM examples, its zero locus and its extremal growth normalization. We also explain three limitations of comparison data: ordinary Hodge--Tate filtrations do not detect complex multiplication of elliptic curves, polarizable rational Hodge structures with a real involution need not be geometric, and rational or locally integral comparisons need not assemble to an integral comparison. The mathematics and text were produced and checked by GPT-6 Astra in ChatGPT Pro mode; no human has checked them.
\end{abstract}

\section{Grothendieck's question and the answer}

This paper determines the algebraic constraint missing from the strong transcendence prediction in Grothendieck's comparison triangle. On handwritten page 10 of Montpellier cote 18 (PDF page 11), Grothendieck writes ``Conj. de transcendance forte'', followed by a list of expected transcendence degrees and the question ``alors ?'' \cite[p.~10; PDF p.~11]{Grothendieck}. His three realizations are linked by the ordinary comparison isomorphism, the splitting provided by the Hodge decomposition, and their inverse composite. The question is whether the ordinary comparison and its canonical Hodge splitting have as many algebraically independent coordinates as their separate parameter spaces permit.

There is a simple obstruction. A motive over a real algebraic field has a rational real Frobenius. It recovers the conjugate Hodge filtration from the original filtration. Consequently it also recovers the Hodge splitting from the ordinary periods, so that adjoining that splitting contributes no transcendence. This contradicts the prediction whenever its splitting torsor has positive dimension. The obstruction persists, in a measured form, when the motive has no real model: its periods and conjugate periods can share motivic tensors.

Here is the complete dimension statement. Fix an algebraic subfield $k\subset\C$, with its specified embedding. Let $M$ be a finite direct sum of pure motivated motives over $k$, with rational coefficients. We use Andr\'e's semisimple graded Tannakian category \cite[Theorem 0.4, p.~8; \S4.6, pp.~24--25]{Andre}. All tensor functors below are restricted to the tensor category generated by the indicated object. Write
\[
 V=T_{\B},\qquad D=T_{\dR},\qquad H=\gr_FD.
\]
The weight-$n$ summand $\gr_F^pD$ has Hodge bidegree $(p,n-p)$. The comparison triangle consists of
\[
 \alpha:V_{\C}\longrightarrow D_{\C},\qquad
 \beta:D_{\C}\longrightarrow H_{\C},\qquad
 \gamma=(\beta\alpha)^{-1}.
\]
The map $\beta$ sends each Hodge summand isomorphically to its associated-graded summand; in particular $\gr_F\beta=1$. Let $\phi$ be the Hodge filtration on $V_{\C}$ and $\psi$ its Hodge bigrading. Coordinates of tensor isomorphisms are taken in algebraic bases. Coordinates of filtrations and bigradings are taken in their algebraic parameter spaces; algebraic extensions of the field of definition do not affect any degree below. Put $K=k(\alpha,\beta)$.

Let $G=\Aut^{\otimes}(V)$ and let $\mu$ and $w$ denote its Hodge and central weight cocharacters. After geometric identification set
\[
 B=P_G(\mu),\quad B^-=P_G(w\mu^{-1}),\quad
 L=Z_G(\mu),\quad U=R_u(B),
\]
and write $g=\dim G$, $b=\dim B$, $l=\dim L$, $u=\dim U$. Thus
\[
                 g=l+2u,\qquad b=l+u.
\]
Complex conjugation is denoted by $c$. The motive
$N=M\oplus c^*M$ is defined over $k\,c(k)$, and may be viewed over $\Qb$. Conjugation identifies the Betti realizations of its two summands. The joint group is then a subgroup
\[
                   J\subset G\times G
\]
with surjective projections. Let $N_2$ be the kernel of the first projection, viewed as a normal subgroup of the second $G$, and define
\[
             v=\dim N_2-\dim(N_2\cap B^-).
\]
All dimensions concern identity components; no connectedness hypothesis on the motivic groups is imposed.

\begin{theorem}\label{thm:main}
With these hypotheses, $0\leq v\leq u$. On the component of the comparison torsor of $N$ containing its comparison point, the maps defining $\alpha,\beta,\gamma,\phi,\psi$ have image-closure dimensions
\[
\begin{array}{c|rrrrrr}
\text{map}&\alpha&\beta&\gamma&\phi&\psi&(\alpha,\beta)\\ \hline
\text{dimension}&g&u&b+v&u&u+v&g+v.
\end{array}
\]
The corresponding transcendence degrees are bounded above by these numbers. The complete table, including relative degrees, is
\begin{center}
\begin{tabular}{@{}lcc@{}}
\toprule
Field or extension&Upper bound&Strong prediction\\
\midrule
$k(\alpha)/k$&$g$&$g$\\
$k(\beta)/k$&$u$&$u$\\
$k(\gamma)/k$&$b+v$&$g$\\
$\Q(\phi)/\Q$&$u$&$u$\\
$\Q(\psi)/\Q$&$u+v$&$2u$\\
$K/k$&$g+v$&$g+u$\\
$K/k(\alpha)$&$v$&$u$\\
$K/k(\beta)$&$b+v$&$g$\\
$K/k(\gamma)$&$u$&$u$\\
\bottomrule
\end{tabular}
\end{center}
Every upper bound is an equality if the comparison point of $N$ is Zariski generic in its motivated comparison-torsor component. The strong prediction therefore follows from this period hypothesis exactly when $v=u$. If $v<u$, its five larger entries are false unconditionally.
\end{theorem}

The last assertion compares numerical predictions, not the truth values of two unproved arithmetic conjectures. In particular, this theorem does not prove the ordinary Grothendieck period conjecture. The hypothesis we use is its usual genericity assertion for the named joint motive. The algebraic-cycle version of the period conjecture implies it: Bost and Charles's Corollary 2.11 and Conjecture 2.12 \cite[pp.~11--12]{BostCharles} put the motivated torsor between the comparison closure and the algebraic-cycle torsor.

\begin{corollary}\label{cor:real}
If $M\simeq c^*M$ in the motivated category over $\Qb$, then $v=0$. If $M$ has a model over a real algebraic field, then
\[
                 K=k(\alpha),\qquad \Q(\psi)=\Q(\phi)
\]
after extending the field of definition of that model when necessary. In particular, for
\[
               M=h^1(E),\qquad E:y^2=x^3-x+1
\]
over $\Q$, one has $(g,b,u,v)=(4,3,1,0)$. The five strong predictions concerning $\gamma,\psi,K,K/k(\alpha)$ and $K/k(\beta)$ fail for this motive.
\end{corollary}

The correction matters because a Hodge splitting depends on both a filtration and its conjugate. Treating those two filtrations as independently variable overestimates the dimension unless the common motivic quotient carries only central Hodge weights. Section~\ref{sec:images} makes this obstruction precise by two stabilizer computations. Reductivity supplies opposite pairs of Hodge weights; purity supplies an opposed pair of filtrations; normality of $N_2$ supplies uniform fibre dimensions, including at nongeneric points.

Further assertions clarify the scope of comparison arguments in the same notes.
\begin{theorem}\label{thm:other}
The following statements hold.
\begin{enumerate}
\item The non-CM elliptic curve $y^2=4x^3-4x-1$ has ordinary good reduction at $5$. Over $\mathcal O_{\C_5}$ its $5$-divisible group is isomorphic to $\mu_{5^\infty}\oplus\Q_5/\Z_5$, and its Hodge--Tate line is defined over $\Q_5$.
\item If $k$ is perfect of characteristic $p>3$ and $E:y^2=P(x)=4x^3-g_2x-g_3$ is smooth, put
\[
 H_p=[x^{p-1}]P(x)^{(p-1)/2},\quad
 K_p=[x^{p-2}]P(x)^{(p-1)/2}.
\]
For $a\in k$, on $k(E)(t)$ the derivations
\[
 D_a=y\partial_x+(6x^2-g_2/2)\partial_y-(x+a)\partial_t,
 \qquad V=-\partial_t
\]
satisfy $D_a^p=H_pD_a-(K_p+aH_p)V$. The line $kD_a$ is restricted-stable exactly when $K_p+aH_p=0$.
\item There is a fully specified polarizable rational Hodge structure of type $(3,0)+(0,3)$, with a rational anti-Hodge involution, which is not the realization of any pure motive defined by algebraic correspondences over $\C$.
\item There are two complex embeddings of $\overline{\Qb(t)}$, fixing $\Qb$ and not related by conjugation, whose integral Betti correspondence diagrams are naturally isomorphic. There is also another pair, agreeing already on $\Qb(t)$, for which every finite subdiagram has a comparison but the whole abelian-variety and isogeny diagram has none.
\item For $\mathcal O=\Z[\sqrt{-5}]$ and $I=(2,1+\sqrt{-5})$, the CM elliptic curves $\C/\mathcal O$ and $\C/I$ occur as conjugate realizations of an algebraic elliptic curve with its labelled endomorphisms. Their integral Betti modules are locally isomorphic and rationally isomorphic, but are not isomorphic as $\mathcal O$-modules.
\end{enumerate}
\end{theorem}

The first statement is an application of the ordinary case already described by Howe \cite[Theorem 1.2 and Remark 1.3, p.~673]{Howe}; the point here is the distinction between endomorphisms of the $p$-divisible group after base change and endomorphisms of the elliptic curve. The second is a coordinate form of Cartier's identity \cite[Proposition 1.2.4, equation (1.2.4.5), p.~24]{Ogus}. The third combines the rigidity of Hodge structures with gaps, emphasized by Totaro \cite[pp.~1--2]{Totaro}, and Yoshinaga's computability theorem for periods \cite[Theorem 18, p.~11]{Yoshinaga}. The last statement is an ideal-class obstruction, together with the fundamental theorem of complex multiplication \cite[Theorem 3.13, pp.~23--24]{Milne}. These antecedents are ingredients and are not asserted as new theorems here.

The elliptic analysis in Section~\ref{sec:elliptic-splitting} gives a complementary description of the canonical splitting: an explicit CM example has parameter $1/2$, its vanishing locus is exactly $j=0,1728$, and the Hodge parameter uniquely minimizes the quadratic growth of a normalized sigma function. These facts are compatible with an open arithmetic question: they do not prove that every algebraic splitting parameter comes from CM. The arithmetic tests in Sections~\ref{sec:elliptic-arithmetic}--\ref{sec:elliptic-equations} separate algebraicity of this scalar from algebraic compact points, restricted powers, crystalline Frobenius fixity, bounded denominators and ordinary overconvergent corrections. The local Frobenius criterion is proved using the precise deformation-theoretic lifting theorem; explicit CM examples demonstrate denominator failures. Known results on unitary connections, analytic subgroups and horizontal integrality are credited at the point of use. The later sections also distinguish abstract real Hodge pairs from geometric ones, show precisely what adding labelled Hodge filtrations recovers, and prove the cycle, Tate and integral Lefschetz normalization statements used in these comparisons.

The same universal extension gives an analytic quotient onto the conjugate elliptic curve whose actual differential scale cannot be algebraic together with the splitting parameter. We classify the CM cases among the associated second-kind lattices and give a rank bound for algebraic points on the actual quotient graph. Herblot's quantitative theorem yields a restricted-Lie density bound with the field degree retained. The first prime-degree denominator then determines the exact local geometric size at the complementary primes and proves positive arithmetic coefficient growth. These results, together with the ordinary isogeny denominator formula and the two-jet obstruction for the elliptic period product, explain why the cited algebraicity criteria do not by themselves settle the non-CM splitting question.

The real and adelic comparisons lead to further exact distinctions. Real abelian absolute-Hodge motives are fully faithfully realized as the indicated Hodge pairs, and twisting the Betti functor by real Frobenius has different answers for Artin, Tate and CM objects. Relative tensor frames have dimension $(n-1)\dim G$; compact-open adelic images attain this bound after suitable changes of places, using the known elliptic open-image theorems. Finally, the actual automorphism-cocycle image has the same algebraic locus as the comparison point, while passing to its generated subgroup can enlarge the closure. A CM period example over a transcendental field realizes that strict enlargement, including a degree-one augmentation with minimum field transcendence degree one. A Tate-square example separately refutes universal openness of the combined finite and complex image.

\subsection*{Relation to period theorems}
The Tannakian framework is due to Grothendieck, Saavedra, Deligne and Milne; Andr\'e supplies the unconditional semisimple category used here. Kontsevich and Zagier formulate relations among periods through algebraic operations on integrals \cite[\S1.2 and \S4]{KontsevichZagier}. Huber and M\"uller-Stach relate this period algebra to Nori's category \cite[Theorem 13.1.4, pp.~284--285]{HuberMuellerStach}. The question studied here keeps the canonical Hodge splitter tied to the joint comparison point throughout the dimension calculation.

Brown and Dupont's single-valued comparison transports real Frobenius to de Rham cohomology \cite[Definition 2.6 and Proposition 2.9, pp.~10--11]{BrownDupont}. That established construction is the matrix $\alpha f_\infty\alpha^{-1}$ used in the real-model proof below. Our dimension question additionally retains the Hodge filtration and, for a general complex embedding, the joint motivic relations with the conjugate motive. We make no priority claim for the ingredients or for their combination.

The functional period theorem of Bakker and Tsimerman \cite[Theorem 1.3; Theorem 4.25]{BakkerTsimerman} concerns periods varying in algebraic families. It supplies a geometric analogue of period genericity, not the numerical genericity assumption for a fixed $M$ in Theorem~\ref{thm:main}. Ayoub identifies the motivic Galois group with the automorphism object of a functor of geometric local systems for a chosen embedding \cite[Theorem 4.4.2, p.~234]{Ayoub}. The integral questions in Section~\ref{sec:integral} ask instead for isomorphisms between two embeddings with specified lattices; their existence is additional data. The recent machine-assisted Montpellier edition \cite{Hua} also separates modernized readings of the manuscripts from what its formal proofs establish. The explicit orbit calculation below is independent of any assertion that a formalized modern reading exhausts the handwritten conjecture.

\subsection*{Printed and dated formulations}
The complex period question has a precise printed antecedent. In the 1966 paper on algebraic de Rham cohomology, Grothendieck's note 10 discusses the four elliptic periods and quasi-periods, calling the displayed Legendre identity ``The only known general algebraic relation'' and asking for algebraic independence in the non-CM case \cite[note 10, p.~102; PDF p.~9]{GrothendieckDeRham}. The sign of its quasi-period convention is opposite to that of $\eta_i$ in Section~\ref{sec:elliptic-splitting}, because that note integrates $x\,dx/y$ itself. Note 13, explicitly marked July 1965, proposes the canonical integrable connection on relative de Rham cohomology \cite[p.~103; PDF p.~10]{GrothendieckDeRham}.

The Tate-module reconstruction question is present in a May 1966 letter to Tate. Section 3.2 asks ``quel est le module de Tate correspondant'' to the generic fibre from Dieudonne data and a lifted filtration \cite[p.~30, \S3.2]{GrothendieckTateLetter}. Its immediate setup retains the reduction condition in \S3.1: an extension of an ind-etale group by a toric $p$-divisible group; its final sentence asks about more general groups. This is a dated written question, not a theorem asserting the later crystalline equivalence.

In the Nice address, Remarks 2--3 ask for both arrows between the filtered Dieudonne object and the generic-fibre Galois representation, including how to ``reconstituer plus ou moins alg\'ebriquement le couple'' from that representation \cite[p.~435; PDF p.~5]{GrothendieckNice}. The final paragraph extends the question to longer filtrations and higher-dimensional bases. These dated and printed passages establish the questions in those documents; no relative date for an individual Montpellier leaf, and no claim of earliest priority, is inferred.

The local comparison section also distinguishes arbitrary generic-fibre $p$-divisible groups from good-reduction objects, and records the crystalline, semistable and integral reconstruction theorems with their exact hypotheses. Two elementary obstructions complement those known theorems. Under both global centralizations and compatibility, the equalizer of two comparison sections is normal. Nevertheless, a compatible inverse system of good-reduction Artin--Tate realizations can have a Galois image that is not a $p$-adic analytic group. The finite-dimensional hypothesis is essential.

\section{The comparison triangle as algebraic geometry}
This section constructs the parameter spaces and explains why complex conjugation must be included in the algebraic object controlling the splitting.

\subsection{Normalized splittings}
Let $\mathcal C$ be the tensor category generated by $M$. Strictness of morphisms of pure Hodge structures makes $H=\gr_FD$ an exact tensor fibre functor. Let
\[
 P=\Isom^{\otimes}(V_k,D),\qquad
 P''=\Isom^{\otimes}(H,V_k).
\]
Let $Q$ be the space of tensor isomorphisms $s:D\to H$ preserving the filtration and inducing the identity on its associated graded. In the de Rham form of the group, write
\[
 B_D=\Aut^{\otimes}(D,F),\qquad
 U_D=\ker(B_D\to\Aut^{\otimes}(H)).
\]
Then $Q$ is a $U_D$-torsor. Indeed, for $s_1,s_2\in Q$, the quotient $s_2^{-1}s_1$ is in $U_D$, uniquely, and every element of $U_D$ gives another splitting by composition. Existence after algebraic closure follows by choosing a Levi subgroup of the parabolic $B_D$. Thus $\dim Q=u$.

Over an algebraic closure, the spaces of filtrations and bigradings of the given tensor type are $G/B$ and $G/L$. If the type is not defined over $\Q$, one first works over its number field of definition, or takes the finite union of its conjugates. These two choices have the same transcendence dimensions. The morphism
\[
                       G/L\longrightarrow G/B
\]
has fibre $B/L$, a torsor under the unipotent radical of the stabilizer of the chosen flag. The unipotent group varies with the flag; no fixed subgroup of $G$ is required to act on every fibre.

The Lie algebra of $G$ decomposes into $\mu$-weight spaces. Its nonzero weights occur in pairs of opposite weights because $G^0$ is reductive. The zero-weight subspace is $\Lie L$. Choosing the parabolic convention for which nonnegative weights give $B$, the positive-weight spaces form $\Lie U$. Consequently $g=l+2u$, $b=l+u$, and $\dim(G/B)=u$. Components of $G$ are finite in number and do not change these dimensions.

\begin{remark}[Normalization and the missing reverse arrow]\label{rem:normalized-torsor}
If one instead allows a splitting to induce an element of a nontrivial central one-dimensional torus $T_0$ on the associated graded, its space $Q_H$ maps to $T_0$ with fibre $Q$ over the identity. Its group is the inverse image of $T_0$ in $B_D$, and after choosing a splitting $Q_H\simeq Q\times T_0$. Thus its dimension is $u+1$, whereas the canonical $\beta$ belongs to the normalized fibre of dimension $u$. When a cocharacter has finite kernel, $T_0$ means its image torus.

The map $G/L\to G/B$ has no $G$-equivariant section unless $u=0$: a section would select a $B$-fixed point of $B/L$, but $U$ acts simply transitively on that fibre. For $\GL_2$ and a line filtration even a nonequivariant regular section is impossible. Here the map is $(\mathbf P^1\times\mathbf P^1)\setminus\Delta\to\mathbf P^1$. A section would give a fixed-point-free morphism $f:\mathbf P^1\to\mathbf P^1$. If $f=[P:Q]$ has positive degree, the homogeneous equation $XQ-YP=0$ has a projective zero over the algebraic closure (or vanishes identically). A constant map fixes its own value. Forgetting a grading therefore has no general algebraic inverse on these parameter spaces.
\end{remark}

\begin{lemma}\label{lem:tate-membership}
If $\Q(-n)$ belongs to the tensor category generated by $M$, then $(2\pi i)^n\in k(\alpha)$.
\end{lemma}
\begin{proof}
For $\Q(-1)=H^2(\mathbf P^1)$, the positive integral Betti generator maps to $(2\pi i)^{-1}$ times the residue-normalized de Rham generator. Hence the comparison coefficient on $\Q(-n)$ is $(2\pi i)^{-n}$. Every comparison in a tensor construction, dual or algebraic summand is a rational function of the entries of $\alpha$: tensor products multiply entries, duals invert matrices, and summands use algebraic projectors. This coefficient lies in $k(\alpha)$, and inversion gives the assertion. Proposition~\ref{prop:cycle-normalization} gives the same normalization for all cycle classes.
\end{proof}

\subsection{Algebraic extension of the base field}
We may pass from $k$ to $\Qb$ for dimension and transcendence-degree calculations. The reason for dimensions deserves mention. A motivated tensor defined over an algebraic extension is defined over a finite extension. After a finite Galois enlargement, its finitely many conjugates form a finite set in a rational Betti tensor construction. Their symmetric tensors descend, using traces and products of the algebraic cycles in the motivated construction. The group over $k$ therefore permutes that finite set. Its identity component fixes each tensor. The group over $\Qb$ is a subgroup of the group over $k$ and has the same identity component. The dimensions of their Hodge-weight spaces coincide.

For a coordinate tuple $z$, the extension $\Qb(z)/k(z)$ is algebraic. Transcendence degree is invariant under algebraic extension of either field in the relevant tower. These observations justify all base extensions used below without assuming that the original group is connected.

\subsection{The joint torsor}
Conjugation gives a rational tensor isomorphism $j_{\B}:V\to V_c$ between the Betti realizations of $M$ and $c^*M$, and a conjugate-semilinear map $j_{\dR}:D\to D_c$. If $\kappa_V$ denotes coefficient conjugation on $V_{\C}$, their comparisons satisfy
\[
                  \alpha_c j_{\B}=j_{\dR}\alpha\kappa_V.
\]
Both sides are linear: the right side has two semilinear factors. This identity follows by integrating conjugate forms over conjugate cycles. In conjugate bases the matrix of $\alpha_c$ is $\overline\alpha$.

Let $T$ be the comparison torsor of $N=M\oplus c^*M$, with projection maps $a:T\to P$ and $d:T\to P_c$. Define two Betti flags on $T$ by
\[
 F_1(t)=a(t)^{-1}F_D,\qquad
 F_2(t)=j_{\B}^{-1}d(t)^{-1}F_{D_c}.
\]
Let $T^{\mathrm{opp}}$ be the open subset where these flags are opposed in each weight. It is nonempty: at $t_0=(\alpha,\alpha_c)$ the second flag is the conjugate of the first, and a pure Hodge structure is the direct sum of their intersections. A nonempty open variety over $\Qb$ has algebraic points, so algebraic trivializations on this open set exist even though $t_0$ itself need not be algebraic.

For $t\in T^{\mathrm{opp}}$, the intersections
\[
                 F_1^p(t)\cap F_2^{n-p}(t)
\]
give a tensor bigrading. Transport it by $a(t)$ and project to $H$ to obtain $s(t)\in Q$. Put
\[
 \gamma(t)=(s(t)a(t))^{-1},\quad
 \phi(t)=F_1(t),\quad
 \psi(t)=\bigl(F_1^p(t)\cap F_2^{n-p}(t)\bigr)_{p,n}.
\]
These maps are regular on $T^{\mathrm{opp}}$. On a chart where the necessary minors are invertible, intersections of subbundles of constant rank and the projection to the associated graded are expressed by matrix inversion. They specialize at $t_0$ to $\beta,\gamma,\phi,\psi$. In particular
\[
             k(\alpha,\beta)\subset k(\alpha)\,c(k(\alpha))
\]
after extending $k$ to contain $c(k)$.

\section{The six image dimensions}\label{sec:images}
This section proves Theorem~\ref{thm:main} by computing stabilizers before making any assertion about transcendence.

\begin{lemma}\label{lem:parabolic}
Let $n=\dim N_2$ and $n_B=\dim(N_2\cap B^-)$. The joint Hodge cocharacter is $\lambda=(\mu,w\mu^{-1})\in J_{\C}$. Its parabolic is
\[
          B_J=P_J(\lambda)=J\cap(B\times B^-).
\]
The first projection $B_J^0\to B^0$ is surjective and has kernel of dimension $n_B$. Moreover $0\leq n-n_B\leq u$.
\end{lemma}
\begin{proof}
Motivated tensors are Hodge tensors, so the pair of Hodge cocharacters preserves all tensors defining $J$. This proves $\lambda\in J_{\C}$ and the formula for its parabolic. The differential of the surjection $J^0\to G^0$ is surjective and respects cocharacter weights. It is therefore surjective on the sum of the nonnegative-weight spaces. The image of the connected algebraic subgroup $B_J^0$ in $B^0$ is closed and has its full Lie algebra. In characteristic zero it is all of $B^0$. Its kernel is the identity-component contribution from $N_2\cap B^-$, of dimension $n_B$.

The Lie algebra of $N_2$ is stable under the second Hodge cocharacter. Its quotient by the nonnegative-weight subspace for $w\mu^{-1}$ has dimension $n-n_B$. Since $w$ is central, these are some of the positive $\mu$-weight spaces in $\Lie G$. Their total dimension is $u$. This proves both inequalities.
\end{proof}

\begin{proposition}\label{prop:images}
On the component of $T$ containing $t_0$, the closures of the $a,\phi,\psi,(a,s),\gamma,s$ images have dimensions
\[
                      g, u, u+v, g+v, b+v, u.
\]
Under the isomorphism
\[
 P\times Q\longrightarrow Q\times P'',\qquad
 (a,s)\longmapsto\bigl(s,(sa)^{-1}\bigr),
\]
the closure of the $(a,s)$ image is $Q\times\Gamma$, where $\Gamma$ is the closure of the $\gamma$ image in the corresponding component of $P''$.
\end{proposition}
\begin{proof}
We work geometrically and use the component through $t_0$. Since both projections of $J$ are surjective, $\dim J=g+n$. The first projection of the joint torsor is surjective, so its restriction to the nonempty opposed open set has dense image in its torsor component, of dimension $g$. Its flag quotient has stabilizer of dimension $b$, and hence has dimension $g-b=u$.

The pair $(F_1,F_2)$ has stabilizer $B_J$. On the opposed open set, a pair of flags is equivalent to the resulting bigrading: the latter recovers both flags by summing its components. By Lemma~\ref{lem:parabolic}, $\dim B_J=b+n_B$. Thus
\[
            \dim\im\psi=(g+n)-(b+n_B)=u+v.
\]

For the remaining maps use the de Rham form of $J$. Its filtration stabilizers are defined over the algebraic base field. After identifying them over $\C$, set
\[
 S_\gamma=J\cap(U\times B^-),\qquad
 S_K=J\cap(1\times B^-).
\]
Write two points as $t=(a,d)$ and $t'=(xa,yd)$ under the simply transitive de Rham action. If $x\in U$ and $y\in B^-$, both Betti flags are unchanged. Their opposite de Rham flag is transported by $x$, so the normalized splitter changes by $s'=sx^{-1}$. Hence $(s'a')^{-1}=(sa)^{-1}$ and $\gamma' =\gamma$.

Conversely, if $\gamma'=\gamma$, then $s'a'=sa$, so $x=a'a^{-1}=s'^{-1}s\in U$. The common map $\gamma$ transports the fixed grading of $H$ to a common Betti bigrading; therefore it determines the same second flag, which forces $y\in B^-$. The $\gamma$ fibres are exactly the $S_\gamma$-cosets on the open set. Imposing $a'=a$ gives exactly the $S_K$-cosets for the $(a,s)$ fibres.

The surjection $B_J^0\to B^0$ implies that $S_\gamma\to U$ is surjective: take the inverse image of the connected subgroup $U$. Its kernel is $S_K$, and
\[
          \dim S_K=n_B,\qquad \dim S_\gamma=u+n_B.
\]
Subtracting from $g+n$ yields
\[
 \dim\im(a,s)=g+v,\qquad
 \dim\im\gamma=g-u+v=b+v.
\]
On each $\gamma$ fibre, the coordinate $s$ ranges through all of $Q$, because the projection to $U$ is surjective and $U$ acts simply transitively on normalized splittings. The displayed change of coordinates therefore identifies the image with the product of $Q$ and the $\gamma$ image. Taking closures gives $Q\times\Gamma$, and also proves $\dim\im s=u$. Finite components do not alter any fibre dimension or the connected $U$ action.
\end{proof}

\begin{corollary}\label{cor:central}
The image of $(a,s)$ is dense in the corresponding component of $P\times Q$ if and only if $v=u$. Equivalently, the projected Hodge cocharacter is central on the identity component of $G/N_2$.
\end{corollary}
\begin{proof}
The ambient dimension is $g+u$ and the irreducible image closure has dimension $g+v$. Equality is equivalent to density. On the reductive quotient Lie algebra $\Lie(G/N_2)$, the positive Hodge-weight subspace has dimension $u-v$. Reductivity pairs positive and negative weights. Consequently $u-v=0$ precisely when every weight is zero, which is precisely centrality of the projected cocharacter on the identity component. Replacing $N_2$ by its identity component only changes a finite quotient and leaves the criterion unchanged.
\end{proof}

\section{Transcendence degrees and the elliptic counterexample}
This section passes from algebraic image dimensions to bounds at the comparison point and proves the counterexample without a period conjecture.

\begin{proof}[Proof of Theorem~\ref{thm:main}]
The Zariski closure of a point over an algebraic field has dimension equal to the transcendence degree of its coordinate field. Proposition~\ref{prop:images} therefore gives all six absolute upper bounds.

Relative bounds require uniform fibres, not only generic fibres. The product $Q\times\Gamma$ proves that adjoining $\gamma$ to $\beta$ contributes at most $b+v$, while adjoining $\beta$ to $\gamma$ contributes at most $u$. The coordinate change in Proposition~\ref{prop:images} and its inverse identify both joint coordinate fields with $K$.

For the remaining bound, we verify that closing the joint image adds no exceptional first-coordinate fibres. Work geometrically on the chosen component and identify it with $J^0$. Write $F^+,F^-$ for the fixed opposed flags in these identifications, and put
\[
 R=\ker(J^0\longrightarrow G^0),\qquad R^0=N_2^0,
 \qquad X^-=G^0/(B^-\cap G^0).
\]
View $R$ in the second factor. It is normal in $G^0$, since the second projection of $J^0$ is surjective. The subgroup $N_2^0\cap B^-$ is parabolic in $N_2^0$: up to central isogeny a connected normal reductive subgroup consists of a central subtorus and semisimple factors, and the cocharacter selects the corresponding parabolics in those factors. Thus $N_2^0\cdot F^-$ is projective. The orbit $R\cdot F^-$ is a finite union of its translates, hence is closed and of pure dimension $v$.

Let $W'$ be the image of
\[
 J^0\longrightarrow G^0\times X^-,\qquad
 (x,y)\longmapsto(x,y^{-1}F^-).
\]
For a lift $(x,y_0)$ of $x$, its fibre is $y_0^{-1}(R\cdot F^-)$. Pulling back along the smooth surjection $J^0\to G^0$ and multiplying the flag by $y_0$ identifies the relative image with $J^0\times(R\cdot F^-)$. It is closed there, so closedness descends to $W'$. This proves global closedness, not merely closedness of individual fibres.

The incidence open set where the second flag is opposed to $x^{-1}F^+$ is isomorphic to $G^0\times U$ by
\[
 (x,s)\longmapsto\bigl(x,(sx)^{-1}F^-\bigr).
\]
Its intersection with $W'$ is exactly the actual $(a,s)$ image. Consequently that image is already closed in the corresponding $P\times Q$ component, and every nonempty first-coordinate fibre is an open subset of a translate of $R\cdot F^-$, of dimension $v$. These closedness and dimension statements descend to the algebraic field of definition. In particular the fibre over the point with coordinate field $k(\alpha)$ has dimension $v$. The residue-field dimension bound gives $\trdeg_{k(\alpha)}K\leq v$, even if $\alpha$ is nongeneric.

If $t_0$ is generic, every rational map defined at $t_0$ sends it to the generic point of its image: the induced inclusion of the image function field into the torsor function field is injective. The six absolute upper bounds become equalities. Subtracting transcendence degrees in the three field towers gives the relative equalities. Algebraic base extension, justified in Section~2, transfers the conclusions to $k$. The comparison with the last column of the table now follows from $g=b+u$.
\end{proof}

\begin{proof}[Proof of Corollary~\ref{cor:real}]
After an algebraic enlargement defining the motivated isomorphism $M\simeq c^*M$, the tensor categories generated by $M$ and by $N$ are equal. Their Betti automorphism groups are then identified by the first projection, so the geometric kernel $N_2$ is trivial. Over the original field its connected dimension is still zero, by the algebraic-extension invariance in Section~2.2; hence $v=0$ there as well. An isomorphism of unadorned vector spaces, or even a realization isomorphism not known to be motivated, would not justify this argument.

For a real model, let $f_\infty$ be rational cohomological real Frobenius. The conjugate Hodge flag is $f_\infty\phi$. Its intersections with $\phi$ therefore recover $\psi$ rationally, while forgetting the second grading recovers $\phi$. Hence $\Q(\psi)=\Q(\phi)$. On de Rham cohomology the opposite flag is
\[
                    \alpha f_\infty\alpha^{-1}F_D.
\]
Its intersections with $F_D$ recover the normalized splitter by rational matrix operations. Thus $\beta\in k(\alpha)$ and $K=k(\alpha)$.

For $E:y^2=x^3-x+1$, the standard Weierstrass formula gives
\[
            \Delta=-16\bigl(4(-1)^3+27\bigr)=-368,
 \qquad j=1728\frac{4(-1)^3}{4(-1)^3+27}=-\frac{6912}{23}.
\]
A CM elliptic $j$-invariant is an algebraic integer \cite[Theorem 12.14 and Corollaries 12.15--12.16, p.~125]{MilneMF}, so $E$ has no CM. For completeness, its Hodge group is $\SL_2$. Indeed it is a connected reductive subgroup of the symplectic group $\SL_2$ and contains the nontrivial Hodge circle. A proper connected reductive subgroup of $\SL_2$ is a torus. Here its real points contain the anisotropic Hodge circle, so its rational commuting algebra is an imaginary quadratic field. A nonrational element of that field is a rational endomorphism preserving the Hodge line; after clearing denominators it preserves the period lattice and defines complex multiplication on $E$. This is excluded. The central weight cocharacter then makes the Mumford--Tate group $\GL_2$. The motivated group contains it, since motivated tensors are Hodge tensors, and is already a subgroup of $\GL(H^1(E))$, so it too is $\GL_2$. The Hodge filtration stabilizes a line, giving $g=4$, $b=3$, $u=1$. The rational model gives $v=0$. The five bounds are $3,1,4,0,3$, strictly below the predictions $4,2,5,1,4$.
\end{proof}

\subsection*{A matrix calculation of the obstruction}
The rank-two case can be seen without group terminology. Choose a rational eigenbasis for $f_\infty=\operatorname{diag}(1,-1)$ and an algebraic de Rham basis adapted to its Hodge line. On the open chart $cd(ad-bc)\ne0$, write
\[
 A=\begin{pmatrix}a&b\\c&d\end{pmatrix},\qquad
 h=-\frac{ad+bc}{2cd},\qquad
 S=\begin{pmatrix}1&h\\0&1\end{pmatrix}.
\]
The opposite line is $Af_\infty A^{-1}$ applied to the first coordinate line; multiplication by $S$ sends it to the second coordinate line. With $\Delta_A=ad-bc$ one obtains
\[
 SA=\begin{pmatrix}\Delta_A/(2d)&-\Delta_A/(2c)\\c&d\end{pmatrix}.
\]
Its entries satisfy $z_{11}z_{22}+z_{12}z_{21}=0$. Conversely $c,d,\Delta_A$ are independent coordinates on this image chart. Thus the inverse composite has dimension three, while $A$ has dimension four; $h$ is already a rational function of its entries. This is the smallest instructive case of the general calculation.

\subsection*{The obstruction without a real model}
The condition $v=0$ is broader than existence of a real model. Consider
\[
 E_7/\Q(i):\qquad y^2=x^3+2x^2+(\tfrac12+i)x.
\]
For $a=2$ and $b=\tfrac12+i$, set $x'=x+a+b/x$ and $y'=y(1-b/x^2)$. The identity
\[
 x'\bigl((x'-a)^2-4b\bigr)=x'(x-b/x)^2
                         =y^2(1-b/x^2)^2
\]
shows that the target is $y'^2=x'^3-4x'^2+(2-4i)x'$. The map extends between the smooth projective curves, preserves infinity and has degree two: the original $x$ satisfies $x^2+(a-x')x+b=0$. Scaling the conjugate curve by $(x,y)\mapsto(u^2x,u^3y)$ with $u^2=-2$ gives that target. Hence $E_7$ is isogenous to its conjugate and $h^1(E_7)\simeq h^1(\overline E_7)$ with rational coefficients.

Its invariants give
\[
 j(E_7)=\frac{-56512+38144i}{25},\qquad
 N_{\Q(i)/\Q}(j(E_7))=\frac{929714176}{125}.
\]
The nonintegral norm excludes CM; the nonreal $j$ excludes a real model. Thus $v=0$ although real descent fails. If $\tau$ is an oriented period ratio, the filtration and bigrading fields are $\Q(\tau)$ and $\Q(\tau,\bar\tau)$: these are the affine coordinates of the Hodge line and its conjugate. An isogeny to the conjugate identifies their period lattices up to commensurability and a scalar. Dividing its two period equations gives a rational fractional-linear relation between $\tau$ and $\bar\tau$, with nonzero determinant. The bigrading field therefore has transcendence degree at most one. Schneider's theorem \cite[\S1, result IIc, p.~5]{Schneider} says that an algebraic $\tau$ with algebraic $j(\tau)$ is imaginary quadratic, which would give CM. Consequently $\tau$ is transcendental and the bigrading field has transcendence degree exactly one, without a period conjecture.

\section{The elliptic splitting parameter}\label{sec:elliptic-splitting}
The dimension formula concerns the splitting as a point of a torsor. For an elliptic curve the same point is a single scalar, which makes several questions and their limits explicit.

\subsection{Periods, quasi-periods and the opposite line}
Let $E/k$ have equation
\[
 y^2=4x^3-g_2x-g_3,\qquad \omega=dx/y,\qquad \xi=x\,dx/y,
\]
where $k\subset\C$ is algebraic and $g_2^3-27g_3^2\ne0$. The two classes $\omega,\xi$ form a basis of $H^1_{\dR}(E/k)$. Choose an oriented period basis $\omega_1,\omega_2$, put $\tau=\omega_2/\omega_1\in\mathfrak H$, and write $\eta_1,\eta_2$ for the quasi-periods of the Weierstrass zeta function. Thus the periods of $\xi$ are $-\eta_i$ and
\[
 \eta_1\omega_2-\eta_2\omega_1=2\pi i.
\]
Define $h$ by $H^{0,1}_{\dR}=\C(\xi+h\omega)$.

\begin{proposition}\label{prop:h-formula}
Put $A=\operatorname{Im}(\overline\omega_1\omega_2)>0$. Then
\[
 \eta_i=h\omega_i+\frac{\pi}{A}\overline\omega_i,\qquad
 h=\frac{\eta_1\overline\omega_2-\eta_2\overline\omega_1}
          {\omega_1\overline\omega_2-\omega_2\overline\omega_1}
   =\frac{\pi^2}{3\omega_1^2}E_2^*(\tau),
\]
where $E_2^*(\tau)=E_2(\tau)-3/(\pi\operatorname{Im}\tau)$.
If $\omega'=r\omega$ and $\xi'=s\xi+t\omega$, then $h'=(sh-t)/r$.
For an isogeny $f:E\to E'$ of degree $n$, written
\[
 f^*\omega'=\lambda\omega,\qquad f^*\xi'=a\omega+b\xi,
\]
one has $\lambda b=n$ and $h'=nh/\lambda^2-a/\lambda$.
\end{proposition}
\begin{proof}
Under $x=\wp(u)$ and $y=\wp'(u)$, the two differentials are $du$ and $\wp(u)du=-d\zeta(u)$. The periods of an antiholomorphic class are proportional to $\overline\omega_i$. Hence $-\eta_i+h\omega_i=-c\overline\omega_i$ for one $c$. Substitution in the Legendre relation gives $c=\pi/A$, and solving these two equations gives the displayed determinant. The classical Fourier identity $\eta_1=\pi^2E_2(\tau)/(3\omega_1)$ gives the last formula. This is also Katz's description of the opposite line: his scalar $S$ is $-12h$ \cite[Lemma 1.3.6 and (1.3.7), p.~467; (1.3.10), p.~468]{Katz}.

The basis-change formula follows by expressing $\xi'+h'\omega'$ as $s(\xi+h\omega)$. The normalized cup product of $\omega,\xi$ is unchanged by changing the elliptic equation. Pullback multiplies it by $n$, so $\lambda b=n$. Since an isogeny preserves the two Hodge lines, $a+\lambda h'=bh$, which proves the last assertion.
\end{proof}

\begin{proposition}\label{prop:cm-h}
If $E$ has complex multiplication, then $h\in k$. For
\[
                 E:y^2=4x^3-30x-28
\]
one has $j(E)=8000$ and $h=1/2$. For every algebraic CM elliptic curve,
\[
 \Qb(\omega_1,\omega_2,\eta_1,\eta_2)=\Qb(\omega_1,\pi),
 \qquad \trdeg_{\Qb}\Qb(\omega_1,\pi)=2.
\]
\end{proposition}
\begin{proof}
Over $\overline k$, the algebra $\End(E_{\overline k})\otimes\Q$ has two distinct simultaneous eigenspaces on $H^1_{\dR}(E/k)\otimes_k\overline k$. Functoriality under conjugation of endomorphisms makes this unordered pair Galois stable. One eigenline is $F^1$, which is defined over $k$, so every Galois element fixes it and hence fixes the other eigenline. Under the chosen complex embedding the latter is $H^{0,1}$, because the two CM characters are distinct. That line therefore descends to $k$. Its unique generator with $\xi$-coefficient one has $\omega$-coefficient in $k$, proving $h\in k$. The eigenspace argument over a field defining all endomorphisms is Katz's Lemma 4.0.7 \cite[p.~501]{Katz}; the preceding descent also handles a smaller field of definition of the equation.

For the displayed example put $u=x+2$, $v=y/2$ and $b=9/2$. Then
\[
 v^2=u^3-6u^2+bu.
\]
The map $(u,v)\mapsto(u-6+b/u,\,v(1-b/u^2))$ is a degree-two isogeny to $v'^2=u'^3+12u'^2+18u'$. Direct substitution verifies the equation; its kernel is the origin together with $(0,0)$. If $s^2=-2$, the scaling $(u,v)\mapsto(s^2u,s^3v)$ is an isomorphism from the original curve to this target. Composing the isogeny with its inverse gives an endomorphism $f$ with $f^*\omega=s\omega$. The exact differential
\[
 d(v/u)=(u-b/u)\omega
\]
and substitution in $x=u-2$ give $f^*\xi=-s\xi-s\omega$ in cohomology. Its eigenvalue on $H^{0,1}$ is $-s$. Thus $s(h-1)=-sh$, so $h=1/2$. The formula for $j$ gives $8000$.

For a CM lattice, $\tau$ and $\operatorname{Im}\tau$ are algebraic. Proposition~\ref{prop:h-formula} therefore expresses every quasi-period in $\Qb(\omega_1,\pi)$; the Legendre relation gives the reverse inclusion. For the transcendence degree we use Nesterenko's theorem: for $0<|q|<1$, at least three of $q,E_2,E_4,E_6$ are algebraically independent \cite[Theorem 1, p.~1320]{Nesterenko}. Here $q=e^{2\pi i\tau}$ and
\[
 E_4=\frac{3g_2\omega_1^4}{4\pi^4},\quad
 E_6=\frac{27g_3\omega_1^6}{8\pi^6},\quad
 E_2=\frac{3h\omega_1^2}{\pi^2}+\frac{3}{\pi\operatorname{Im}\tau}.
\]
If $\omega_1$ and $\pi$ were algebraically dependent, these three Eisenstein values would generate a field of transcendence degree at most one; adjoining $q$ raises it by at most one. This contradicts the cited theorem. The reverse bound is immediate from two generators.
\end{proof}

For $M=h^1(E)$ in the CM case, the motivated group over $\Qb$ has dimension two: it contains the two-dimensional Mumford--Tate torus and commutes with the CM algebra, whose centralizer on this rank-two space is that same torus. Here $u=0$ and $E$ is isogenous to its conjugate, so $v=0$. Proposition~\ref{prop:cm-h} makes the comparison point generic in the two-dimensional joint torsor. Thus the six degrees of Theorem~\ref{thm:main} are unconditionally $(2,0,2,0,0,2)$ in this case; the three displayed relative degrees are $(0,2,0)$. The genericity is supplied by the stated transcendence theorem, not by the dimension computation alone.

\begin{proposition}[Zero locus of the splitting parameter]\label{prop:h-zero}
One has $h=0$ if and only if $j(E)=0$ or $1728$.
\end{proposition}
\begin{proof}
This is the zero-locus statement for $E_2^*$, proved in \cite[Lemma 8.1, pp.~25--26]{SebbarBesrour}. We include the elementary bounds. Weight-two modularity reduces to $\tau=x+iy$ in the closed standard fundamental domain, with $|x|\le1/2$ and $y\ge\sqrt3/2$. Put $r=e^{-2\pi y}<1/200$. For $0<x<1/2$, the leading term in
\[
 \operatorname{Im}E_2(x+iy)=-24\sum_{n\ge1}\sigma_1(n)r^n\sin(2\pi nx)
\]
has strictly greater magnitude than the sum of the others: use $|\sin(n\theta)|\le n\sin\theta$, $\sigma_1(n)\le n^2$, and
\[
 S_3(r):=\sum_{n\ge1}n^3r^n=\frac{r(1+4r+r^2)}{(1-r)^4},
 \qquad S_3(r)/r-1<1.
\]
Thus a zero lies on $x=0$ or $x=\pm1/2$. On $x=0$, $y\ge1$ and
\[
 \frac d{dy}E_2^*(iy)=48\pi\sum_{n\ge1}n\sigma_1(n)e^{-2\pi ny}
                  +\frac3{\pi y^2}>0.
\]
The modular stabilizer forces a zero at $i$, so it is the only zero on that ray. On $x=1/2$, each $y^2e^{-2\pi ny}$ decreases for $y\ge\sqrt3/2$, whence
\[
 y^2\frac d{dy}E_2^*(1/2+iy)\ge\frac3\pi-36\pi S_3(1/200)>0.
\]
Indeed $S_3(1/200)<1/180$ and $\pi^2<15$. The order-three stabilizer forces the unique zero at $1/2+i\sqrt3/2$; translation handles $x=-1/2$. These are precisely the elliptic points with $j=1728$ and $0$. All estimates are strict. For example $e^{\pi\sqrt3}>200$ follows already from the first 25 terms of the exponential series at $(157/50)(173/100)<\pi\sqrt3$.
\end{proof}

\subsection{Masser's formula and an almost-holomorphic invariant}
For $j(\tau)\ne0,1728$ put
\[
 \Psi(\tau)=\frac{3E_4(\tau)E_2^*(\tau)}{2E_6(\tau)},\qquad
 \chi^*(\tau)=\frac{1728E_2^*(\tau)E_4(\tau)E_6(\tau)}
 {E_4(\tau)^3-E_6(\tau)^2}.
\]
\begin{proposition}[Masser's formula in the chosen basis]\label{prop:masser-normalization}
One has
\[
 h=\frac{g_3}{g_2}\Psi
   =\frac{3g_3}{2g_2(j-1728)}\chi^*,\qquad
 \chi^*=\frac23(j-1728)\Psi.
\]
If $\tau$ is imaginary quadratic, then
\[
 \Psi=9j\gamma+\frac{3(7j-6912)}{2(j-1728)},\qquad \gamma\in\Q(j),
\]
where $\gamma$ is specified by a modular polynomial as follows. Choose a primitive equation $C\tau^2+B\tau+A=0$, $C>0$, and set
\[
 \mu=\begin{cases}
 \left(\begin{smallmatrix}-B&-2A\\2C&B\end{smallmatrix}\right),& B\text{ odd},\\
 \left(\begin{smallmatrix}-B/2&-A\\C&B/2\end{smallmatrix}\right),& B\text{ even}.
 \end{cases}
\]
Put $D=\det\mu$ and expand
$\Phi_D(j+X,j+Y)=\sum\beta_{mn}X^mY^n$.
Except when $B$ is odd and $D=3d^2$ for an integer $d>1$,
\[
 \gamma=\frac{\beta_{20}-\beta_{11}+\beta_{02}}{\beta_{10}},
 \qquad \beta_{10}=\beta_{01}\ne0.
\]
In that exceptional case,
\[
 \gamma=\frac{\beta_{40}-\beta_{31}+\beta_{22}-\beta_{13}+\beta_{04}}
                  {\beta_{30}},
 \qquad \beta_{30}=\beta_{03}\ne0,\quad\beta_{21}=\beta_{12}=0.
\]
\end{proposition}
\begin{proof}
The assertions about the modular polynomial and $\Psi$ are Masser's known formula: Theorem A1, Lemmas A1--A2 and equations (105)--(106) \cite[pp.~114, 116--119]{Masser}. The coefficients $\beta_{mn}$ lie in $\Q(j)$, so the specified quotients do also. To match the normalization, substitute
\[
 g_2=\frac{4\pi^4}{3\omega_1^4}E_4,\qquad
 g_3=\frac{8\pi^6}{27\omega_1^6}E_6
\]
into Proposition~\ref{prop:h-formula}. This gives $h=(g_3/g_2)\Psi$. The identity
$j-1728=1728E_6^2/(E_4^3-E_6^2)$ gives the remaining formulas. In particular, for algebraic $g_2,g_3\ne0$, algebraicity of $h$ is equivalent to algebraicity of $\chi^*(\tau)$. The displayed rational formula is used only away from the two elliptic values, since its derivation divides by $j'(\tau)$. At $j=0,1728$, Proposition~\ref{prop:h-zero} instead gives $h=0$. Masser's table at p.~121 has $\Psi(i\sqrt2)=15/28$; multiplying by $28/30$ recovers $h=1/2$ for Proposition~\ref{prop:cm-h}.
\end{proof}

\subsection{Conjugate period ratios and an algebraic example}
\begin{proposition}\label{prop:conjugate-ratio}
Let $E/\Qb$ be non-CM and let $\tau\in\mathfrak H$ be its period ratio. Then
\[
       \overline\tau\in\Q(\tau)
       \quad\Longleftrightarrow\quad E\text{ is isogenous to }\overline E.
\]
If $\tau$ and $\overline\tau$ are algebraically dependent but the curves are not isogenous, their irreducible relation over $\Q$ is symmetric, with equal degree at least two in each variable. The fixed field of conjugation in $\Q(\tau,\overline\tau)$ is
\[
                 \Q(\tau+\overline\tau,\tau\overline\tau).
\]
There is an algebraic elliptic curve not isogenous to its conjugate, with $j=1/(2+i)$.
\end{proposition}
\begin{proof}
Schneider's theorem makes $\tau$ transcendental for non-CM $E/\Qb$ \cite[Theorem IIc, p.~5]{Schneider}. If $\overline\tau=f(\tau)$ with $f\in\Q(X)$, conjugation gives $\tau=f(f(\tau))$, hence the rational-function identity $f\circ f=X$. Degrees multiply, so $f$ has degree one. A rational fractional-linear relation between the two ratios means that their lattices are commensurable up to scale, which is exactly isogeny. The same lattice criterion proves the converse (the ratio of the conjugate lattice is $-\overline\tau$). In this case the two flag fields are equal, and their common transcendence degree is one. Moreover $h\in\Qb(\omega_1,\omega_2,\eta_1,\eta_2)$: an isogeny $E\to\overline E$ is defined over $\Qb$, since its finite kernel is torsion and the quotient and resulting isomorphism have algebraic models. Its rational Betti matrix and algebraic de Rham matrix express the conjugate period matrix in terms of the original one. The determinant formula for $h$ then belongs to the stated field.

For the dependent case, the kernel of $\Q[X,Y]\to\Q(\tau,\overline\tau)$ is a height-one prime generated by an irreducible $P$. Conjugation interchanges $X,Y$, so $P(Y,X)=\pm P(X,Y)$ after scaling. The minus sign would make $X-Y$ a factor, forcing $\tau=\overline\tau$, which is impossible. Symmetry gives equal bidegrees; degree one would put $\overline\tau$ in $\Q(\tau)$ and has just been excluded. The two variables are roots of $T^2-(\tau+\overline\tau)T+\tau\overline\tau$. Conjugation is nontrivial, so the displayed quadratic extension has degree two and its base is the fixed field.

For the example put $j_0=1/(2+i)$, $c=j_0/(j_0-1728)$ and take $g_2=3c,g_3=c$. Its invariant is $j_0$. At the Gaussian prime $(2+i)$, $j_0$ has negative valuation whereas $\overline j_0$ is integral. Isogenous elliptic curves have simultaneous integrality of their $j$-invariants: factor an isogeny into cyclic isogenies and use the modular polynomial, which is integral and monic in either variable \cite[Theorem 6.1 and proof, pp.~88--90]{MilneMF}. The two curves cannot therefore be isogenous. Its nonintegral $j$ also excludes CM.
\end{proof}

\subsection{The joint flag quotient and the generated torsor}
\begin{proposition}\label{prop:elliptic-joint-flags}
Suppose that $E/\Qb$ is not isogenous to $\overline E$. The motivated group of $h^1(E)\oplus h^1(\overline E)$ is
\[
 \Delta_2=\{(g_1,g_2)\in\GL_2^2:\det g_1=\det g_2\},\qquad\dim\Delta_2=7.
\]
Its comparison torsor has a surjective map to $\mathbf P^1\times\mathbf P^1$ with fibres of dimension five, obtained by pulling back the two de Rham Hodge lines. At the comparison point its value, in oriented bases, is
$([1:\tau],[1:-\overline\tau])$.
Thus the period-genericity hypothesis for this joint motive implies algebraic independence of $\tau$ and $\overline\tau$.
\end{proposition}
\begin{proof}
A CM elliptic curve is isogenous to its conjugate, so both factors here are non-CM. The known Hodge-group calculation for pairwise nonisogenous elliptic curves gives $\SL_2\times\SL_2$; all Hodge classes on their powers are algebraic \cite[Proposition 5.14, pp.~20--21]{KSV}. Adjoining their common weight character therefore gives the stated motivated group. This also explains why there is one common determinant rather than two independent determinants.

Each Hodge line is algebraic on the de Rham side, so pullback defines an algebraic morphism. The group $\SL_2\times\SL_2$ acts transitively on the target, proving surjectivity and the fibre dimension $7-2=5$. For the conjugate curve choose periods $(\overline\omega_1,-\overline\omega_2)$; these have the required orientation. This gives the asserted coordinates. A dominant morphism sends a generic point to a generic point, since pullback of a nonzero rational function is nonzero. This proves the conditional independence assertion without claiming that the required period hypothesis is known.
\end{proof}

The tensor relation behind this implication is explicit. If
$P(X,Y)=\sum c_{ij}X^iY^j$ has bidegree $(m,n)$, then
\[
 P(\tau,-\overline\tau)=0
 \quad\Longleftrightarrow\quad
 \sum_{i,j}(-1)^jc_{ij}\omega_1^{m-i}\omega_2^i
       \overline\omega_1^{\,n-j}\overline\omega_2^{\,j}=0.
\]
The sign disappears if one instead writes the polynomial in $(\tau,\overline\tau)$, replacing $P(X,Y)$ by $P(X,-Y)$. These products are periods of
$\operatorname{Sym}^m H^1(E)\otimes\operatorname{Sym}^nH^1(\overline E)$: use the product holomorphic differential and product cycles, and then the symmetric projector. The geometric permutation action includes the usual cohomological signs. The summand has extreme Hodge types $(m+n,0)$ and $(0,m+n)$, hence level $m+n$. A Tate twist preserves this level, whereas a pure constituent of a 1-motive has level at most one. For the unresolved equal bidegrees $m=n\ge2$, the displayed tensor cannot itself be replaced by such a constituent. This is a limitation on that direct use of a linear 1-period theorem.

There is a known result about a generated torsor which is weaker than period genericity. We make the distinction precise.
\begin{proposition}[A known elliptic de Rham--Betti consequence]\label{prop:generated-elliptic-torsor}
Let $P_E=\Isom^{\otimes}(V_{\Qb},D)$ in the category generated by $h^1(E)$ over $\Qb$, and let $Z_E$ be the $\Qb$-Zariski closure of its comparison point. For any $z\in Z_E(\Qb)$, the smallest closed subgroup containing $z^{-1}Z_E$ is the full motivated group. This conclusion does not assert $Z_E=P_E$.
\end{proposition}
\begin{proof}
A nonempty finite-type variety over the algebraically closed field $\Qb$ has a $\Qb$-point, so the choice of $z$ is possible. A closed subset of an isomorphism torsor is a subtorsor precisely when it is nonempty and is closed under the ternary composition $(x,y,z')\mapsto xy^{-1}z'$. Upon choosing one point, this is exactly closure of its translate under multiplication and inversion. Therefore $z\langle z^{-1}Z_E\rangle$ is the smallest subtorsor containing $Z_E$.

Kreutz--Shen--Vial identify this generated subtorsor with the de Rham--Betti torsor after extending Betti coefficients to $\Qb$ \cite[Proposition 3.1, p.~8]{KSV}. Their Theorem 5.13(ii), pp.~20 and 22, identifies that torsor with the motivated torsor for one elliptic factor; the condition allowing at most one CM field is automatic for a single curve. The theorem is formulated for $h(E)$, but
$h(E)=\mathbf1\oplus h^1(E)\oplus\bigwedge^2h^1(E)$, with algebraic projectors, so $h(E)$ and $h^1(E)$ generate the same tensor category. Their comparison is in the reverse direction to ours; inversion identifies the two torsors and preserves ternary closure. These observations give exactly the assertion as written. Changing $z$ to $z'=za$ with $a\in\langle z^{-1}Z_E\rangle$ does not change the generated subgroup, because the new generating set $a^{-1}z^{-1}Z_E$ contains $a^{-1}$ and recovers the old set after multiplication by $a$.
\end{proof}

\subsection{An algebraic splitting is not a subtorsor condition}
In the real case choose rational Betti coordinates in which real Frobenius is $\operatorname{diag}(1,-1)$, and write the de-Rham-to-Betti period matrix as
\[
                  P=\begin{pmatrix}x&p\\y&q\end{pmatrix}.
\]
The first column gives $\omega$ and the second gives $\xi$. The conjugate of the Hodge line $\C(x,y)$ is $\C(x,-y)$.
\begin{proposition}\label{prop:splitting-hypersurface}
The equation $h=a$ is $xq+yp+2axy=0$. Its locus $V_a\subset\GL_2$ is an irreducible three-dimensional variety isomorphic to $\Gm^3$, but is not contained in any coset of a proper algebraic subgroup of $\GL_2$.
\end{proposition}
\begin{proof}
The condition that $(p+ax,q+ay)$ be proportional to $(x,-y)$ gives the equation. Set $p'=p+ax$, $q'=q+ay$. The determinant on the equation $xq'+yp'=0$ is $-2yp'$, so invertibility forces $x,y,p',q'$ all nonzero and $q'=-yp'/x$. The three remaining coordinates identify the locus with $\Gm^3$.

It suffices to handle $a=0$, since the change of columns is multiplication by a fixed invertible matrix. Put $H=\left(\begin{smallmatrix}1&1\\1&-1\end{smallmatrix}\right)\in V_0$. The variety $H^{-1}V_0$ contains the identity and both the diagonal torus and $H^{-1}TH$, where $T$ is the diagonal torus. Their Lie algebras contain $I$, $Z=\operatorname{diag}(1,-1)$ and $X=\left(\begin{smallmatrix}0&1\\1&0\end{smallmatrix}\right)$. Together with $[Z,X]$ these span $\mathfrak{gl}_2$. Hence no proper algebraic subgroup contains $H^{-1}V_0$. If $V_0$ lay in any coset, translation by its point $H$ would give just such a subgroup containment.
\end{proof}

The de Rham projector onto $H^{0,1}$ along $H^{1,0}$ has matrix $\left(\begin{smallmatrix}0&h\\0&1\end{smallmatrix}\right)$. In Betti coordinates normalized by the period ratio it is
\[
 \frac1{\overline\tau-\tau}
 \begin{pmatrix}-\tau&1\\-\tau\overline\tau&\overline\tau\end{pmatrix}.
\]
Its image and kernel recover $\overline\tau$ and $\tau$. Thus its Betti entries are algebraic exactly in the CM case, by Schneider. This does not imply the corresponding converse for its de Rham entries: algebraicity of $h$ is a different assertion, and the preceding hypersurface calculation prevents a direct reduction to membership in a smaller comparison subtorsor.

The same projector detects all algebraic Hodge endomorphisms. Let
\[
 \mathcal D_E=\{T\in\End_{\Qb}H^1_{\dR}(E/\Qb):
                    T_{\C}\text{ preserves both Hodge lines}\}.
\]
In the ordered basis $(\omega,\xi)$, preserving $\C\omega$ and
$\C(\xi+h\omega)$ gives exactly
\[
 T=\begin{pmatrix}a&(d-a)h\\0&d\end{pmatrix}.
\]
Consequently $\mathcal D_E=\Qb\,1$ if $h$ is transcendental, and
$\mathcal D_E\simeq\Qb\times\Qb$ if $h$ is algebraic. A nonscalar such endomorphism has algebraic Betti matrix precisely when the Betti projector does: it equals $(T-a)/(d-a)$. Schneider's theorem therefore forces CM in that case. Conversely CM endomorphisms, after extension of their coefficient field to $\Qb$, give both eigenspace projectors and all of $\mathcal D_E$. The projector itself is never a rational Betti endomorphism: its nonzero image is $H^{0,1}$, which contains no nonzero real vector, since a real vector in that line would also belong to $H^{1,0}$.

For later arithmetic tests it is useful to express $h=a$ on a tensor of rank three. In the real-model coordinates above let $F=\operatorname{diag}(1,-1)$ and $D_P=\det P$. Direct multiplication gives
\[
 S=P^{-1}FP=\frac1{D_P}
 \begin{pmatrix}xq+yp&2pq\\-2xy&-xq-yp\end{pmatrix}.
\]
The equation from Proposition~\ref{prop:splitting-hypersurface} says that
$S_{11}/S_{21}=h$. For the algebraic functional
$\ell_a(T)=T_{11}-aT_{21}$ it follows that
\[
 h=a\quad\Longleftrightarrow\quad
 \ell_a(S)=\frac{xq+yp+2axy}{D_P}=0.
\]
Here $F$ is a rational Betti vector of the trace-free part of $\End(H^1)$, and $S$ is its inverse-comparison image. Polarization gives $(H^1)^\vee\simeq H^1(1)$; decomposing $H^1\otimes H^1$ into its symmetric and alternating summands identifies the trace-free part with $\operatorname{Sym}^2H^1(1)$. This is the precise tensor in which the displayed period equation lies.

\subsection{Theta normalization and exact growth}
Let $\sigma(u)$ be the Weierstrass sigma function for the chosen lattice and put
\[
          \theta_a(u)=e^{-au^2/2}\sigma(u),\qquad a\in\C.
\]
\begin{proposition}\label{prop:theta-growth}
For every integer $n$, the quotient $\theta_a(nu)/\theta_a(u)^{n^2}$ is independent of $a$. If $c=\pi/A$ and $M_a(r)=\max_{|u|\le r}|\theta_a(u)|$, then
\[
             \lim_{r\to\infty}\frac{\log M_a(r)}{r^2}
                    =\frac{c+|h-a|}{2}.
\]
Consequently $a=h$ is the unique parameter minimizing this quadratic growth. In the normalization $\sigma(u)=u+O(u^5)$, the fields generated by the Taylor coefficients of $\theta_a$ over $k$, and by the Laurent coefficients of
\[
        e^{-auv}\frac{\sigma(u+v)}{\sigma(u)\sigma(v)},
\]
are both exactly $k(a)$.
\end{proposition}
\begin{proof}
The exponential factors cancel in the quotient because $(nu)^2=n^2u^2$. Integrating $\sigma'/\sigma=\zeta$ and $\zeta(u+\lambda)-\zeta(u)=\eta_\lambda$, with the normalization that $\sigma$ is odd, gives sigma's usual translation rule. Its absolute value, together with Proposition~\ref{prop:h-formula}, gives for every lattice vector $\lambda$
\[
 |\theta_h(u+\lambda)|
   =e^{c\operatorname{Re}(\overline\lambda u)+c|\lambda|^2/2}
       |\theta_h(u)|.
\]
Thus $B(u)=|\theta_h(u)|e^{-c|u|^2/2}$ is continuous and lattice-periodic. It is bounded, and is positive at any chosen $u_0$ outside the lattice. Hence
\[
 |\theta_a(u)|=B(u)\exp\!\left(\frac{c|u|^2+\operatorname{Re}((h-a)u^2)}2\right)
\]
gives the upper bound. For the lower bound choose a direction maximizing the last real part. Every large point on that ray is within a fixed distance of $u_0+\lambda$ for some lattice vector $\lambda$. Choose the point at distance $r-O(1)$ from the origin; then $|u_0+\lambda|\le r$, $B(u_0+\lambda)=B(u_0)>0$, and the quadratic exponent differs from $(c+|h-a|)r^2/2$ by $O(r)$. This proves the limit, including $a=h$.

The differential equation for $\wp$ recursively gives its Laurent coefficients in $\Q[g_2,g_3]$, and $\zeta'=-\wp$, $\sigma'/\sigma=\zeta$ give the normalized sigma coefficients in $k$. Now $[u^3]\theta_a=-a/2$, while
\[
 e^{-auv}\frac{\sigma(u+v)}{\sigma(u)\sigma(v)}
       =\frac1u+\frac1v-a(u+v)+\text{terms of higher total degree}.
\]
Every coefficient lies in $k(a)$ and the displayed coefficients recover $a$.
\end{proof}

The division quotients therefore cannot determine the quadratic exponential normalization. The growth minimization does determine it analytically, but supplies no arithmetic implication from $h\in\Qb$ to complex multiplication. Katz explicitly asks the latter question in \cite[Question 4.0.8, p.~501]{Katz}.

\subsection{Torsion translations and the Poincar\'e section}
\begin{proposition}\label{prop:theta-torsion}
The function
\[
 \Theta_h(u,v)=e^{-huv}\frac{\sigma(u+v)}{\sigma(u)\sigma(v)}
\]
represents the normalized meromorphic section of
\[
 m^*\mathcal O_E([0])\otimes p_1^*\mathcal O_E([0])^{-1}
                         \otimes p_2^*\mathcal O_E([0])^{-1}
\]
on $E\times E$, with residue one along $u=0$. For a nonzero torsion point represented by $u_0\in\Lambda\otimes\Q$, put
\[
 T_{u_0}(u)=e^{-c\overline u_0(u+u_0/2)}\theta_h(u+u_0),\qquad c=\pi/A.
\]
The field generated by its Taylor coefficients over $\Qb$ is exactly $\Qb(h)$.
\end{proposition}
\begin{proof}
Sigma's transformation rule is
\[
 \sigma(u+\lambda)=\epsilon(\lambda)
 e^{\eta_\lambda(u+\lambda/2)}\sigma(u),\qquad
 \epsilon(\lambda)\in\{1,-1\}.
\]
By Proposition~\ref{prop:h-formula}, the rule for $\theta_h$ is
$\theta_h(u+\lambda)=\epsilon(\lambda)
 e^{c\overline\lambda(u+\lambda/2)}\theta_h(u)$.
This is the reduced-theta normalization of Bannai--Kobayashi: their area parameter is $A/\pi$ \cite[Examples 1.9--1.10, pp.~10--12]{BannaiKobayashi}. Taking the quotient under addition gives the stated line bundle. The leading term $1/u$ gives its residue and fixes the multiplicative constant.

Here is a direct proof of the algebraic constant needed for the torsion assertion. For an integer $n\ge2$, the function
\[
 F_n(P(u))=\frac{\sigma(nu)}{\sigma(u)^{n^2}},\qquad
 P(u)=(\wp(u),\wp'(u)),
\]
is elliptic. Its automorphy factors cancel, including the sign since $n-n^2$ is even. Its divisor is $[n]^*[O]-n^2[O]$ and its leading term at $O$ is $n u^{1-n^2}$. These data determine a rational function over $\Q(g_2,g_3)$: the one-dimensional space with the specified divisor descends by base change for Riemann--Roch, and the algebraic leading coefficient fixes its scalar. In this normalization $F_2=-y$.

Choose $n$ such that $\lambda=nu_0\in\Lambda$, and extend the quasi-period map rationally to $\Lambda\otimes\Q$. Differentiating at $u_0$ gives
\[
 D_EF_n(P(u_0))
 =\frac{n\sigma'(\lambda)}{\sigma(u_0)^{n^2}},\qquad
 D_E=y\partial_x+(6x^2-g_2/2)\partial_y.
\]
The left side is a nonzero algebraic number: $P(u_0)$ is a nonzero algebraic torsion point and the zero of $F_n$ there is simple. Also
$\sigma'(\lambda)=\epsilon(\lambda)e^{\eta_\lambda\lambda/2}$.
Thus
\[
 C_0=e^{-\eta_{u_0}u_0/2}\sigma(u_0),\qquad
 C_0^{n^2}=\frac{n\epsilon(\lambda)}{D_EF_n(P(u_0))}
\]
is a nonzero algebraic number. Notice that $C_0=T_{u_0}(0)$.

Write $Z(u)=\zeta(u)-hu-c\overline u$, and put
$z_0=Z(u_0)$. The zeta addition formula shows that $z_0$ is algebraic for torsion points, as also follows from Proposition~\ref{prop:unitary-lift} below. More explicitly, if $u_0$ has order $n>2$, the rational addition term for $(n-1)u_0$ gives $z_0$ as $-1/n$ times an algebraic number; for order two, oddness and periodicity give $z_0=0$.
Now
\[
 \frac{T_{u_0}'(u)}{T_{u_0}(u)}
 =\zeta(u+u_0)-h(u+u_0)-c\overline u_0
 =z_0+\text{a series in }\Qb(h)[[u]]
\]
whose derivative at zero is $-\wp(u_0)-h$; all subsequent derivatives are algebraic derivatives of $\wp$ at an algebraic point. Solving this first-order formal equation with constant $C_0$ proves that all coefficients belong to $\Qb(h)$. Conversely
\[
 h=z_0^2-\wp(u_0)-\frac{T_{u_0}''(0)}{T_{u_0}(0)}
\]
recovers $h$ from those coefficients. This proves equality of fields. The algebraicity of the corresponding sigma translation is also the all-elliptic-curve statement of \cite[Theorem 2.12, pp.~26--27]{BannaiKobayashi}; no CM hypothesis is used here.
\end{proof}

These torsion constants, like division quotients, are unchanged by an arbitrary Gaussian normalization. Indeed, with
$\eta^{(a)}_{u_0}=\eta_{u_0}-au_0$,
\[
 \frac{\theta_a(nu)}{\theta_a(u)^{n^2}}=F_n(P(u)),\qquad
 e^{-\eta^{(a)}_{u_0}u_0/2}\theta_a(u_0)
       =e^{-\eta_{u_0}u_0/2}\sigma(u_0).
\]
Both identities follow by cancellation of quadratic exponents. They hold formally over every field of characteristic zero. Thus these data do not select $a=h$; the growth criterion in Proposition~\ref{prop:theta-growth} does so by a different, analytic condition.

\section{Arithmetic tests of an elliptic splitting}\label{sec:elliptic-arithmetic}
The complex splitting scalar has several arithmetic interpretations. The hypotheses in those interpretations are different: an algebraic tangent line, an algebraic point on a real compact subgroup, and a Frobenius-stable line are distinct conditions.

\subsection{The universal vector extension and its compact subgroup}
Let $E^\dagger$ be the universal vector extension of $E$. In the chosen differential normalization it has the analytic description
\[
 E^\dagger(\C)=\C^2/\Lambda^\dagger,\qquad
 \Lambda^\dagger=\Z(\omega_1,\eta_1)+\Z(\omega_2,\eta_2).
\]
Above $E\setminus\{O\}$ an algebraic chart is
\[
 (u,s)\longmapsto(x,y,t)=(\wp(u),\wp'(u),\zeta(u)-s),\qquad
 du=dx/y,\quad ds=-dt-x\,dx/y.
\]
These are the chart and invariant coframe in \cite[Sections 4.1--4.2, (4.1) and Proposition 4.2, pp.~7--8]{FonsecaMatthes}; their period of $x\,dx/y$ is the negative of our quasi-period.

\begin{proposition}\label{prop:unitary-lift}
Put $c=\pi/A$. The unique real Lie-group section with compact image is
\[
 \kappa:E(\C)\longrightarrow E^\dagger(\C),\qquad
 [u]\longmapsto[(u,hu+c\overline u)].
\]
Its image is the maximal compact subgroup. For an algebraic point $P=[u]\ne O$,
\[
 t(\kappa(P))=Z(u):=\zeta(u)-hu-c\overline u\in\Qb
 \quad\Longleftrightarrow\quad P\text{ is torsion}.
\]
The complex-linear part of $d\kappa$ has image $L_h=\{s=hu\}$, whose annihilator represents $-(\xi+h\omega)$ in de Rham cohomology.
\end{proposition}
\begin{proof}
The real span of the period lattice is
$W=\{(u,hu+c\overline u):u\in\C\}$, by the quasi-period formula. The quotient $W/\Lambda^\dagger$ is compact and projects isomorphically to $\C/\Lambda$. Every compact subgroup of $\C^2/\Lambda^\dagger$ maps trivially into the real vector group $\C^2/W$, so it lies in this subgroup. This also proves the uniqueness of the compact section. The function $Z$ is periodic by the same quasi-period formula.

We use the known theorem that an algebraic point of a universal vector extension lying in its maximal compact subgroup at one complex embedding is torsion \cite[Theorem 3.1.2, p.~386]{BostKunnemann}. If $Z(u)$ is algebraic, the three chart coordinates of $\kappa(P)$ are algebraic, so the theorem applies. Conversely, multiplication by $n$ is an isomorphism on the additive kernel. Hence $E^\dagger[n]\to E[n]$ is an isomorphism: a lift of an $n$-torsion point can be uniquely adjusted in that kernel to have order dividing $n$. This unique lift is algebraic, and belongs to the maximal compact subgroup. It equals $\kappa(P)$, proving the reverse implication. The last assertion follows by differentiating the real-linear section and using $ds=-dt-\xi$.
\end{proof}

For torsion points the algebraic coordinate can also be computed directly. Zeta's addition formula gives
\[
 Z(u+v)=Z(u)+Z(v)
 +\frac{\wp'(u)-\wp'(v)}{2(\wp(u)-\wp(v))}.
\]
Its doubled value is interpreted by the limit
\[
 Z(2u)=2Z(u)+\frac{\wp''(u)}{2\wp'(u)}.
\]
For a point $P$ of order $n>2$, successive addition gives
$Z(mu)=mZ(u)+R_m(P)$ for $1\le m\le n-1$, where $R_m$ is a rational expression in the coordinates of the multiples of $P$. The doubled denominator is nonzero; all later denominators are nonzero because $mP\ne\pm P$ for $2\le m\le n-2$. Since $(n-1)P=-P$ and $Z$ is odd and periodic, $Z(u)=-R_{n-1}(P)/n$. A point of order two has $Z(u)=0$ directly. In particular, at a point of order three,
\[
                    Z(u)=-\frac{12x(P)^2-g_2}{12y(P)}.
\]

For an explicit non-CM instance, let $E_*:y^2=4x^3-4x-1$, and let $r$ be the root in $(1,2)$ of
$3r^4-6r^2-3r-1=0$. This root is unique: the polynomial is strictly convex on $[1,\infty)$, is negative at $1$ and $3/2$, has positive derivative at $3/2$, and is positive at $2$. In particular $r>3/2$, so $s=(4r^3-4r-1)^{1/2}>0$. Put $P=(r,s)$. The defining equation gives
$((3r^2-1)/s)^2=3r$; the tangent formula on $Y^2=x^3-x-1/4$ therefore gives $2P=-P$. Thus $P$ has exact order three and
\[
       \kappa(P)=\left(r,s,-\frac{3r^2-1}{3s}\right)
\]
is algebraic. The invariant $j(E_*)=110592/37$ is nonintegral, so the curve is not CM. Algebraic compact points are therefore not by themselves a criterion for CM; here they are exactly the torsion points.

\subsection{Algebraic tangent lines and restricted powers}
\begin{proposition}\label{prop:no-algebraic-splitting-line}
For every $a\in\Qb$, the line $L_a=\{s=au\}\subset\Lie E^\dagger$ is not the Lie algebra of an algebraic subgroup. Its reductions fail to be closed under restricted $p$-th powers at infinitely many good primes of every number field of definition. Moreover
\[
             \exp(L_a)\cap E^\dagger(\Qb)=\{0\}.
\]
Consequently, for $P=[u]\in E(\Qb)\setminus\{O\}$, the number $\zeta(u)-au$ is transcendental.
\end{proposition}
\begin{proof}
A connected algebraic subgroup with this Lie algebra would project onto $E$. Its kernel is a finite subgroup of $\mathbf G_a$, hence is trivial in characteristic zero. It would give an algebraic section of the universal extension. On analytic universal covers such a group section is linear, $u\mapsto(u,au)$, and would require $\eta_i=a\omega_i$ for both periods. This contradicts the Legendre relation. Thus no such subgroup exists.

Bost's restricted-Lie criterion states that a Lie subalgebra defined over a number field is algebraic if its good reductions are closed under $p$-th powers at all but finitely many primes \cite[Theorem 2.3, p.~172]{BostLeaves}. Its contrapositive proves the assertion about infinitely many primes, after any finite field extension. The analytic subgroup theorem says that an algebraic Lie line whose analytic exponential contains a nonzero algebraic point contains the Lie algebra of a positive-dimensional algebraic subgroup \cite[Hauptsatz, p.~3]{WustholzAnalytic}. This contradicts the first paragraph, proving the intersection assertion. If $\zeta(u)-au$ were algebraic, the chart would make $[(u,au)]$ an algebraic point of that intersection; its projection is $P\ne O$, which is impossible.
\end{proof}

In the chart above $L_a$ is generated by
\[
 D_a=y\partial_x+(6x^2-g_2/2)\partial_y-(x+a)\partial_t,
 \qquad V=-\partial_t.
\]
The all-prime calculation below specializes in characteristic five to
\[
 D_a^5=2g_2D_a+(3g_3-2g_2a)V.
\]
For the CM curve $(g_2,g_3,a)=(30,28,1/2)$ the discriminant is $5832$, a unit at five, and this gives $D_a^5=4V$. Thus even the actual algebraic CM splitting line need not be restricted-stable at every good prime.

The tangent condition and the compact-point condition cannot be substituted for one another. Indeed $\exp(L_h)$ meets the maximal compact subgroup only at its identity: membership of $(u,hu)$ in $W$ requires $c\overline u=0$, and adding a period vector does not change membership in $W$. Algebraicity of $h$ specifies the tangent line; the compact-point theorem instead requires the complete value $\zeta(u)-hu-c\overline u$ to be algebraic.

\subsection{Frobenius fixity with its coefficient field}
\begin{proposition}\label{prop:projector-frobenius-density}
Let $E$ and a rank-one projector $e\in\End_LH^1_{\dR}(E/L)$ be defined over a number field $L$. Its crystalline-Frobenius fixity set has lower weighted density at most $3/4$, with the convention
\[
 \underline\delta(e)=\liminf_{X\to\infty}
 \frac{\displaystyle\sum_{\substack{v:p_v\le X\\\varphi_v(e)=e}}
          [L_v:\Q_{p_v}]\log(p_v)/(p_v-1)}
 {[L:\Q]\displaystyle\sum_{p\le X}\log(p)/(p-1)}.
\]
Finitely many ramified and bad places are omitted. At an unramified good place write $\varphi_v=F_v\sigma_v$, where
$F_v=\left(\begin{smallmatrix}A&B\\C&D\end{smallmatrix}\right)$ and $\sigma_v$ acts on coefficients. For $e_h=\left(\begin{smallmatrix}0&h\\0&1\end{smallmatrix}\right)$,
\[
 \varphi_v(e_h)=e_h
 \quad\Longleftrightarrow\quad C=0,\qquad A\sigma_v(h)+B=hD.
\]
\end{proposition}
\begin{proof}
Tang's theorem for geometrically simple abelian varieties makes an endomorphism fixed with lower weighted density greater than $3/4$ an element of $\End_{\Qb}(E)\otimes\Q$ \cite[Definition 2.3.1, p.~7; Corollary 5.2.3, p.~23]{Tang}. This is a division algebra for an elliptic curve: every nonzero endomorphism is an isogeny. It has no nontrivial idempotent. Thus a rank-one projector cannot satisfy the strict density inequality. In particular it cannot be a $1$-absolute-Tate endomorphism. The coefficient field matters: $\End(E)\otimes\Qb$ in the CM case is a product and does have projectors, but it is not the rational endomorphism algebra in the cited theorem.

For the equation, Frobenius fixity is $F_v\sigma_v(e_h)=e_hF_v$. The two sides are
\[
 \begin{pmatrix}0&A\sigma_v(h)+B\\0&C\sigma_v(h)+D\end{pmatrix},
 \qquad
 \begin{pmatrix}hC&hD\\C&D\end{pmatrix}.
\]
Their equality is precisely the claimed condition.
\end{proof}

The same CM curve $y^2=4x^3-30x-28$ gives a concrete failure at five for its rational projector $e_{1/2}$. Its reduction is $y^2=4x^3+2$. The counts of solutions above $x=0,1,2,3,4$ are $0,2,2,1,0$, so there are six points including infinity. Frobenius has polynomial $T^2+5$ on crystalline $H^1/\Q_5$. This is irreducible, since a square in $\Q_5$ has even valuation. There is no invariant $\Q_5$-line, hence no fixed rank-one projector. Coefficient Frobenius on $\Q_5$ is the identity; no extension of the coefficient field has been made in this example.

\begin{proposition}\label{prop:one-place-frobenius}
If $E/L$ is non-CM, then at every unramified good-reduction place of residue characteristic $p>2$, absolute crystalline Frobenius does not preserve $\operatorname{Fil}^1H^1_{\dR}(E/L_v)$. In particular, it does not fix a rank-one projector with that kernel.
\end{proposition}
\begin{proof}
Write the residue field as $\mathbf F_q$, $q=p^f$, and $W=W(\mathbf F_q)=\mathcal O_{L_v}$. On
$M=H^1_{\dR}(\mathcal E/W)=H^1_{\mathrm{cris}}(E_{\mathbf F_q}/W)$ let $F$ be absolute crystalline Frobenius. If the fraction-field Hodge line is stable, its saturated integral line is stable because $F(M)\subset M$. In a basis starting on that line, $F$ has semilinear matrix
\[
                        \begin{pmatrix}a&b\\0&d\end{pmatrix},\qquad a,b,d\in W.
\]
Modulo $p$, Frobenius annihilates invariant differentials, so $a\in pW$. Cup product makes the determinant have valuation one. Therefore $v_p(a)=1$ and $v_p(d)=0$. The linear iterate $F^f$ has eigenvalues of valuations $f$ and zero, and is nonscalar.

The special-fibre endomorphism $\pi_q$ has crystalline action $F^f$, which preserves the integral Hodge line. The Serre--Tate--Grothendieck--Messing morphism criterion therefore lifts this particular endomorphism to $\mathcal E/W$. More explicitly, for $p>2$ the canonical divided powers on $pW/p^n$ are nilpotent, because
\[
 v_p(p^m/m!)\ge m\frac{p-2}{p-1}\longrightarrow\infty.
\]
Grothendieck--Messing lifts the morphism of $p$-divisible groups respecting the filtration at each level. Serre--Tate's fully faithful deformation functor, with the original special-fibre endomorphism specified, gives compatible abelian-scheme morphisms. Formal existence algebraizes their graphs. The exact criterion and its explanation are \cite[Theorem 3.4, Remark 3.5, pp.~12--13; Appendix B, pp.~35--37]{LiftingTraces}; the original deformation functor is \cite[Theorem 1.2.1, pp.~143--145]{KatzSerreTate}. Thus this step uses an actual special-fibre endomorphism, not a general assertion about morphisms of local Tate modules.

The lifted endomorphism is noninteger because its crystalline action is nonscalar. It descends to $\Qb$ after choosing an algebraic closure inside $\overline{L_v}$: its finite kernel consists of torsion points of the original algebraic curve, the quotient by that kernel is algebraic, and an isomorphism of this quotient with $E$ is defined over $\Qb$. Any difference from the given isomorphism is an automorphism of $E$, also defined over $\Qb$. This produces a noninteger endomorphism in characteristic zero, contradicting the non-CM hypothesis. Finally a fixed projector has Frobenius-stable kernel.
\end{proof}

The conclusion is a criterion for an additional crystalline equation. It does not establish that algebraicity of the complex scalar $h$ implies that equation. The explicit CM failure at five also rules out imposing that equation at every good place solely from algebraicity of $h$.

\subsection{The conjugate quotient and its differential scale}
The scale $c=\pi/A$ is essential when the opposite complex line is compared with an algebraic elliptic curve. We first record an elementary subgroup fact that makes this distinction precise.

\begin{lemma}\label{lem:no-graph-subgroup}
Let $F$ be a complex elliptic curve and $a,b\in\C$ with $b\ne0$. No positive-dimensional connected algebraic subgroup of $E^\dagger\times F$ has Lie algebra contained in
\[
             W(a,b)=\{(u,s,v):s-au-bv=0\}.
\]
\end{lemma}
\begin{proof}
Such a subgroup $H$ has trivial intersection with the vertical additive group: the intersection has zero Lie algebra, and $\mathbf G_a$ has no nontrivial finite subgroup in characteristic zero. The projection identifies $H$ with its algebraic image in $E\times F$, so $H$ is an abelian variety.

Every algebraic homomorphism from an abelian variety to $E^\dagger$ is zero. Indeed its analytic image is a compact complex subgroup and therefore lies in the maximal compact subgroup of Proposition~\ref{prop:unitary-lift}. Its tangent space is contained in
\[
                 W_{\R}=\{(u,hu+c\overline u):u\in\C\}.
\]
This real plane contains no nonzero complex line. If both $(u,s)$ and $i(u,s)$ belong to it, comparison of their second coordinates gives $2ic\overline u=0$, hence $u=s=0$. Thus $H$ lies in the factor $F$, whose tangent intersects $W(a,b)$ trivially because $b\ne0$. This is impossible for a positive-dimensional $H$.
\end{proof}

We use the analytic subgroup theorem also for a nonzero period logarithm. The precise refined form says that, for a nonzero $z\in\Lie G_{\C}$ with $\exp_G(z)\in G(\Qb)$, the smallest algebraic linear subspace containing $z$ is the Lie algebra of a connected algebraic subgroup. In particular, if $z$ lies in an algebraic linear subspace $W$, one can choose such a subgroup with $z\in\Lie H_{\C}\subset W_{\C}$. This is \cite[Theorem 6.2, printed pp.~43--44, PDF pp.~59--60]{HuberWustholz}. The case $\exp_G(z)=0$ is included: divide $z$ by a sufficiently large integer to obtain a nonzero torsion point, then apply the theorem. This is a named transcendence input throughout the following arguments.

\begin{proposition}\label{prop:conjugate-quotient}
For every algebraic elliptic curve in our normalization, there is a surjective analytic group homomorphism
\[
 \Phi_h:E^\dagger(\C)\longrightarrow\overline E(\C),\qquad
 [(u,s)]\longmapsto[(s-hu)/c]\pmod{\overline\Lambda}.
\]
Its kernel is the injective image of $u\mapsto[(u,hu)]$, and
\[
 \Phi_h(\kappa([u]))=[\overline u],\qquad
 \Phi_h^*(dv)=c^{-1}(ds-h\,du).
\]
It is never algebraic. Moreover $(h,c)\notin\Qb^2$.
\end{proposition}
\begin{proof}
For a lattice vector $\lambda$, the equality $\eta_\lambda-h\lambda=c\overline\lambda$ proves descent. The points $(0,cz)$ prove surjectivity. If $(s-hu)/c=\overline\lambda$, subtracting $(\lambda,\eta_\lambda)$ puts the representative on $s=hu$. This proves the kernel assertion; that line meets the period lattice only at zero because $c\overline\lambda=0$ forces $\lambda=0$. The compact-section identity and differential follow by substitution.

On the vertical $\mathbf G_a$, the map is $s\mapsto[s/c]$, with nonzero derivative. An algebraic map from $\mathbf A^1$ to an elliptic curve extends to $\mathbf P^1$ by properness and is constant by Riemann--Hurwitz. Thus $\Phi_h$ cannot be algebraic.

If $h,c$ were algebraic, $W(h,c)\subset\Lie(E^\dagger\times\overline E)$ would be algebraic and would contain the nonzero period logarithm
$(\omega_1,\eta_1,\overline{\omega_1})$. The analytic subgroup theorem would give a positive-dimensional algebraic subgroup with Lie algebra in this plane, contrary to Lemma~\ref{lem:no-graph-subgroup}.
\end{proof}

For the CM example $(g_2,g_3,h)=(30,28,1/2)$, this gives a nonalgebraic analytic homomorphism with algebraic projective differential $[-1/2:1]$. Its actual differential includes the transcendental factor $1/c$. The alternative target $\C/(c\overline\Lambda)$ makes the differential $(-h,1)$, but its raw Weierstrass coefficients are $c^{-4}\overline g_2,c^{-6}\overline g_3$. At least one is nonzero, and if $h$ is algebraic each nonzero one is transcendental. The analytic isomorphism to the algebraic model $\overline E$ is $[z]\mapsto[z/c]$, or $(X,Y)\mapsto(c^2X,c^3Y)$, and restores the differential factor.

The same conclusion about an algebraic differential follows from the known dense-torsion analytic-homomorphism lemma \cite[Lemma 6.5, printed pp.~45--46]{HuberWustholz}. Its differential hypothesis concerns the actual linear map, not its projective class. The direct argument above uses the analytic subgroup theorem with the displayed period vector.

\subsection{Second-kind lattices and algebraic quotient models}
\begin{proposition}\label{prop:second-kind-lattice}
For $a\in\C$ put
\[
 \mu_i=\eta_i-a\omega_i,\qquad
 \Lambda_a=\Z\mu_1+\Z\mu_2.
\]
This is a full lattice exactly when $|h-a|\ne c$. In that case
\[
 \operatorname{Im}\frac{\mu_2}{\mu_1}
       =\frac{A(|h-a|^2-c^2)}{|\mu_1|^2},
 \qquad
 F_a:E\setminus\{O\}\longrightarrow\C/\Lambda_a,
 \quad[u]\longmapsto[\zeta(u)-au]
\]
is holomorphic, with $F_a^*(dz)=-(x+a)\omega$. It has no holomorphic extension at $O$ and is not algebraic. If $a\in\Qb$, the target elliptic curve has CM if and only if $E$ has CM and $a=h$.
\end{proposition}
\begin{proof}
The real-linear map $T_a(z)=(h-a)z+c\overline z$ carries $\Lambda$ to $\Lambda_a$ and has determinant $|h-a|^2-c^2$. When that determinant is nonzero it takes a lattice to a lattice; when it is zero its real rank is one, so its image cannot be a full lattice. The determinant also gives the displayed oriented-area formula. Off the circle, $\mu_1\ne0$.

The quasi-period law proves that $F_a$ descends, and $\zeta'=-\wp$ gives its differential. This differential has a double pole with nonzero leading coefficient at $O$. A holomorphic extension would pull a regular differential back to a regular differential, a contradiction. An algebraic map from a punctured smooth curve to a proper curve extends across the puncture, so $F_a$ is not algebraic.

For the last assertion, a CM target has imaginary-quadratic basis ratio $r=\mu_2/\mu_1$, with either orientation. Thus
\[
                \eta_2-r\eta_1-a\omega_2+ar\omega_1=0.       \tag{$\dagger$}
\]
Suppose first that $E$ is non-CM. Apply the refined analytic subgroup theorem to the nonzero period vector
$(\omega_1,\eta_1,\omega_2,\eta_2)$ in $(E^\dagger)^2$ and to the algebraic hyperplane defined by $(\dagger)$. It gives a proper connected subgroup $H$ containing that vector in its Lie algebra. Its image $B\subset E^2$ cannot be proper: since $\End^0(E)=\Q$, a proper abelian subvariety is killed by a nonzero integer pair $(m,n)$, which would imply $m\omega_1+n\omega_2=0$.

If $H$ surjects onto $E^2$, the quotient $(E^\dagger)^2/H$ is a vector group. As $H$ is proper, this gives a nonzero homomorphism $(E^\dagger)^2\to\mathbf G_a$. There is no such homomorphism. For one factor the exact sequence obtained by applying $\operatorname{Hom}(-,\mathbf G_a)$ to the universal extension has connecting map
\[
 \operatorname{Hom}(\mathbf G_a,\mathbf G_a)
          \longrightarrow\operatorname{Ext}^1(E,\mathbf G_a),
\]
an isomorphism by universality; also $\operatorname{Hom}(E,\mathbf G_a)=0$. Hence $\operatorname{Hom}(E^\dagger,\mathbf G_a)=0$. This contradiction excludes the non-CM case. Equivalently, the four periods are linearly independent over $\Qb$, the known special case of \cite[Definition 15.1 and Lemma 15.6, printed pp.~133, 136--138]{HuberWustholz} for $[0\to E]$.

If $E$ is CM, $\tau$ and $h$ are algebraic. The earlier CM period theorem makes
\[
 U=\omega_1,\qquad V=c\overline{\omega_1}
                      =\frac{\pi}{\omega_1\operatorname{Im}\tau}
\]
linearly independent over $\Qb$: a dependence would make $\pi$ an algebraic multiple of $\omega_1^2$. Writing $b=h-a$, equation $(\dagger)$ becomes
$b(\tau-r)U+(\overline\tau-r)V=0$. Therefore $r=\overline\tau$ and $b=0$. Conversely, $a=h$ gives $\Lambda_h=c\overline\Lambda$, whose elliptic curve is CM. Its displayed ratio has negative imaginary part; reversing the lattice orientation gives the usual upper-half-plane convention.
\end{proof}

\begin{proposition}\label{prop:raw-second-kind}
If $a\in\Qb$ and $\Lambda_a$ is full, its raw invariants $g_2(\Lambda_a)$ and $g_3(\Lambda_a)$ cannot both be algebraic. Every $d\in\C^\times$ for which $d^{-1}\Lambda_a$ has both invariants algebraic is therefore transcendental.
\end{proposition}
\begin{proof}
Otherwise the differential-normalized lattice defines an elliptic curve $F/\Qb$. The nonzero vector
$(\omega_1,\eta_1,\mu_1)$ is a period logarithm in $E^\dagger\times F$ and belongs to the algebraic plane $W(a,1)$. The analytic subgroup theorem contradicts Lemma~\ref{lem:no-graph-subgroup}. Finally
\[
 g_2(d^{-1}\Lambda_a)=d^4g_2(\Lambda_a),\qquad
 g_3(d^{-1}\Lambda_a)=d^6g_3(\Lambda_a),
\]
so an algebraic nonzero $d$ cannot make both invariants algebraic. This argument does not divide by either invariant and includes the cases where one vanishes.
\end{proof}

The raw-period assertion also follows from Movasati's known equal-period-module result: two simple abelian varieties with equal full integral period modules of nonzero de Rham classes are isomorphic with those classes matched \cite[Corollary 2, p.~14]{Movasati}. Here the classes would be $-\xi-a\omega$ and a regular differential on $F$. Pullback preserves $F^1$, whereas the first class has nonzero $\xi$-coefficient.

These conclusions leave a precise possibility open. At $a=h$, $j(\Lambda_h)=\overline{j(E)}$ is algebraic, although every differential scale leading to an algebraic model is transcendental if $h$ is algebraic. In Movasati's shifted family
$Y^2=4(X-a)^3-g_2(X-a)-g_3$, the class $X\,dX/Y$ pulls back to $\xi+a\omega$. At $a=h$ all its nonzero integral-cycle periods equal $-c\overline\lambda$ and are nonzero. The zero-period hypothesis in \cite[Theorem 2(2), p.~3]{Movasati} fails there; the at-most-one-algebraic-point assertion in its part (1) does not exclude one such point.

\subsection{The symmetric-square tensor and conjugate extension}
\begin{proposition}\label{prop:symmetric-conjugate}
For a real algebraic short model with three real roots, choose $\omega_1=w>0$ and $\omega_2=\tau w\in i\R_{>0}$, and set
\[
 T_a=\begin{pmatrix}2a&1\\1&0\end{pmatrix},\qquad
 P=\begin{pmatrix}w&-(h+c)w\\ \tau w&-(h-c)\tau w\end{pmatrix}.
\]
The product-cycle matrix of
\[
 \theta_a=\operatorname{pr}_1^*\omega\wedge\operatorname{pr}_2^*\xi
 +\operatorname{pr}_1^*\xi\wedge\operatorname{pr}_2^*\omega
 +2a\operatorname{pr}_1^*\omega\wedge\operatorname{pr}_2^*\omega
\]
is $PT_aP^{\mathsf T}$, and
\[
 \frac{PT_aP^{\mathsf T}}{2\pi i}
 =\begin{pmatrix}
 (a-h-c)/(c\tau)&(a-h)/c\\
 (a-h)/c&(a-h+c)\tau/c
 \end{pmatrix}.
\]
At $a=h$ it is $\operatorname{diag}(-\tau^{-1},\tau)$. For a non-CM curve, no nonzero scalar multiple of this vector has all Betti coordinates algebraic.

For any algebraic short model, the unique complex-linear map taking the universal-extension period basis to its conjugate is
\[
 M=\begin{pmatrix}-h/c&1/c\\c-h\overline h/c&\overline h/c\end{pmatrix}.
\]
It induces an analytic group isomorphism $E^\dagger\to\overline E^\dagger$ with $\det M=-1$ and $\overline M M=I$. It is not algebraic, and its entries cannot all be algebraic.
\end{proposition}
\begin{proof}
The quasi-period relations give $\eta_1=(h+c)w$, $\eta_2=(h-c)\tau w$ and $c\tau w^2=\pi i$. Thus $\det P=2\pi i$, and multiplication gives the displayed matrix. At $a=h$ the ratio of its two nonzero coordinates is $-\tau^2$. Algebraicity after a common rescaling would make this ratio algebraic and therefore make $\tau$ algebraic, contrary to Schneider's theorem for non-CM $E$.

The exact projector comparison includes a determinant twist. For
$J=\left(\begin{smallmatrix}0&1\\-1&0\end{smallmatrix}\right)$, the map
\[
 \operatorname{Sym}^2V\otimes(\det V)^{-1}\longrightarrow\mathfrak{sl}(V),
 \qquad T\longmapsto TJ
\]
is an isomorphism, with inverse $U\mapsto-UJ$. Since $P^{\mathsf T}JP=(\det P)J$,
\[
 (PTP^{\mathsf T}/\det P)J=P(TJ)P^{-1},\qquad
 T_aJ=\begin{pmatrix}-1&2a\\0&1\end{pmatrix}=2e_a-I.
\]
At $a=h$ this is the earlier Hodge projector. Its centralizer is
\[
 \left\{\begin{pmatrix}r&a(s-r)\\0&s\end{pmatrix}:r,s\in\C^\times\right\},
\]
because a commuting matrix preserves the two distinct eigenlines. An algebraic stabilizer for an algebraic complement $a$ does not make the actual Betti tensor algebraic.

The matrix $M$ sends $(u,hu+cv)$ to $(v,\overline h v+cu)$. On the period basis take $v=\overline u$; its image is the conjugate period basis. Legendre makes this a complex basis, proving uniqueness. Direct substitution gives the determinant and involution identities. The composite with $\overline E^\dagger\to\overline E$ is $\Phi_h$, so the map is not algebraic. If every entry of $M$ were algebraic, its entry $1/c$ and the ratio of its first-row entries would make $c,h$ algebraic, contradicting Proposition~\ref{prop:conjugate-quotient}.
\end{proof}

\subsection{A quantitative restricted-Lie density bound}
\begin{proposition}\label{prop:restricted-density}
Let $K$ contain $g_2,g_3,a$ with $a\in\Qb$, and put $d=[K:\Q]$. Omit finitely many ramified, bad and small primes. Define
\[
 \beta_K(a)=\liminf_{R\to\infty}
 \frac{\displaystyle\sum_{\substack{v\mid p,\ p\le R\\ K_p+aH_p=0\text{ in }k(v)}}
                  [K_v:\Q_p]\log p/(p-1)}
      {\displaystyle\sum_{v\mid p,\ p\le R}[K_v:\Q_p]\log p/(p-1)}.
\]
Then $\beta_K(a)\le1-1/(4d)$. The complementary set has upper weighted density at least $1/(4d)$.
\end{proposition}
\begin{proof}
We apply Herblot's quantitative algebraicity theorem and verify its geometric input. The additive torsor $E^\dagger\to E$ gives an extension
$0\to\mathcal O_E\to\mathcal V\to\mathcal O_E\to0$. The projective bundle of lines
$X=\mathbb P_E(\mathcal V)$ compactifies $E^\dagger$: the complement of the first subline has affine coordinate $s=v_0/v_1$. Analytically its cocycle is
\[
 (u,v_0,v_1)\sim(u+\lambda,v_0+\eta_\lambda v_1,v_1),
 \qquad \Theta_a(u)=[u;(au,1)].
\]
Choose $N$ large enough that
$\mathcal L=\mathcal O_X(1)\otimes\pi^*\mathcal O_E([O])^{\otimes N}$ is ample, and fix a smooth Hermitian metric on $\mathcal V$. The tautological vector $v_a(u)=(au,1)$ has
\[
                  C_1\le\|v_a(u)\|\le C_a(1+|u|).
\]
For the upper bound write $u=w+\lambda$ with $w$ in a compact fundamental parallelogram, and use $|\eta_\lambda|=O(|\lambda|)$. For the lower bound apply the bounded quotient map $\mathcal V\to\mathcal O_E$, which sends the vector to $1$.

The reduced theta multiplier gives a holomorphic frame $e_h$ for the pullback of $\mathcal O_E([O])$ with norm $e^{-c|u|^2/2}$. Indeed its multiplier has modulus
$\exp(c\operatorname{Re}(\overline\lambda u)+c|\lambda|^2/2)$, whose inverse is the ratio of these frame norms. Hence
\[
 s_a=v_a\otimes e_h^{-N}\in\Gamma(\C,\Theta_a^*\mathcal L^{-1}),\qquad
 C_1e^{Nc|u|^2/2}\le\|s_a(u)\|
       \le C_a(1+|u|)e^{Nc|u|^2/2}.                       \tag{*}
\]
This is a nonvanishing holomorphic growth section of order two in the exact sense of \cite[Definition 5.1, p.~29; Definitions 5.4--5.5, p.~31]{Herblot}. Here the quadratic type of a nonzero section $s$ means $\limsup_{r\to\infty}r^{-2}\log\max_{|u|=r}\|s(u)\|$. Formula $(*)$ also proves the optimal type for this bundle and metric: every other nonzero section is $f s_a$ for an entire $f\ne0$. If $f$ has order of vanishing $m$ at zero, Jensen's formula and the lower bound in $(*)$ give
\[
 \log\max_{|u|=r}\|f(u)s_a(u)\|
          \ge Nc\,r^2/2+m\log r+O(1).
\]
The section $s_a$ attains the limiting coefficient $Nc/2$. Thus this least quadratic type is the same for every $a$; algebraicity of the actual $h$ is not needed to construct the metric.

The Zariski closure of $\Theta_a(\C)$ is a connected algebraic subgroup of $E^\dagger$ projecting onto $E$. A one-dimensional closure would have Lie algebra $L_a$, contrary to Proposition~\ref{prop:no-algebraic-splitting-line}. Its dimension is therefore two. The formal leaf has the same closure by the analytic identity theorem.

The restricted-power calculation says that the invariant foliation is $p$-closed precisely at the primes in the numerator defining $\beta_K(a)$. Herblot's Definition 3.5 and Proposition 3.6, pp.~20--21, make the formal leaf $(1-\beta_K(a))$-arithmetic. That definition uses a lower weighted density, as above.

Regard $X$ as a projective integral $\Q$-variety via the finite \'etale morphism $\Spec K\to\Spec\Q$. Its identity is a closed point of degree $d$, and its dimension is still two. Herblot's Lemma 2.6, pp.~12--13, preserves arithmeticity under finite base extension; Definition 4.6 and its conjugacy argument, pp.~27--28, give it at all embeddings of the residue field. The analytic uniformization requires one geometric point above the closed point, supplied by our fixed complex embedding. A $\Q$-closed subset containing the dense chosen complex component contains its conjugates, so the formal closure over $\Q$ has dimension two as well.

Herblot's Theorem 6.1, pp.~31--32, says that for formal curves at closed points of degrees $d_j$, with arithmeticity parameters $\alpha_j$ and a common uniformization of order $\rho$, a closure of dimension $r>1$ satisfies
\[
                    \sum_j\frac1{\alpha_jd_j}
                       \le\frac r{r-1}\rho.
\]
Here $r=2$, $\rho=2$, there is one point, and $\alpha=1-\beta_K(a)$. Apply the theorem with a slightly larger positive $\alpha$ if necessary, and let it decrease. It gives $1/((1-\beta_K(a))d)\le4$, as claimed.
\end{proof}

This is an application of the cited arithmeticity and uniformization theorem. It does not assert that the actual complex $h$, if algebraic, satisfies the ordinary restricted-Lie congruences on a density-one set.

\subsection{An ordinary integral line and isogeny denominators}
\begin{proposition}\label{prop:isogeny-ordinary-orbits}
At a place of good ordinary reduction with $p>3$, the inverse-Verschiebung construction gives an integral line in $\Lie E^\dagger$ projecting isomorphically onto $\Lie E$. Write it as $s=\upsilon_p(E)u$ in a good short N\'eron frame. It is functorial under actual isogenies. If $C\subset E$ has order $N$, $S_C=\sum_{P\in C\setminus\{0\}}x(P)$, and its quotient $f:E\to E'$ has $f^*\omega'=\lambda\omega$, then
\[
 \lambda^2h(E')=Nh(E)+S_C,
 \qquad\lambda^2\upsilon_p(E')=N\upsilon_p(E)+S_C.          \tag{1}
\]
If $h\in K$, the Hodge coefficients of the differential-preserving quotients obey
$\sigma(h(E_C))=h(E_{\sigma C})$ for $\sigma\in\Gal(\Qb/K)$ fixing the original equation.

Suppose $h$ is integral at this place, $h\ne\upsilon_p(E)$, and
$b=v_p(h-\upsilon_p(E))\ge0$, with $v_p(p)=1$. For a cyclic subgroup $C$ of order $p^n$, let $p^m$ be the order of its intersection with multiplicative torsion. In good N\'eron frames,
\[
 v_p(\lambda_C)=m,
 \qquad d_p(C):=\max(0,-v_p(h(E_C)))=\max(0,2m-n-b).       \tag{2}
\]
The number of these subgroups is $p^n$ for $m=0$, $(p-1)p^{n-m-1}$ for $1\le m<n$, and $1$ for $m=n$. Their total is $D_n=(p+1)p^{n-1}$, and their total depth is at most
\[
 S_{2r}=\frac{2(p^r-1)}{p-1},\qquad
 S_{2r+1}=\frac{(p+1)p^r-2}{p-1}.
\]
For non-CM $E$, every $\Gal(\Qb/K)$-orbit $\mathcal O$ of such subgroups satisfies
\[
 \frac1{|\mathcal O|}\sum_{C\in\mathcal O}d_p(C)
       \le I_p\frac{S_n}{D_n}=O_{E,K,p}(p^{-n/2}),
 \quad I_p=[\GL_2(\Z_p):\rho_p(\Gal(\Qb/K))].
\]
\end{proposition}
\begin{proof}
Use the ordinary moduli neighborhood and Verschiebung tower of \cite[\S2.2, pp.~16--19; Theorem 2.2, p.~18]{MochizukiConnections}. Its maps $V^j:E^{F^j}\to E$ are finite \'etale. Identity first-order points lift uniquely and compatibly through this tower. The normalized splittings
$\kappa[p^j]:E^{F^j}_{\Z/p^j}\to E^\dagger_{\Z/p^j}$ of the cited construction give, in the inverse limit, an integral lift of the tangent line of $E$. Its composite with projection is the identity. Pulling the whole construction along the classifying map treats any given ordinary elliptic scheme; the classifying point itself is not assumed \'etale. The theorem identifies this tower with the completed \'etale integral structure on the universal extension.

The classes $\omega,\xi$ form an integral de Rham basis when $p>3$: the principal part of $\xi$ is the differential of an integral simple pole, and its cup pairing with $\omega$ is a unit. Thus the resulting slope is integral. For an isogeny $f$, the maps $f^{F^j}$ commute with the Verschiebung maps because $f$ preserves multiplicative torsion. The two composites with $\kappa[p^j]$ lift the same map to $E'$. Their difference is a morphism from a proper elliptic scheme to its vector kernel, hence is constant, since its global functions are those of the base; it vanishes at the identity. Passage to the inverse limit proves functoriality.

For the differential-preserving quotient, the lattice has index $N$ over $\Lambda$ and the addition formula gives
\[
 \wp_{\Lambda_C}(u)=\sum_{P\in C}\wp_\Lambda(u+u_P)-S_C.
\]
For $P\ne O$, put $F_P=(y-y(P))/(2(x-x(P)))$. Zeta's addition formula gives
$dF_P=(\wp(u)-\wp(u+u_P))du$. Consequently
\[
 f_C^*\xi_C=N\xi-S_C\omega-d\sum_{P\ne O}F_P.
\]
In a general short frame this becomes
\[
 f^*\xi'=\frac N\lambda\xi-\frac{S_C}{\lambda}\omega,
 \qquad (u,s)\longmapsto
 \left(\lambda u,\frac N\lambda s+\frac{S_C}{\lambda}u\right).
\]
Apply this map to the Hodge line and to the functorial integral line to obtain (1). The differential-preserving formula is algebraic in the subgroup coordinates, so applying $\sigma$ fixing $g_2,g_3,h$ proves the asserted Galois covariance of the actual Hodge coefficients.

Factor $f_C$ through its multiplicative intersection of order $p^m$. The dual of this first isogeny is \'etale and has unit differential, while their composite is $[p^m]$, so its differential has valuation $m$. The remaining kernel maps isomorphically onto its subgroup in the \'etale quotient. Integrally, after splitting that quotient, each component receives a finite flat rank-one component, which is the base itself; hence this remaining isogeny is \'etale. This proves the first equality of (2). Subtracting (1) gives
\[
 \lambda_C^2(h(E_C)-\upsilon_p(E_C))
                   =p^n(h(E)-\upsilon_p(E)),
\]
so the difference on the left target has valuation $b+n-2m$. Adding the integral slope $\upsilon_p(E_C)$ gives exactly (2). Changes of good short N\'eron frame multiply the slope by a unit square, preserving its negative valuation.

Choose a geometric torsion basis with multiplicative first axis. Cyclic lines in $(\Z/p^n)^2$ have unique representatives $(a,1)$ or $(1,z)$ with $z\in p\Z/p^n$. The first set has $m=0$ and cardinality $p^n$; in the second $m=v_p(z)$, with $z=0$ corresponding to $m=n$. This proves the counts. Since $b\ge0$, the total depth is at most
\[
             n+(p-1)\sum_{1\le k<n/2}p^{k-1}(n-2k).
\]
Subtract its value at $n-2$. The even and odd increments are respectively $2p^{r-1}$ and $(p+1)p^{r-1}$, with initial values $S_1=1,S_2=2$, proving the formulas.

Serre's open-image theorem \cite[\S4.4, Th\'eor\`eme~3 and $3'$, p.~299]{SerreAdelic} makes $I_p$ finite for non-CM $E$. The full group modulo $p^n$ acts transitively on the $D_n$ lines, and a subgroup of index at most $I_p$ has every orbit of size at least $D_n/I_p$. Its nonnegative depth sum is at most the total sum $S_n$, which proves the stated orbit average and its decay.
\end{proof}

The unique fully multiplicative branch has depth $\max(0,n-b)$, while the orbit averages tend to zero. These are local norms in good N\'eron frames. A global height comparison would require a coherent metrized line at all places, and is not supplied by this calculation.

\subsection{Algebraic points on the actual quotient graph}
\begin{proposition}\label{prop:graph-rank}
If the actual $h$ is algebraic, then
\[
 S_h=\operatorname{graph}(\Phi_h)\cap(E^\dagger\times\overline E)(\Qb)
\]
has torsion-free rank at most two. Its two period logarithms alone do not contradict this hypothesis.
\end{proposition}
\begin{proof}
In $G=E^\dagger\times\overline E$, put
\[
 V=\{s-hu=cv\},\qquad W=\C(1,h,0)\subset V.
\]
The plane $V$ need not be algebraic; the line $W$ is. The exponential image of $V$ is the graph, and $V$ intersects the period kernel in
\[
             \{(\omega,\eta_\omega,\overline\omega):\omega\in\Lambda\},
\]
of rank two. The equation of $V$ determines the third coordinate uniquely, so there are no additional period directions.

For a finitely generated subgroup $\Gamma\subset S_h$, choose logarithms in $V$ and adjoin both period generators. Their generated group $Y$ has $\exp_G(Y)=\Gamma$ and period rank $\kappa=2$. Waldschmidt's Theorem 4.1 \cite[printed pp.~385--386, PDF pp.~6--7]{WaldschmidtSubgroup} gives a proper connected algebraic subgroup $G'$ for which, with
\[
 \delta=\dim(G/G'),\quad
 t=\dim W/(W\cap\Lie G'),\quad
 \lambda=\operatorname{rank}\bigl(\Gamma/(\Gamma\cap G')\bigr),
\]
one has $\delta>t$ and
\[
                 (\lambda+2\delta)(3-2)
                      \le(\delta-t)(6-\kappa),
 \quad\text{hence}\quad\lambda\le2\delta-4t.             \tag{3}
\]
Here the theorem permits the direct additive and multiplicative factors to have dimension zero; its remaining factor is all of $G$. The rank $\lambda$ is measured after exponentiation, while $\kappa$ counts periods separately.

The algebraic points in $\ker\Phi_h=\exp(L_h)$ are only the identity, by Proposition~\ref{prop:no-algebraic-splitting-line}. Therefore
$S_h\cap(E^\dagger\times0)=\{0\}$. The graph property also gives $S_h\cap(0\times\overline E)=\{0\}$.

A one-dimensional connected subgroup of $G$ is either the vertical $\mathbf G_a\times0$ or $0\times\overline E$. An affine subgroup maps trivially to the elliptic factors, and an elliptic subgroup maps trivially to $E^\dagger$ by the compactness argument in Lemma~\ref{lem:no-graph-subgroup}. Neither subgroup contains $W$, and each meets $\Gamma$ trivially. Thus $\delta=2,t=1$, and (3) gives rank zero.

For a two-dimensional $G'$, inequality (3) requires $t=0$, so $W\subset\Lie G'$. Quotient $G$ by the vertical additive group. If the image of $G'$ in $E\times\overline E$ has dimension one, it contains the tangent $\Lie E\times0$, hence is $E\times0$. Its full inverse image has dimension two, so $G'=E^\dagger\times0$. If the image had dimension two, the map would have trivial kernel and identify $G'$ with an abelian variety, whose nonzero projection to $E^\dagger$ is impossible. The only remaining two-dimensional case has $\delta=1,t=0$ and $\Gamma\cap G'=0$, giving rank at most two. Finally $G'=0$ has $\delta=3,t=1$ and gives the same bound. Every finitely generated subgroup of $S_h$ therefore has rank at most two.

If $Y$ consists of just the two period generators, then $\Gamma=0$, $\kappa=2$, and $G'=0$ satisfies the original numerical inequality $6\le8$, with $3>1$. Thus those periods themselves provide no contradiction to algebraicity of $h$.
\end{proof}

\subsection{Local geometric sizes and arithmetic coefficient growth}
\begin{proposition}\label{prop:leaf-sizes}
Let $a\in K$ and $d=[K:\Q]$. Choose fixed \'etale parameters $X,Y$ at the identity of $E^\dagger$, over $K$, satisfying
\[
 X(e)=Y(e)=0,\qquad D_aX(e)=1,\qquad D_aY(e)=0,
\]
and write the formal leaf as $Y=\phi_a(X)=\sum_{n\ge2}b_nX^n$. Outside a finite set of places let
$\mathcal B_a=\{v:K_{p_v}+aH_{p_v}\ne0\pmod v\}$ and $d_v=[K_v:\Q_{p_v}]$.
With $\|p\|_v=p^{-d_v}$, its geometric size satisfies
\[
 \begin{aligned}
 S_v(\widehat C_a)&=\|p\|_v^{1/(p-1)}&& (v\in\mathcal B_a),\\
 S_v(\widehat C_a)&\ge\|p\|_v^{1/(p(p-1))}&& (v\notin\mathcal B_a).
 \end{aligned}
\]
The adelic size product is zero. In Andr\'e's normalization $|p|_v=p^{-d_v/d}$, the arithmetic coefficient invariant satisfies
\[
 \tau(\phi_a)\ge\limsup_{T\to\infty}\frac1T
       \sum_{\substack{v\in\mathcal B_a\\p_v\le T}}\frac{d_v}{d}\log p_v
                 \ge\frac1{4d}>0.
\]
\end{proposition}
\begin{proof}
Spread out the fixed parameters, the smooth group and the primitive vector field over $\mathcal O_K[1/N]$, enlarging $N$ so that the parameter map is \'etale at the identity. Discard ramified places as well. This is one coordinate choice over $K$, not a different choice for each prime.

The formal flow gives
\[
 X(t)=\sum_{n\ge1}\frac{D_a^nX(e)}{n!}t^n,
 \qquad Y(t)=\sum_{n\ge1}\frac{D_a^nY(e)}{n!}t^n.
\]
The coefficients below degree $p$ are integral. The inverse series $g$ of $X(t)$ also has integral coefficients below that degree, since its linear coefficient is one. In $Y(g(X))$, a possibly nonintegral coefficient $g_p$ cannot occur through the linear term because $Y'(0)=0$. Every contribution from degrees $2,\ldots,p-1$ uses only earlier coefficients of $g$. Thus
\[
 b_n\in\mathcal O_v\ (n<p),\qquad
 b_p-\frac{D_a^pY(e)}{p!}\in\mathcal O_v.
\]
In characteristic $p$, $D_a^p$ is again invariant. The line is not restricted-stable exactly when its transverse value $D_a^pY(e)$ is nonzero. At $v\in\mathcal B_a$ it is a unit; since $p!$ has valuation one,
\[
                         \|b_p\|_v=\|p\|_v^{-1}.          \tag{4}
\]

We use the precise geometric size of Bost--Chambert-Loir, which measures the largest radius on which an analytic coordinate automorphism with integral invertible derivative straightens the embedded germ. It is invariant under \'etale localization \cite[Theorem 3.1, pp.~9--10]{BostChambertLoir}. The integral shear $Y\mapsto Y+X$ changes our graph to $X+\phi_a(X)$ and preserves size. Its linear coefficient is a unit. By their Proposition 3.5(2), pp.~11--12, size is the minimum of $1$ and the exponential of the first Newton slope, namely
\[
       \min\left(1,\inf_{n\ge2}\|b_n\|_v^{-1/(n-1)}\right).
\]
The graph has positive convergence radius by local analytic existence for its nonsingular differential equation. Formula (4) gives the upper bound $\|p\|_v^{1/(p-1)}$. Their Proposition 3.6, p.~12, gives the reverse bound for every integral nonsingular rank-one foliation, and the stronger lower bound $\|p\|_v^{1/(p(p-1))}$ at absolutely unramified restricted-stable places. These hypotheses hold on the chosen neighborhood, proving both size statements.

The contribution of the restricted-stable places to $\sum\log S_v^{-1}$ is bounded by the convergent series $d\sum_p\log p/(p(p-1))$. Therefore
\[
 \sum_{p_v\le T}\log S_v^{-1}
    =\sum_{\substack{v\in\mathcal B_a\\p_v\le T}}
                 d_v\frac{\log p_v}{p_v-1}+O(1).
\]
Proposition~\ref{prop:restricted-density} makes the upper weighted density of this complementary set positive. The denominator of that density is asymptotic to $d\log T$, so the nonnegative sum is unbounded and the size product is zero. This is exactly failure of condition (ii) of A-analyticity in \cite[Definition 3.7, p.~13]{BostChambertLoir}, independently of the chosen model or \'etale parameters.

For completeness, write $\epsilon_v=d_v/d$. Andr\'e's exact coefficient definition is
\[
 \tau(\phi_a)=\inf_{P_0}\limsup_{n\to\infty}\frac1n
       \sum_{p_v\ge P_0}\max_{m\le n}\log^+|b_m|_v
 \quad\text{\cite[\S5.1, p.~70; \S5.3, p.~74]{AndreDwork}}.
\]
Each bad place with $p_v\le n$ contributes at least $\epsilon_v\log p_v$ by (4). A fixed threshold removes a bounded sum, proving the first claimed inequality. Put
\[
 \theta_{\mathcal B}(T)=\sum_{\substack{v\in\mathcal B_a\\p_v\le T}}
                         \epsilon_v\log p_v,
 \qquad H_{\mathcal B}(T)=\sum_{\substack{v\in\mathcal B_a\\p_v\le T}}
                         \epsilon_v\frac{\log p_v}{p_v-1}.
\]
Replacing $p-1$ by $p$ changes the second sum by a bounded amount. Partial summation gives
\[
 H_{\mathcal B}(T)=\frac{\theta_{\mathcal B}(T)}T
                +\int_2^T\frac{\theta_{\mathcal B}(t)}{t^2}\,dt+O(1).
\]
If $r=\limsup\theta_{\mathcal B}(T)/T$, then the integrand is at most $(r+\varepsilon)/t$ for all sufficiently large $t$. Consequently
$\limsup H_{\mathcal B}(T)/\log T\le r$. Since $\sum_{v\mid p}\epsilon_v=1$, Proposition~\ref{prop:restricted-density} bounds the left side below by $1/(4d)$. This proves the second inequality using only upper densities.
\end{proof}

Andr\'e's Theorem 5.4.3, p.~76, requires $\tau=0$, finite $\rho$, and simultaneous meromorphic uniformizations whose radii have product greater than one. The actual embedded graph fails its zero-$\tau$ hypothesis. Bost--Chambert-Loir's Theorem 6.1, p.~38, requires A-analyticity, which also fails. At the archimedean place the entire map $u\mapsto[(u,au)]$ does force the canonical tangent seminorm to vanish, by their Proposition 5.15, p.~33, applied to discs of arbitrarily large radius; it does not supply the missing finite-place product. All these conclusions hold for every algebraic $a$, including the known CM splitting coefficients, and therefore give no non-CM exclusion for algebraic $h$ alone.


\section{Modular coefficients and horizontal equations}\label{sec:elliptic-equations}
\subsection{Two modular constructions recover the same scalar}
Let $f(q)=\sum_{n\ge1}b(n)q^n$ be a normalized rational weight-two newform and put
\[
                 \varepsilon(q)=\sum_{n\ge1}\frac{b(n)}nq^n.
\]
For an optimal elliptic curve, use the lattice normalization of Alfes--Griffin--Ono--Rolen \cite[(1.4)--(1.8), p.~3]{AGOR}. Their holomorphic Weierstrass function is
\[
 \widehat{\mathfrak Z}^{+}(q)=\zeta(\varepsilon(q))-S(\Lambda)\varepsilon(q).
\]
The scalar $S(\Lambda)$ equals $h$. Indeed their corrected zeta function
$\zeta(u)-S(\Lambda)u-(\pi/A)\overline u$ is periodic \cite[(2.2), p.~9]{AGOR}; its two period increments give precisely the linear equations defining $h$ in Proposition~\ref{prop:h-formula}. Hence
\[
 [q]\widehat{\mathfrak Z}^{+}
             =\frac{b(2)^2}{4}-\frac{b(3)}3-h.
\]
To verify this, write $\varepsilon=q+b(2)q^2/2+b(3)q^3/3+O(q^4)$ and use $\zeta(u)=u^{-1}+O(u^3)$. All coefficients of the latter Laurent series are rational in the algebraic differential normalization. Thus every coefficient of $\widehat{\mathfrak Z}^{+}$ is rational minus a rational multiple of $h$, and its first coefficient recovers $h$. Algebraicity of that one coefficient is equivalent to algebraicity of all of them and of $h$.

There is an exact geometric comparison for a general modular parametrization.
\begin{proposition}\label{prop:mock-splitting-scalar}
Let $\pi:X_0(N)\to E$ be defined over $\Q$, of degree $d$, and suppose
$\pi^*\omega=m\omega_f$, where $m\in\Q^\times$ and $\omega_f=f(q)dq/q$.
The class
\[
                   \eta_f^{\mathrm{alg}}=\frac m d\pi^*\xi
\]
is cup-dual to $\omega_f$. In its complex Hodge decomposition
$\eta_f^{\mathrm{alg}}=a\omega_f+b\overline{\omega_f}$, one has
\[
                         a=-\frac{m^2}{d}h.
\]
Candelori's weight-zero harmonic Maass form attached to this class has all holomorphic coefficients algebraic exactly when $h$ is algebraic, provided its constant term is algebraic. With rational constant term its coefficient field is exactly $\Q(h)$.
\end{proposition}
\begin{proof}
Our basis has normalized cup product $\langle\omega,\xi\rangle=1$, as follows from its period determinant $2\pi i$. Pullback multiplies the pairing by $d$, giving
$\langle m^{-1}\pi^*\omega,(m/d)\pi^*\xi\rangle=1$.
In particular pullback is injective and its image is the two-dimensional $f$-component. The decomposition
$\xi=-h\omega-c\overline\omega$, with $c=\pi/A$, gives
\[
 \eta_f^{\mathrm{alg}}
   =-\frac{m^2}{d}h\omega_f-\frac{m^2}{d}c\overline{\omega_f},
\]
which proves the scalar formula.

Candelori's construction represents the algebraic class by a rational differential $\phi$ of the second kind, regular away from the cusp at infinity, and solves
$dF=a\omega_f+b\overline{\omega_f}-\phi$
\cite[(6.5) and Proposition 6.2, pp.~27--29]{Candelori}.
Write $\phi=\sum d(n)q^n\,dq/q$ with $d(n)\in\Q$. Equation (6.6) of the same source, p.~30, gives in weight two
\[
                 [q^n]F^+=\frac{ab(n)-d(n)}n\qquad(n\ne0).
\]
All nonconstant coefficients belong to $\Q(a)$, and $b(1)=1$ gives
$a=[q]F^++d(1)$. The two assertions follow, including the specified condition on the constant term.
\end{proof}

\subsection{Bounded denominators impose another equality}
\begin{proposition}\label{prop:ordinary-denominator}
Use the period lattice of Guerzhoy's construction for $f$, and write
$Z_a(q)=\zeta(\varepsilon(q))-a\varepsilon(q)$.
For every good ordinary prime outside a finite exceptional set, there is a unique $\lambda_p\in\Z_p$ such that, for all $a\in\C_p$,
\[
          \inf_n v_p([q^n]Z_a)>-\infty
                         \quad\Longleftrightarrow\quad a=\lambda_p.
\]
Algebraicity of $h$ alone does not imply bounded denominators. On the CM curve $E_0:y^2=4x^3-4x$, where $h=0$, put $t=-2x/y$ and let $u(t)$ be the formal logarithm for $dx/y$. Then
\[
 \zeta(u(t))=\frac1t-\frac12\sum_{n\ge1}
          \frac{(-1)^n\binom{2n}{n}}{4n-1}t^{4n-1}.
\]
For $n_r=(3^{2r+1}+1)/4$, the coefficient of $t^{3^{2r+1}}$ has $3$-adic valuation $-r-1$.
\end{proposition}
\begin{proof}
Guerzhoy's Theorem 2 gives
$\zeta(\varepsilon)-\lambda_p\varepsilon-\mu_p\varepsilon(q^p)\in\Z_p((q))$,
with $\lambda_p\in\Z_p$ and $\mu_p=0$ when $b(p)\not\equiv0\pmod p$ \cite[pp.~4--5]{Guerzhoy}. At an ordinary prime the Hecke recurrence
$b(p^r)=b(p)b(p^{r-1})-p b(p^{r-2})$ gives
$b(p^r)\equiv b(p)^r\pmod p$. Thus the coefficient of $q^{p^r}$ in $\varepsilon$ has valuation $-r$. If $a\ne\lambda_p$, its contribution to
$Z_a=Z_{\lambda_p}+(\lambda_p-a)\varepsilon$
has valuation $v_p(\lambda_p-a)-r$, which for large $r$ is strictly below zero and cannot cancel the integral first summand. This proves the equivalence, also for coefficients in $\C_p$.

For the example set $R=(1+4t^4)^{1/2}$ with constant term one. The curve equation gives
\[
 x=\frac{1+R}{2t^2},\quad \frac{dx}{y}=\frac{dt}{R},\quad
 x\frac{dx}{y}=\frac{1+R^{-1}}{2t^2}\,dt.
\]
Expand $R^{-1}=\sum_{n\ge0}(-1)^n\binom{2n}{n}t^{4n}$ and integrate $d\zeta=-x\,du$. Oddness fixes the integration constant and gives the displayed series. For $r\ge0$,
\[
 n_r=1+2\sum_{j=0}^{r-1}3^{2j+1}.
\]
Adding this number to itself in base three creates exactly $r$ carries, one at each odd-position digit two. The factorial-valuation formula therefore gives
$v_3\binom{2n_r}{n_r}=r$. The denominator $4n_r-1=3^{2r+1}$ gives valuation $-r-1$. Finally $t$ is an integral formal parameter with unit linear term at three; boundedness is preserved under an integral invertible change of parameter, by substitution in both directions. The failure is therefore not an artifact of a nonintegral coordinate.
\end{proof}

\subsection{The ordinary cohomological correction}
For this subsection $p\ge5$ is a good ordinary prime of an elliptic curve over $\Q$. Let
\[
 T^2-b(p)T+p=(T-\beta)(T-\beta'),\quad v_p(\beta)=0,\quad v_p(\beta')=1.
\]
Write $Vf(q)=f(q^p)$, and $\theta=q\,d/dq$. Overconvergent forms are sections on a sufficiently small wide-open neighborhood of the ordinary locus, with the cusps removed, as in \cite[Definition 3.1, pp.~5--6]{CandeloriCastella}. For weight two the natural quotient by $\theta$ is its de Rham cohomology; proper parabolic cohomology injects into this wide-open cohomology \cite[Theorem 3.3 and (6), pp.~5--6]{CandeloriCastella}.

\begin{proposition}\label{prop:ordinary-theta-criterion}
The following are equivalent:
\[
 E\text{ has CM},\qquad f-\beta Vf\in\operatorname{im}\theta,\qquad
       \Phi(\operatorname{Fil}^1)=\operatorname{Fil}^1,
\]
where $\Phi$ is crystalline Frobenius on the $f$-component, identified with $H^1_{\dR}(E/\Q_p)$.

If $E$ is non-CM and $r$ is a weakly holomorphic weight-two form over the chosen $p$-adic coefficient field with zero residues and class $\delta$ in that component, let $B(q)$ be a formal primitive, $\theta B=r$. There is exactly one $\alpha\in\C_p$ such that
\[
              (B-\alpha\varepsilon)\left|\left(1-\frac{\beta'}pV\right)\right.
\]
is overconvergent. Its exact defining equation is
\[
                      (1-\Phi/\beta)(\delta-\alpha[f])=0.
\]
It has $\alpha=0$ exactly when $\delta$ belongs to the unit-root line.
\end{proposition}
\begin{proof}
Coleman's quotient identifies $U$ with crystalline Verschiebung and gives
$\Phi\,\mathrm{Ver}=p$ in weight two \cite[Theorems 5.1 and 5.4, pp.~223 and 225]{Coleman}. Since $V=U^{-1}$ on cohomology, $[Vf]=\Phi[f]/p$. Therefore
\[
 [f-\beta Vf]=0\quad\Longleftrightarrow\quad\Phi[f]=\beta'[f].
\]
The class $[f]$ spans the Hodge line. If that line is stable, Frobenius is divisible by $p$ on its saturated integral lattice, so its eigenvalue is $\beta'$, not the unit root. The quotient theorem and restriction injectivity prove the equivalence of the second and third conditions. Proposition~\ref{prop:one-place-frobenius} proves that either condition forces CM. Conversely a rational CM newform at a good ordinary prime satisfies the CM construction in \cite[Proposition 7.1, p.~231]{Coleman}, which supplies an overconvergent primitive of $f-\beta Vf$. The rational weight-two converse is also the coefficient-field specialization of Emerton's known Theorem (1.3) \cite[pp.~3--4]{EmertonCM}.

For the last assertion, $\theta V=pV\theta$ gives
\[
 \theta\left((B-\alpha\varepsilon)\left|\left(1-\frac{\beta'}pV\right)\right.\right)
                      =(1-\beta' V)(r-\alpha f).
\]
Its class is $(1-\Phi/\beta)(\delta-\alpha[f])$. This vanishes if and only if the displayed primitive is overconvergent: the quotient theorem produces a primitive on the wide-open domain, and its expansion differs from the specified formal primitive by a constant. The operator $1-\Phi/\beta$ has one-dimensional kernel, the unit-root line, and one-dimensional image. Its value on $[f]$ is nonzero by the integral Hodge-line observation. Hence its image equation has exactly one solution $\alpha$. That solution is zero exactly when $\delta$ lies in its kernel.
\end{proof}

For non-CM $E$, the nonvanishing hypothesis $f-\beta Vf\notin\operatorname{im}\theta$ in Candelori--Castella's Theorem 5.3 is consequently satisfied \cite[p.~11]{CandeloriCastella}. The preceding proof uses the cohomological equation itself. For a class obtained from a complex algebraic splitting by Proposition~\ref{prop:mock-splitting-scalar}, vanishing of its correction still requires equality of its transported antiholomorphic line with the crystalline unit-root line. Algebraicity of the coefficients does not supply that equality.

\subsection{The defect between two differential operators}
On a fine modular curve choose a local symplectic frame $(\omega,\xi)$ with $\omega$ generating the Hodge line. Write
\[
 \nabla\omega=\omega\otimes A+\xi\otimes B,\qquad
           \operatorname{KS}(\omega^{\otimes2})=B.
\]
Thus $B$ is a local cotangent generator. Project along $L_a=\langle\xi+a\omega\rangle$ and apply the inverse Kodaira--Spencer identification. On a section of weight $k$ the resulting operator is
\[
 D_a(f\omega^{\otimes k})=
 \left(\frac{df}{B}+k\left(\frac AB-a\right)f\right)\omega^{\otimes(k+2)}.
\]
Indeed the projection sends $\xi$ to $-a\omega$, and the formula follows directly from the Leibniz rule. It gives the exact difference
\[
       (D_a-D_b)(f\omega^{\otimes k})
                     =k(b-a)f\omega^{\otimes(k+2)}.
\]
At an ordinary point let $L_{u_p}$ be the unit-root complement. After fixing an embedding into $\C_p$ of a number field containing an algebraic $h$, equality of these two operators on one section of nonzero weight and nonzero value is exactly $\iota_p(h)=u_p$. This calculation specifies the further condition needed to identify the complex and ordinary operators.

\subsection{An elliptic primitive is a \texorpdfstring{$G$}{G}-function}
\begin{proposition}\label{prop:elliptic-gfunction}
Let $E:y^2=P(x)=4x^3-g_2x-g_3$ be smooth over a number field $K$, and let $(x_0,y_0)\in E(K)$ with $y_0\ne0$. For $a\in K$ the germ
\[
 F_a(t)=-\int_0^t\frac{x_0+s+a}{\sqrt{P(x_0+s)}}\,ds,
                         \qquad\sqrt{P(x_0)}=y_0,
\]
is a $G$-function. In the coordinate $x=x_0+t$ it satisfies
\[
 2P(x)(x+a)F_a''+
              \bigl(P'(x)(x+a)-2P(x)\bigr)F_a'=0.
\]
For some integer $C\ge1$, a common denominator for its first $n$ Taylor coefficients divides
$(4C)^{n+1}\operatorname{lcm}(1,\ldots,n)$.
If $x=\wp(u)$ and $x_0=\wp(u_0)$, then
\[
 F_a(x-x_0)=\zeta(u)-au-\zeta(u_0)+au_0.
\]
The analytic specialization at $a=h$ has additive monodromy $c\overline\lambda\ne0$ on a nonzero elliptic period $\lambda$; it is asserted to be a $G$-function here only if $h$ is algebraic.
\end{proposition}
\begin{proof}
Write $P(x_0+t)=y_0^2(1+b_1t+b_2t^2+b_3t^3)$ and expand
\[
 (1+z)^{-1/2}=\sum_{m\ge0}\frac{(-1)^m\binom{2m}{m}}{4^m}z^m.
\]
Choose $C$ clearing the denominators in $\mathcal O_K$ of all $b_i$, $(x_0+a)/y_0$, and $1/y_0$. If $e_n=[t^n]F_a'$, each summand uses at most $n$ of the $b_i$ and one further cleared factor, so $(4C)^{n+1}e_n\in\mathcal O_K$. Integration divides $e_{j-1}$ by $j$, giving the claimed common denominator. The least common multiple grows at most exponentially; for example the elementary binomial-coefficient estimate for the product of prime powers between $n$ and $2n$ and iteration give such a bound. Each complex conjugate of $F_a'$ is analytic on a nonzero disc about zero, so Cauchy's estimate bounds all conjugates of its coefficients exponentially. Together with the linear differential equation, obtained by differentiating $F_a'=-(x+a)/\sqrt{P(x)}$, these are the defining $G$-function properties.

Finally $du=dx/y$ and $d\zeta=-x\,du$ give the uniformization identity with the stated constant. Translation by $\lambda$ changes its value by $\eta_\lambda-a\lambda$. At $a=h$, Proposition~\ref{prop:h-formula} makes this $c\overline\lambda$, which is nonzero.
\end{proof}

\subsection{A functional period theorem at one fibre}
Consider the Legendre family
\[
 E_t:Y^2=4X(X-1)(X-t),\qquad t\in S=\mathbf A^1\setminus\{0,1\},
\]
and the short coordinate $x=X-(1+t)/3$. Its invariants are
\[
 g_2=\tfrac43(t^2-t+1),\qquad
 g_3=\tfrac4{27}(2-3t-3t^2+2t^3).
\]
Choose periods near a real $t_0\in(0,1)$ with the first row real and the second purely imaginary on the real interval. For a fixed algebraic $a$, the equation $h(t_0)=a$ becomes
\[
 R_a(P)=xq+yp+2axy=0,\qquad
                  P=\begin{pmatrix}x&p\\y&q\end{pmatrix},
\]
where the letters in this matrix denote periods, not the fibre coordinates.

The algebraic monodromy is $\SL_2$. One can read this from the two nodal degenerations at $0$ and $1$: Picard--Lefschetz monodromy in a symplectic basis gives the two nonparallel unipotents
$\left(\begin{smallmatrix}1&2\\0&1\end{smallmatrix}\right)$ and
$\left(\begin{smallmatrix}1&0\\-2&1\end{smallmatrix}\right)$, up to simultaneous change of basis. Their Zariski closures contain the two root groups, whose Lie algebras and bracket span $\mathfrak{sl}_2$.
The period torsor $\Omega\to S$ therefore has three-dimensional determinant-fixed fibres and total dimension four. For $D=\det P\ne0$, the equation $R_a=0$ has the fibre parametrization
\[
 P=\begin{pmatrix}x&-D/(2y)-ax\\y&D/(2x)-ay\end{pmatrix},
                    \qquad x,y\ne0.
\]
Thus its hypersurface $W_a$ has total dimension three. Its intersection with a horizontal period leaf is discrete: if $R_a(P(t))$ vanished identically, continuation and Zariski-dense monodromy would force an entire determinant fibre into the proper two-dimensional hypersurface fibre. Any intersection through $t_0$ is consequently isolated and has codimension three in $W_a$.

Bakker--Tsimerman's functional theorem requires this codimension to be strictly less than the monodromy dimension \cite[Theorem 1.3, p.~5]{BakkerTsimerman}. Here it is equal to that dimension. Adding $t=t_0$ makes the codimension inequality strict, but the conclusion then holds for every point: restriction to a point has trivial monodromy, and their Corollary 2.9 characterizes containment in a proper weak Mumford--Tate subvariety by that drop \cite[p.~8]{BakkerTsimerman}. A constant family likewise has trivial relative monodromy, so their relative period theorem gives no positive lower bound from it \cite[Theorem 4.25, p.~18]{BakkerTsimerman}. These are exact checks of the indicated applications; no assertion about other applications of functional transcendence is needed.

\subsection{Horizontal jets of an actual CM splitting}
\begin{proposition}\label{prop:cm-horizontal-jets}
In the Legendre family with unshifted basis $\omega=dx/y$, $\eta=x\,dx/y$, the Gauss--Manin matrix, whose columns give the derivatives of the basis, is
\[
 C(t)=\begin{pmatrix}A&A\\B&-A\end{pmatrix},\qquad
 A=\frac1{2(1-t)},\quad B=-\frac1{2t(1-t)}.
\]
At $t=1/2$, the antiholomorphic class is $\eta-\omega/2$. Put $s=t-1/2$. Its horizontal extension $F(s)\omega+G(s)\eta$ is
\[
 G(s)=\sum_{m\ge0}\frac{\prod_{j=0}^{m-1}(4j+1)^2}{(2m)!}s^{2m},\qquad
 F(s)=\left(\frac12-2s^2\right)G'(s)-\left(\frac12+s\right)G(s).
\]
For every prime $p\equiv3\pmod4$, the coefficient of $s^{p+1}$ in $G$ has valuation $-1$ at $p$. Thus this algebraic CM initial vector has horizontal coefficients with infinitely many denominator primes; it also fails $\omega(p)$-integrality, meaning integrality of the first $\omega(p)$ coefficients with $\omega(p)/p\to\infty$.
\end{proposition}
\begin{proof}
Put $P=4x(x-1)(x-t)$. Differentiating at fixed $x$ gives
$\partial_t\omega=\omega/(2(x-t))$ and
$\partial_t\eta=\omega/2+t\partial_t\omega$.
The following two identities of rational differentials give an explicit reduction modulo exact forms:
\[
 \begin{split}
 \partial_t\omega-(A+Bx)\omega
       &=d_x\left(\frac{x(x-1)}{t(1-t)y}\right),\\
 \partial_t\eta-A(1-x)\omega
       &=d_x\left(\frac{x(x-1)}{(1-t)y}\right).
 \end{split}
\]
They follow by differentiating $y^2=P$. This proves the matrix with its signs. At $t=1/2$, the translation $x=X+1/2$ gives $y^2=4X^3-X$. Proposition~\ref{prop:h-zero} gives $h=0$ in this short basis, so the antiholomorphic class is $X\,dX/y=\eta-\omega/2$.

Horizontality is $(F,G)'=-C(F,G)$. Its second equation gives the asserted expression for $F$. Substitution in the first gives
\[
       (1-4s^2)G''-8sG'-G=0,\qquad G(0)=1,\quad G'(0)=0.
\]
The coefficient recurrence is
\[
 [s^{2m+2}]G=\frac{(4m+1)^2}{(2m+2)(2m+1)}[s^{2m}]G,
\]
which proves the series. For $m=(p+1)/2$, its numerator factors $4j+1$ lie strictly below $2p$ and none equals $p$, because $p\equiv3\pmod4$. The denominator $(p+1)!$ has valuation one. Hence the claimed coefficient has valuation $-1$. There are infinitely many primes congruent to three modulo four: a prime factor congruent to three modulo four of four times any proposed finite product minus one is new. Finally $p+1\le\omega(p)$ for all sufficiently large $p$, so the same coefficients contradict $\omega(p)$-integrality.
\end{proof}

The failure extends beyond this example by a known theorem. Lam--Litt prove that for every non-isotrivial smooth projective genus-one family over a finitely generated $\Z$-algebra, any nonzero algebraic initial vector in a positive symmetric power of relative $H^1_{\dR}$ has a horizontal extension which is neither integral nor $\omega(p)$-integral \cite[Theorem 4.1.1, p.~21]{LamLitt}. For any algebraic nonsingular Legendre parameter, enlarge $\mathcal O_K[1/N]$ so that the initial data and $t_0,(1-t_0)$ and their inverses are defined. The Legendre family over that ring is smooth, projective, connected and has nonconstant $j$. These verify every hypothesis of the cited theorem. Its conclusion applies to every nonzero initial vector and every positive symmetric power, independently of whether the initial fibre is CM.

\subsection{Two jets of the actual period product}
\begin{proposition}\label{prop:period-product-jets}
Let the algebraic real short model have three distinct real roots $e_1>e_2>e_3$, and put
\[
 \delta=e_1-e_3,\qquad m_0=\frac{e_2-e_3}{\delta}\in(0,1),
 \qquad F(m)=K(m)K(1-m),
\]
where $K(m)=\int_0^{\pi/2}(1-m\sin^2\theta)^{-1/2}\,d\theta$ uses the parameter convention. At $m_0$,
\[
 h=\delta\left(m(1-m)\frac{F'}F-\frac{2m-1}{3}\right),
 \qquad c=\frac{\delta\pi}{4F}.
\]
The two numbers $F'(m_0)/F(m_0)$ and $F''(m_0)/F(m_0)$ cannot both be algebraic. No nonzero constant multiple of the actual germ $F(m_0+z)$ has all its Taylor coefficients algebraic; in particular it cannot be an ordinary $G$-function germ over $\Qb$.
\end{proposition}
\begin{proof}
Vary $m$ while keeping $\delta$ fixed, with roots
$(2-m)\delta/3,(2m-1)\delta/3,-(1+m)\delta/3$. At $m_0$ this is the given model. Its periods are $w_1=2K(m)/\sqrt\delta$ and $w_2=2iK(1-m)/\sqrt\delta$, so the area is $4F/\delta$, proving the formula for $c$.

Put $L(m)=K(1-m)$. The change of variable
$x=\delta(X-(1+m)/3)$ from the unshifted Legendre family transforms the Gauss--Manin formula of Proposition~\ref{prop:cm-horizontal-jets} into
\[
 \nabla_{\partial_m}\omega
 =\frac{2m-1}{6m(1-m)}\omega
            -\frac1{2\delta m(1-m)}\xi.
\]
Indeed its unshifted coefficients are $A=1/(2(1-m))$ and $B=-1/(2m(1-m))$, and the new coefficients are $A+B(1+m)/3$ and $B/\delta$. Since the periods of $\xi$ are $-\eta_i$, this gives
\[
 \frac{\eta_i}{w_i}
   =\delta\left(2m(1-m)\frac{w_i'}{w_i}-\frac{2m-1}{3}\right).
\]
The two ratios on the left are $h+c$ and $h-c$. Taking their average gives the formula for $h$; subtracting them and using the already established formula for $c$ gives
\[
 \frac{K'}K-\frac{L'}L=\frac{c}{\delta m(1-m)},
 \qquad KL'-K'L=-\frac\pi{4m(1-m)}.
\]

Both $K$ and $L$ solve $y''+py'+qy=0$, where
$p=(1-2m)/(m(1-m))$ and $q=-1/(4m(1-m))$. Expanding the derivatives of $F=KL$ gives
\[
 (F')^2-2FF''-2pFF'-4qF^2=(KL'-K'L)^2.
\]
Thus, with $r=F'/F$ and $s=F''/F$,
\[
 s=\frac{r^2-2pr-4q}{2}
           -\frac{c^2}{2\delta^2m^2(1-m)^2}.             \tag{5}
\]
If $r,s$ were algebraic at $m_0$, the formula for $h$ would make $h$ algebraic and (5) would make $c^2$, hence $c$, algebraic. This contradicts Proposition~\ref{prop:conjugate-quotient}. A scalar multiple with all algebraic Taylor coefficients would have a nonzero algebraic constant term and algebraic first two derivatives, forcing $r,s$ to be algebraic as well.
\end{proof}

At the CM midpoint $m_0=1/2$, $r=h=0$ and $s=2-8c^2/\delta^2$ is transcendental. At the original singular center $m=0$, the product instead has a logarithmic singularity: $K(m)\to\pi/2$ and $K(1-m)\sim-\tfrac12\log m$. The latter follows by comparing the endpoint part of the integral with $\int dt/\sqrt{t^2+m}$. Multiplication of the two second-order equations gives the rational differential equation
\[
 F'''+3pF''+(p'+4q+2p^2)F'+(2q'+4pq)F=0,
\]
but this does not remove the actual germ's coefficient obstruction.

An algebraic logarithmic derivative at one ordinary algebraic point would not itself settle a general $G$-function question. For a direct example set
\[
                   g(z)=\log(1-z)-\tfrac14\log(1-2z).
\]
Its rational coefficients are $(2^{n-2}-1)/n$ for $n\ge1$. Their absolute values grow at most exponentially and their common denominator through degree $n$ divides $4\operatorname{lcm}(1,\ldots,n)$, which has an elementary exponential bound. The vector $(g,1)^{\mathsf T}$ solves the rational Fuchsian system with matrix
\[
 \begin{pmatrix}0&(3z-1)/(2(1-z)(1-2z))\\0&0\end{pmatrix}.
\]
Its nonzero logarithmic monodromy proves that $g$ is transcendental as a function. Nevertheless $z=1/3$ is an ordinary point inside the convergence disc, and
\[
                     g'(1/3)=0,\qquad
                     g(1/3)=\tfrac14\log(16/27)\ne0.
\]
Hence $g'/g$ has the algebraic value zero there. This example and Proposition~\ref{prop:period-product-jets} identify distinct obstructions: one pointwise value need not control a transcendental function, and the actual elliptic product cannot be turned into an algebraic-coefficient germ by constant normalization.


\section{\texorpdfstring{$p$}{p}-adic comparison and restricted Lie powers}
We first separate the loss of information under ordinary $p$-divisible-group comparison from an exact calculation in characteristic $p$. We then state the classical local reconstruction theorems with their precise data, and examine what is lost on forgetting those data or passing to infinite inverse systems.

\subsection{An ordinary curve does not acquire elliptic CM}
For $E_*:y^2=4x^3-4x-1$, the discriminant in this normalization is $37$ and $j(E_*)=110592/37$. The latter is nonintegral, so $E_*$ is non-CM. At $p=5$, the numbers of affine $y$-solutions for $x=0,1,2,3,4$ are respectively $2,2,0,1,2$. Including the point at infinity gives $\#E_*(\mathbf F_5)=8$, hence $a_5=-2$ and ordinary good reduction.

Over $\mathcal O_{\C_5}$, the connected and etale height-one factors of its $5$-divisible group have no remaining unramified twists. The extension has a Serre--Tate parameter $q\in1+\mathfrak m_{\C_5}$. Choose compatible roots $q_n$ with $q_n^{5^n}=q$ and $q_{n+1}^5=q_n$. At level $n$, the Kummer model over component $a\bmod 5^n$ is $x^{5^n}=q^a$. The points $(q_n^a,a)$ define a section: the carry in component addition is canceled by the equality $q_n^{5^n}=q$. They are compatible under $a\mapsto5a$ at the next level. Thus
\[
 E_*[5^\infty]_{\mathcal O_{\C_5}}\simeq
                   \mu_{5^\infty}\oplus\Q_5/\Z_5.
\]
Its projections supply the algebra $\Q_5\times\Q_5$ of $p$-divisible-group endomorphisms. The Hodge--Tate line, in the basis of these two Tate-module summands, is the multiplicative line and has coordinates $[1:0]$. Its field of definition is $\Q_5$. This is the ordinary situation of Howe's theorem \cite[pp.~672--673]{Howe}.

By contrast, extension between algebraically closed fields does not create endomorphisms of an elliptic curve in characteristic zero. One direct argument takes the finite kernel of a nonzero endomorphism, which is defined over the algebraic closure of the original field, forms the quotient there, and descends the resulting isomorphism using equality of $j$-invariants. Automorphisms of the elliptic curve are also defined there. Hence $E_*\otimes\C_5$ remains non-CM. An inference from algebraicity of this ordinary Hodge--Tate line to CM of the elliptic curve is therefore false.

\subsection{The all-prime identity}
We prove the second assertion of Theorem~\ref{thm:other}. Let
\[
              \omega=\frac{dx}{y},\qquad
              \xi=x\frac{dx}{y},\qquad \eta=-dt-\xi.
\]
The forms $\omega,\eta$ are a rational coframe on $k(E)(t)$. Their contractions are
\[
 \omega(D_a)=1,\quad\eta(D_a)=a,
 \qquad\omega(V)=0,\quad\eta(V)=1.
\]
Use the classical $p^{-1}$-linear Cartier operator $C$, with convention
\[
 C(f^px^m dx)=0\ (m\not\equiv p-1\bmod p),\qquad
 C(f^px^{p-1}dx)=f\,dx.
\]
Since $y^{p-1}=P(x)^{(p-1)/2}$, the numerator degrees in
\[
 \omega=\frac{P(x)^{(p-1)/2}dx}{y^p},\qquad
 \xi=\frac{xP(x)^{(p-1)/2}dx}{y^p}
\]
are less than $2p-1$. Only exponent $p-1$ contributes. Therefore
\[
 C(\omega)=H_p^{1/p}\omega,\qquad
 C(\xi)=K_p^{1/p}\omega,\qquad C(\eta)=-K_p^{1/p}\omega.
\]
The Cartier contraction identity for a closed rational form $\zeta$ and a derivation $D$ is
\[
 \zeta(D^p)-D^{p-1}(\zeta(D))=(C(\zeta)(D))^p
 \quad\text{\cite[eq.~(1.2.4.5), p.~24]{Ogus}}.
\]
Applying it to the coframe gives $\omega(D_a^p)=H_p$ and $\eta(D_a^p)=-K_p$. A derivation is determined by these two contractions, which proves
\[
                   D_a^p=H_pD_a-(K_p+aH_p)V.
\]

We also verify that $H_p$ and $K_p$ cannot both vanish for a smooth elliptic curve. The classes of $\omega$ and $\xi$ form a basis of $H^1_{\dR}(E/k)$. To check the second class at infinity, set $z=-2x/y$. Then
\[
 \omega=(1+O(z^4))dz,\qquad
 \xi=(z^{-2}+O(z^2))dz.
\]
Its principal part is $d(-1/z)$; subtracting it gives a regular local representative and a Cech transition function $1/z$. The residue $\Res_O((1/z)\omega)=1$ makes the two classes linearly independent. The conjugate filtration gives a surjection
\[
 H^1_{\dR}(E/k)\longrightarrow H^0(E^{(p)},\Omega^1_{E^{(p)}/k})
\]
induced by Cartier. Its target is one-dimensional, so the Cartier images of a basis cannot both vanish. Ordinary reduction has $H_p\ne0$, and then the unique restricted-stable line in this family has $a=-K_p/H_p$. Supersingular reduction has $H_p=0$ and $K_p\ne0$, so no such line is stable.

For example, in characteristic five direct coefficient extraction gives $H_5=2g_2$ and $K_5=2g_3$, hence
\[
                 D_a^5=2g_2D_a+(3g_3-2g_2a)V.
\]
The companion program \emph{Period comparison certificates} checks this identity and the analogous identities at seven and eleven by finite local-algebra matrices, independently of the coefficient-ring recurrence. These checks test the signs and specialization formulas; the Cartier argument proves the statement for every good prime greater than three.

\subsection{Generic fibre and good reduction are different categories}
In the rest of this section let $K$ be complete for a discrete valuation, of characteristic zero, with perfect residue field $\kappa$ of characteristic $p$. Put
\[
 K_0=W(\kappa)[1/p],\qquad C_K=\widehat{\overline K},\qquad
 G_K=\Gal(\overline K/K).
\]
The finite absolute ramification is unrestricted. We use the covariant Tate module; the exponent of $\Q_p(1)$ in a Hodge--Tate decomposition is $+1$.

\begin{proposition}\label{prop:local-categories}
The Tate functor gives equivalences
\[
 \operatorname{BT}(K)\simeq\operatorname{Rep}^{\mathrm{free,cont}}_{\Z_p}(G_K),
 \qquad
 \operatorname{BT}(K)_{\Q}\simeq\operatorname{Rep}^{\mathrm{fd,cont}}_{\Q_p}(G_K).
\]
For groups over $\mathcal O_K$ it is fully faithful, and its image consists precisely of the free stable lattices in crystalline representations with exponents in $\{0,1\}$. The corresponding potentially good generic-fibre category consists of representations acquiring this crystalline property after a finite extension.
\end{proposition}
\begin{proof}
In characteristic zero, finite group schemes are etale. Finite etale descent therefore identifies a height-$d$ $p$-divisible group with the compatible finite Galois modules arising from $(\Q_p/\Z_p)^d$. A free continuous lattice $T$ gives the groups $p^{-n}T/T$, including their inclusions and multiplication maps. After extension to $\overline K$ the divisibility axioms are the usual ones. Taking inverse limits in the other direction recovers $T$, and the constructions also invert the Hom maps.

The compact image of a continuous rational representation preserves a lattice: the orbit of any lattice is bounded inside $p^{-a}T$, and its $\Z_p$-span is a full stable lattice. Every rational intertwiner becomes an integral one after multiplication by a power of $p$. This proves the second equivalence.

The assertion over $\mathcal O_K$ uses classical classification theorems. Integral full faithfulness is Tate's Theorem~4 and Corollary~1 in \S4.2 \cite[pp.~180--181]{TateLocal}. Kisin's Corollary~2.2.6 \cite[author p.~26]{KisinLocal} gives rational existence for every crystalline representation with these exponents, including $p=2$. His proof uses Corollary~2.1.14, author p.~24, to recover the full $G_K$-action from the comparison over $G_{K_\infty}$. The integral classification is Kisin's for $p>2$ and Kim's Theorem~4.1 for $p=2$ \cite[pp.~4--5, 10]{KimLocal}. These are imported classification theorems, with the same perfect-residue-field hypothesis as here.

One can also pass from rational existence to any prescribed stable lattice directly. After scaling, the prescribed lattice $T$ contains a realized lattice $T_0$ with finite index. The finite subgroup $T/T_0$ of the generic $p$-divisible group has finite flat schematic closure in a suitable $G_0[p^n]$: its coordinate algebra is finite and torsion-free over the DVR. Quotienting by this finite flat subgroup gives a $p$-divisible group with Tate lattice $T$. The generic-fibre equivalence now proves both implications of the potentially good assertion, by applying the good-reduction result over the finite extension.
\end{proof}

Arbitrary generic-fibre groups need not become Hodge--Tate. For odd $p$ take $K=\Q_p(\mu_p)$, and for $p=2$ take $K=\Q_2(i)$. The character
\[
 \rho(g)=\exp\left(\frac{\log\chi_p(g)}{p+1}\right)
 \quad\text{satisfies}\quad \rho^{p+1}=\chi_p.
\]
The logarithm and exponential converge on $1+p\Z_p$ for odd $p$, and on $1+4\Z_2$ for $p=2$. Thus $\rho$ is a continuous integral character and gives a generic-fibre $p$-divisible group. If it became Hodge--Tate over a finite extension, its one-dimensional decomposition would have an integer exponent $r$. Tensoring $p+1$ times would give $(p+1)r=1$, which is impossible. Here distinct exponents are distinguished by Tate's vanishing $H^0(G_L,C_K(m))=0$ for nonzero integers $m$ \cite[\S3, Theorem~2, p.~176]{TateLocal}.

The prime and the connected part also matter. If $\ell\ne p$, multiplication by $\ell$ is invertible on every $p^n$-torsion group, so the $\ell$-adic Tate functor of a $p$-divisible group is zero and is not faithful. In equal characteristic $p$, even the geometric $p$-adic Tate functor kills the nonzero group $\mu_{p^\infty}$.

For a Hodge--Tate representation the canonical evaluation map is
\[
 \bigoplus_r (V\otimes C_K(-r))^{G_K}\otimes_K C_K(r)
       \xrightarrow{\sim}V\otimes_{\Q_p}C_K.
\]
Retaining the action and adjoining this grading changes no Hom space, by functoriality. If only the underlying $\Q_p$-space and its $C_K$-grading are retained, the functor remains faithful, since its underlying maps are unchanged. It need not be full: over $\Q_p$, the trivial character and its unramified quadratic twist come from etale $p$-divisible groups over $\Z_p$. Both underlying spaces are $\Q_p$, with grading concentrated in degree zero. Their Hom space in the bare graded category is $\Q_p$, while their equivariant Hom space is zero. Hodge--Tate-ness itself descends across finite extensions: after making the extension Galois, each graded invariant space is a semilinear space over that extension, and finite Galois descent preserves its dimension. Thus the same Hom distinction applies to the potentially good category. Descent of the representation cannot be discarded after finite extension: the trivial character and the ramified quadratic character of $\Q_5(\sqrt5)/\Q_5$ become equal over that field but are different over $\Q_5$.

\subsection{Filtered Frobenius data and the period-ring inverse}
A filtered Frobenius module is a finite-dimensional $K_0$-space $D$ with a bijective Witt-semilinear map $\varphi$, together with an exhaustive separated decreasing filtration on $D_K=D\otimes_{K_0}K$. Define
\[
 t_H(D)=\sum_i i\dim_K\gr^i D_K,\qquad t_N(D)=v_p(\det\varphi).
\]
The second number is independent of a basis, since changing basis multiplies the determinant by $\sigma(\det A)/\det A$, of valuation zero. Weak admissibility means $t_H(D)=t_N(D)$ and $t_H(D')\le t_N(D')$ for every $\varphi$-stable subspace with its induced filtration.

Colmez--Fontaine's theorem identifies weakly admissible filtered $(\varphi,N)$-modules with semistable representations; $N=0$ gives the crystalline category \cite[Theorem~A, p.~2; Proposition~4.2, Theorem~4.3 and corollary, pp.~23--24]{ColmezFontaineLocal}. Together with Proposition~\ref{prop:local-categories}, it gives the anti-equivalence
\[
 \operatorname{BT}(\mathcal O_K)_{\Q}^{\mathrm{op}}
 \simeq
 \left\{D\text{ weakly admissible}:\operatorname{Fil}^0D_K=D_K,
                              \ \operatorname{Fil}^2D_K=0\right\},
 \quad G\longmapsto D_{\mathrm{cris}}(V_pG^\vee).
\]
Dualization accounts for the filtration jumps $0,1$.

Here are the comparison rings specifying the inverse. Let $A_{\inf}=W(\mathcal O_{C_K}^{\flat})$ with its map $\theta$ to $\mathcal O_{C_K}$. Complete $A_{\inf}[1/p]$ at $\ker\theta$ to obtain $B_{\dR}^+$, and invert $t=\log[\epsilon]$ to obtain $B_{\dR}$. Let $A_{\mathrm{cris}}$ be the $p$-completed divided-power envelope of $\ker\theta$, compatible with the canonical divided powers of $p$, and put $B_{\mathrm{cris}}=A_{\mathrm{cris}}[1/p,1/t]$. The semistable ring $B_{\mathrm{st}}$ has monodromy $N$ with $N\varphi=p\varphi N$; its embedding into $B_{\dR}$ uses a uniformizer/logarithm convention.

For $B=B_{\dR},B_{\mathrm{cris}},B_{\mathrm{st}}$, the invariant space $D_B(V)=(B\otimes_{\Q_p}V)^{G_K}$ is over $K$ in the first case and $K_0$ in the other cases. Admissibility means that this invariant space has dimension $\dim_{\Q_p}V$ over its indicated base. On that category, evaluation is an isomorphism $B\otimes D_B(V)\simeq B\otimes V$, respecting each structure present. The semistable inverse is
\[
 V_{\mathrm{st}}(D)=
 (B_{\mathrm{st}}\otimes_{K_0}D)^{\varphi=1,N=0}
 \cap\operatorname{Fil}^0(B_{\dR}\otimes_KD_K).
\]
Use $B_{\mathrm{cris}}$ and omit $N$ for the crystalline inverse. These formulas and their quasi-inverse property are supplied by the cited admissibility theorem. The two unramified characters above give the same degree-zero filtered $K$-line but distinct crystalline Frobenius data; retaining the filtration alone loses their distinction.

For instance, $\varphi(e)=e$ and $\operatorname{Fil}^1D_K=D_K$, $\operatorname{Fil}^2D_K=0$ give $t_N=0<t_H=1$, so these raw data are not admissible. Nor can one always put Frobenius on $K$: take
\[
 K_0=\Q_2(\zeta_3),\qquad K=K_0(\pi),\qquad \pi^3=2\zeta_3.
\]
An extension of Witt Frobenius to $K$ would give a unit $z=\sigma(\pi)/\pi$ with $z^3=\zeta_3$. Reduction to $\mathbf F_4$ is impossible, since every nonzero element there has cube one. The scalar-extension map from $D$ to $D_K$ is therefore essential.

\subsection{Geometric comparisons and the integral refinement}
For proper smooth $X/K$, the de Rham comparison is
\[
 H^i_{\mathrm{et}}(X_{\overline K},\Q_p)\otimes B_{\dR}
       \simeq H^i_{\dR}(X/K)\otimes_KB_{\dR},
\]
with its filtration and Galois action. No reduction assumption is imposed. Scholze's Theorem~1.6 and Corollary~1.8 give this result for proper smooth rigid spaces \cite[pp.~4--5]{ScholzeLocal}. For a proper smooth morphism of smooth rigid spaces and a de Rham lisse $\Z_p$-sheaf, his Theorem~1.10 gives the relative filtered connection, provided the higher direct image is lisse; this holds in the algebraic situation. The constant sheaf yields the Gauss--Manin connection. The definitions of the pro-etale covers and the completed relative period sheaf are those of the corrigendum \cite{ScholzeErratumLocal}.

Good reduction supplies a crystalline comparison with rational crystalline cohomology of the special fibre. A proper regular semistable model, with reduced normal-crossing special fibre, instead supplies Hyodo--Kato cohomology, $N$, and the $B_{\mathrm{st}}$ comparison \cite[Theorem~1.1, pp.~323--324]{TsujiLocal}. Proper cohomology is potentially semistable after finite extension \cite[Corollary~2.2.3, p.~209]{TsujiICMLocal}. This concerns the representation; it does not assert that a semistable model always exists after finite extension without an alteration.

Integrally the comparison is derived. For a proper smooth formal scheme $\mathfrak X/\mathcal O_{C_K}$, Bhatt--Morrow--Scholze construct a perfect Frobenius complex $R\Gamma_{A_{\inf}}(\mathfrak X)$. Derived base change to $W(\kappa_{C_K})$, $\mathcal O_{C_K}$ and $A_{\mathrm{cris}}$ gives the corresponding crystalline and de Rham complexes. After inverting $\mu=[\epsilon]-1$ it gives the generic fibre's $\Z_p$-etale cohomology. If crystalline cohomology in degrees $i$ and $i+1$ is $p$-torsion-free, the pair consisting of $H^i_{\mathrm{et}}$ and its comparison $B_{\dR}^+$-lattice recovers integral crystalline $H^i$ with Frobenius. These are Theorems~1.4 and 1.6--1.8 of \cite[pp.~4--6]{BMSLocal}. Omitting the adjacent-degree hypothesis or replacing derived base change by ordinary tensor products would change the statement.

At the base $\operatorname{Spf}(\mathcal O_K)$, a prismatic $F$-crystal is a vector bundle on the absolute prismatic site with Frobenius invertible after inverting the prism ideal. Its etale realization gives precisely the free $\Z_p$-lattices in crystalline $G_K$-representations, by Bhatt--Scholze \cite[Definition~1.1 and Theorem~1.2, p.~2]{PrismaticCrystalsLocal}. This includes all primes and the stated perfect residue fields. The compatible crystal data, rather than an unstructured lattice alone, supply the integral reconstruction.

\subsection{The characteristic-\texorpdfstring{$p$}{p} replacements}
For smooth proper $X/\kappa$, crystalline cohomology is a perfect complex over $W(\kappa)$ with finite, possibly torsion, cohomology. Its Frobenius becomes invertible after inverting $p$. A smooth lift gives its de Rham realization and Hodge filtration; relative crystals give the corresponding integrable connection. Torsion can prevent the filtration on cohomology from being a filtration by direct summands.

For separated finite-type $X/\kappa$, rigid cohomology, with or without compact support, is finite-dimensional over $K_0$ \cite[Theorems~1.2.1--1.2.2, pp.~3--4]{KedlayaLocal}, and its linearized Frobenius is invertible \cite[Proposition~6.2 and Remark~6.3, p.~12]{DegliseLocal}. An integral lattice is additional data. The connection belongs to the coefficient isocrystal; a Hodge filtration requires additional lift data. For smooth proper $X$, rigid cohomology identifies with crystalline cohomology after inverting $p$ \cite[Proposition~1.9, author p.~13]{BerthelotLocal}.

More generally, for a bounded prism $(A,I)$ and a smooth $p$-adic formal scheme $X/(A/I)$, the prismatic complex is perfect when $X$ is proper. Its linearized Frobenius becomes invertible after inverting $I$. The crystalline specialization at $I=(p)$ and the de Rham specialization include the Frobenius base twist. The Nygaard filtration is on the Frobenius twist $R\Gamma_\Delta(X/A)^{(1)}$ \cite[Theorem~1.8, pp.~4--5; Theorem~1.16, p.~8]{PrismsLocal}. At $(W(\kappa),(p))$, Witt Frobenius is invertible, so the inverse scalar twist recovers the ordinary crystalline complex. The prismatic etale comparison requires a perfect prism and concerns the characteristic-zero generic fibre. That fibre is empty for a characteristic-$p$ scheme at the crystalline prism; it does not identify crystalline cohomology with the special fibre's geometric $p$-adic etale cohomology.

\subsection{Two sections: cocycles, compatibility and normality}
The group identities used with comparison torsors have precise additional hypotheses. Let
\[
 1\longrightarrow H\longrightarrow E\longrightarrow\pi\longrightarrow1
\]
be an extension of Hausdorff topological groups, with continuous homomorphic sections $s_1,s_2$. Let $u:\pi\to H$, $j:\Q_p^\times\to H$ and $\chi:\pi\to\Q_p^\times$ be continuous homomorphisms, and put $c=j\chi$.

\begin{proposition}\label{prop:normal-equalizer}
For arbitrary sections, $\phi(g)=s_1(g)^{-1}s_2(g)$ satisfies
\[
 \phi(gd)=s_1(d)^{-1}\phi(g)s_1(d)\phi(d).
\]
Suppose further that $s_1(\pi)$ centralizes $u(\pi)$, that $s_2(\pi)$ centralizes $c(\pi)$, and that
\[
                       s_1(g)u(g)=s_2(g)c(g)\quad(g\in\pi).
\]
Then $\phi=u c^{-1}$, its identity fibre $N$ is a closed normal subgroup, and $u(N)$ is central in $u(\pi)$. Moreover
\[
 \ker\chi\cap N=\ker\chi\cap\ker u.
\]
This intersection equals $\ker u$ if and only if $\ker u\subset\ker\chi$, which is exactly the condition that $\chi$ descend to $\pi/\ker u$. A character $\varepsilon$ with $\chi=\varepsilon\circ u$ suffices.
\end{proposition}
\begin{proof}
Substitute $s_2(g)=s_1(g)\phi(g)$ into the two section laws to obtain the cocycle formula. Cancellation in the compatibility equation gives $\phi=u c^{-1}$. The identity fibre is the equalizer of two continuous homomorphisms, hence a closed subgroup.

For normality take $g\in N$, so $u(g)=c(g)$. As $c(\pi)$ is abelian and $s_2(d)$ centralizes it,
\[
 \begin{aligned}
 u(g)&=s_2(d)c(g)s_2(d)^{-1}\\
 &=s_1(d)u(d)c(d)^{-1}c(g)c(d)u(d)^{-1}s_1(d)^{-1}\\
 &=u(d)u(g)u(d)^{-1}.
 \end{aligned}
\]
The last step uses the first global centralization. Thus $u(N)$ is central in $u(\pi)$, and $u(dgd^{-1})=u(g)=c(g)=c(dgd^{-1})$, proving normality. On $\ker\chi$, the formula for $\phi$ reduces to $\phi=u$; the intersection identity follows. Its left side is contained in $\ker\chi$, giving the necessary inclusion, and that inclusion is also sufficient. The elementary factorization criterion for a homomorphism proves the descent assertion.
\end{proof}

The normality conclusion uses the two global centralizations. For unconstrained sections, take $E=S_3\times S_3\to S_3$, $s_1(g)=(g,g)$ and $s_2(g)=(aga^{-1},g)$ for a transposition $a$. Their equalizer is the nonnormal subgroup $\{1,a\}$.

Centralization itself does not give compatibility. In $\Q_p^\times\times C_2$, put $s_1(g)=(1,g)$, $s_2(g)=((-1)^g,g)$, $u=1$, $\chi=1$, $j=\mathrm{id}$. All centralizations hold but the modified sections differ. Compatibility itself does not force character descent: instead take $\chi(g)=(-1)^g$, $u=1$ and $s_2(g)=(\chi(g)^{-1},g)$. Then $s_1u=s_2j\chi$ but $\ker u=C_2$ and $\ker\chi\cap N=\{0\}$.

\subsection{Finite dimensionality and the unramified direction}
If $\pi$ is profinite and $u:\pi\to G(\Q_p)$ is continuous for a finite-dimensional linear algebraic group $G$, its compact image is closed. Compactness identifies $\pi/\ker u$ homeomorphically with that image, and the $p$-adic closed-subgroup theorem makes it analytic \cite[Corollary~19.14, p.~97]{PstragowskiLocal}. No Hodge--Tate hypothesis enters this argument.

A Hodge--Tate grading defines a cocharacter and its differential, acting by the exponent on each summand. The assertion that its image generates the relevant smallest algebraic group is a further theorem. Serre's formulation of Sen's result assumes algebraically closed residue field \cite[\S1.1, p.~156; Theorems~1, $1'$ and 2, pp.~158--159]{SerreLocal}. For a general residue field the unramified direction cannot simply be included in that assertion.

Indeed, over $\Q_5$ choose $s\in\Z_5$ with $s^2=6$ and $s\equiv1\pmod5$. Hensel's lemma supplies it, and $s\notin\Q$. The unramified representation whose arithmetic Frobenius is
\[
              \operatorname{diag}(6,\exp(s\log6))
\]
has compact image $z\mapsto\operatorname{diag}(\exp(z\log6),\exp(sz\log6))$, $z\in\Z_5$. It has Hodge--Tate exponent zero, since its two characters come from etale $p$-divisible groups. Its grading is trivial, whereas its Lie algebra is the nonzero line $\Q_5\operatorname{diag}(1,s)$. This line is not algebraic: a one-dimensional connected subgroup of the diagonal torus is annihilated by a nonzero integral character $x^my^n$, which would force $m+ns=0$. An algebraic subgroup of $\GL_2$ with this Lie line has identity component in the diagonal torus, since its intersection with that torus has the same Lie algebra and dimension. This gives the same contradiction. The image is Zariski dense in the two-dimensional torus, by the identical character argument.

\begin{proposition}\label{prop:proalgebraic-local}
The finite-dimensional analyticity assertion fails for a proalgebraic Artin--Tate category, even with both sections, both centralizations, compatibility and character descent.
\end{proposition}
\begin{proof}
Fix $\ell\ne p$. The unramified extensions of $\Q_p$ of degrees $\ell^m$ give finite Artin motives and a quotient $r_\ell:G_{\Q_p}\twoheadrightarrow\Z_\ell$. Adjoin the Tate line and set
\[
 A=\varprojlim_m(\Z/\ell^m)_{\Q_p},\qquad
 G=A\times\Gm,\qquad u(g)=(r_\ell(g),\chi_p(g)).
\]
Each finite Artin object extends etale over $\Z_p$; its representation is crystalline of exponent zero, and the Tate line is crystalline of exponent one. Every finite tensor construction is therefore crystalline. Since cyclotomic extensions of $\Q_p$ are totally ramified, they have trivial intersection with the finite unramified extensions. Consequently
\[
                         u(G_{\Q_p})=\Z_\ell\times\Z_p^\times.
\]
At each finite Artin level this is Zariski dense in $(\Z/\ell^m)\times\Gm$, so $G$ is the actual proalgebraic group of the category.

Let $H=G(\C_p)=\Z_\ell\times\C_p^\times$, with its inverse-limit topology on the first factor, and form $E=H\rtimes G_{\Q_p}$ using the natural coefficient action. Put
\[
 \begin{gathered}
 j(a)=(0,a),\qquad\varepsilon(b,z)=z,\\
 s_1(g)=((0,1),g),\qquad
 s_2(g)=((r_\ell(g),1),g).
 \end{gathered}
\]
The coefficient action fixes $u(G_{\Q_p})$ and $j(\Q_p^\times)$ pointwise; $H$ is abelian. Thus both centralizations hold, $s_1u=s_2j\chi_p$, $\phi(g)=(r_\ell(g),1)$ and $\chi_p=\varepsilon u$.

A compact $p$-adic analytic group has an open pro-$p$ subgroup \cite[Theorem~19.11, p.~96]{PstragowskiLocal}. Every continuous map from a pro-$p$ group to $\Z_\ell$ is zero: each finite image is both a $p$-group and an $\ell$-group. Hence every pro-$p$ subgroup of the displayed compact image lies in $\{0\}\times\Z_p^\times$, which has empty interior. There is no compatible $p$-adic analytic structure. Each finite Artin level is analytic; the obstruction arises on taking the infinite inverse limit.
\end{proof}

\section{A non-geometric Hodge pair}
This section constructs the third assertion of Theorem~\ref{thm:other} with a specified real parameter, and proves its exclusion from the geometric image.

For any real $a>0$, take $V=\Q e_1\oplus\Q e_2$ with
\[
 V^{3,0}=\C(e_1+ia e_2),\qquad
 V^{0,3}=\C(e_1-ia e_2),\qquad
 f=\operatorname{diag}(1,-1).
\]
The two lines are conjugate and complementary, and $f$ exchanges them. A rational alternating form is
\[
 E=\begin{pmatrix}0&1\\-1&0\end{pmatrix},\qquad
 Q(x,y)=x^{\mathsf t}Ey.
\]
It vanishes on each of the two equal-type lines. We use the polarization convention that
$(-1)^{n(n-1)/2}i^{p-q}Q(z,\bar z)>0$ for nonzero $z\in V^{p,q}$ of weight $n$.
For $z=e_1+ia e_2$ one has $Q(z,\bar z)=-2ia$, so this expression is $2a>0$ in weight three. The conjugate line gives the same positive value. This proves polarizability, and interpreting $Q$ as the coefficient of a rational generator of $\Q(-3)$ makes it a Tate-valued Hodge morphism. Finally $f^{\mathsf t}Ef=-E$, exactly the sign of the involution on that Tate target.

To choose $a$, enumerate deterministic Turing machines as follows. The tape alphabet is $0,1,\#,\square$ in that order. For each $r\ge1$, take states $0,\ldots,r$, with initial state $0$ and halt state $r$. Enumerate all tables for the $4r$ nonhalting state-symbol pairs lexicographically, with instruction entries ordered as output symbol, next state, and move left, stay, right. Order the tables first by $r$. On input $n$, the tape contains the binary expansion of $n$ between two $\#$ symbols, with the head on its first digit. This fixes a sequence $(T_n)_{n\ge1}$ containing all machines of this model. Set
\[
 \mathcal H=\{n:T_n\text{ halts on input }n\},\qquad
 a=1+\sum_{n\ge1}\frac{\mathbf1_{\mathcal H}(n)}{4^n}.
\]

\begin{proposition}\label{prop:nong}
The polarized Hodge pair $(V,f)$ with this parameter is not the realization of a pure motive defined by algebraic correspondences over $\C$.
\end{proposition}
\begin{proof}
The number $a$ belongs to $[1,4/3]$ and is noncomputable. The diagonal halting set $\mathcal H$ is undecidable: a purported decider would allow a machine which halts on input $n$ exactly when the decider says that $T_n(n)$ does not halt, contradicting its answer at that machine's own index. If $a$ were computable, its base-four digits in $\{0,1\}$ would decide $\mathcal H$. After recovering the first $n-1$ digits, multiply the remaining tail by $4^n$. It lies in $[0,1/3]$ or $[1,4/3]$ according as the next digit is zero or one. An approximation of error less than $1/6$ distinguishes the two intervals, giving the forbidden decider.

Suppose that $V$ were geometric in the stated sense. A finite presentation uses finitely many smooth projective varieties, algebraic correspondences, rational projectors and Tate twists. Descend this data to a finitely generated field over $\Qb$ and spread it over a smooth integral finite-type $\Qb$-scheme $S$. After shrinking, the varieties are smooth projective families and the cycle actions are horizontal endomorphisms of their rational cohomology local systems. Idempotence holds on the connected component containing the given point, because it is an equality of locally constant endomorphisms. The projector images give a rational variation of Hodge structures; the Hodge ranks are locally constant.

For the rank-two weight-three image, $F^3=F^2=F^1$. Griffiths transversality yields
\[
                  \nabla F^3\subset\Omega^1_S\otimes F^2
                            =\Omega^1_S\otimes F^3.
\]
Thus $F^3$ is horizontal, as is its conjugate, and parallel transport preserves the rational Hodge structure. This is the rigidity mechanism for Hodge structures with gaps discussed by Totaro \cite[pp.~1--2]{Totaro}. Every analytic connected component of the smooth $\Qb$-scheme contains a $\Qb$-point: in an etale coordinate chart, the inverse-function theorem and density of $\Qb$ in $\C$ give such points arbitrarily close to any point. Parallel transport along a path consequently identifies $V$ with the realization of a motive over $\Qb$.

For a motive over $\Qb$, comparison entries in algebraic de Rham and rational Betti bases are periods, multiplied by integer powers of $2\pi i$ for Tate twists. Yoshinaga's theorem says their real and imaginary parts are computable \cite[Theorem 18, p.~11]{Yoshinaga}. The computable complex numbers form a field: to invert a known nonzero element, first find a positive rational lower bound for its modulus. It follows by solving the finite linear systems defining $F^3$ that its coordinate ratio is computable in every rational basis. A rational change of basis acts by an invertible fractional-linear transformation and preserves this property. The ratio $ia$ would therefore be computable, contrary to the construction.
\end{proof}

The qualification ``defined by algebraic correspondences'' is used in spreading the projector. The argument does not replace an algebraic projector by an arbitrary Hodge projector over an unspecified complex coefficient field. It also excludes a real geometric realization, since complex base change would give one over $\C$.

\subsection{An algebraic parameter outside the abelian tensor category}
The noncomputable example excludes every geometric realization of the stated kind. A different, algebraic parameter gives a simpler exclusion from the tensor category generated by abelian varieties.

\begin{proposition}\label{prop:nonabelian-pair}
On $V=\Q e_1\oplus\Q e_2$ put
\[
 V^{3,0}=\C(e_1+i\sqrt[3]{2}\,e_2),\qquad
 V^{0,3}=\C(e_1-i\sqrt[3]{2}\,e_2),\qquad
 f=\operatorname{diag}(1,-1).
\]
This is a polarizable rational Hodge structure with an anti-Hodge involution, but it does not belong to the tensor category generated by $H^1$ of complex abelian varieties.
\end{proposition}
\begin{proof}
The polarization calculation used above applies to every positive parameter $a$, hence to $a=\sqrt[3]{2}$. An endomorphism $\left(\begin{smallmatrix}r&s\\t&u\end{smallmatrix}\right)$ preserves the $(3,0)$ line only if
\[
             t+(u-r)ia+s\sqrt[3]{4}=0.
\]
Taking real and imaginary parts gives $u=r$ and $t+s\sqrt[3]{4}=0$. Rationality forces $s=t=0$, so $\End_{\mathrm{HS}}(V)=\Q$. The Mumford--Tate group cannot be commutative: otherwise its Lie algebra lies in this endomorphism algebra, whereas the weight scaling and nontrivial Hodge circle give two independent real tangent directions.

For completeness, the obstruction to abelian generation is a Lie-algebra weight obstruction. The Lie algebra of the Mumford--Tate group of $H^1(A)$ is a sub-Hodge structure of $\End H^1(A)$, whose types are $(-1,1),(0,0),(1,-1)$. Passing to an object in its tensor category gives a quotient of Mumford--Tate groups and, in characteristic zero, a surjection of their Lie algebras; it cannot add new Hodge types. But $\End V$ has only $(-3,3),(0,0),(3,-3)$. If $V$ belonged to that category, its Mumford--Tate Lie algebra would have only type $(0,0)$. The Hodge homomorphism would act trivially on it by conjugation. By the defining rational Zariski closure of this homomorphism, the connected Mumford--Tate group would be commutative, a contradiction. This is the mechanism of Totaro's Theorem 4.1 and Corollary 4.2 \cite[p.~18]{Totaro}.
\end{proof}

\subsection{Real forms and the Artin sign}\label{sec:real-artin-sign}
An anti-Hodge involution means a rational map $f$ with $f^2=1$ and $fV^{p,q}=V^{q,p}$. It should be distinguished from coefficient conjugation $c$ on $V_\C$. For a real motive their product gives the real structure of de Rham cohomology: under comparison,
\[
 D_\R=\alpha\bigl(V^+_\R\oplus iV^-_\R\bigr).
\]
Indeed a vector $x+iy$ is fixed by $fc$ exactly when $fx=x$ and $fy=-y$. Both $f$ and $c$ interchange the Hodge labels, so $fc$ preserves every Hodge summand. Projection onto those summands, and hence the normalized splitter $\beta$, is defined over $\R$. Thus $\R(\gamma)=\R(\alpha)$, since $\gamma=(\beta\alpha)^{-1}$ and $\beta$ has real entries.

Let $A$ be the one-dimensional pair of Hodge type $(0,0)$ with $f=-1$, and let $T=\Q(-1)$ with its geometric involution $f=-1$. In the category of polarizable rational Hodge structures with an anti-Hodge involution, the objects with $f=1$ are exactly finite direct sums of $(A\otimes T)^p$, $p\in\Z$. In fact $f=1$ forces all types to be $(p,p)$; on a type-$(p,p)$ summand the Tate involution is $(-1)^p$, so the multiplicity space must carry the sign $(-1)^p$. The converse is immediate. Proposition~\ref{prop:real-fullness} below transports this classification to real abelian absolute-Hodge motives, and to the indicated real homological motives under its explicit Hodge-conjecture hypothesis.

The failure of an ``even Tate powers'' description already has an algebraic example. On $H^2(\mathbf P^1_\R\times\Spec\C)$ the deck transformation is $d=\left(\begin{smallmatrix}0&1\\1&0\end{smallmatrix}\right)$ and real Frobenius is $-d$. The algebraic projector $(1-d)/2$ has image $A\otimes T$, with trivial real Frobenius and weight two. Sums of powers of $\Q(2)$ have weights divisible by four and cannot contain it.

There is a related centrality criterion, transported to those same motivic categories by Proposition~\ref{prop:real-fullness}. Let $G$ be the smallest rational algebraic group containing the Hodge homomorphism and $f$. Then the following are equivalent: $f$ is central in $G$; every Hodge type is diagonal; the pair is a sum of $T^p$ and $A\otimes T^p$; and $G$ is commutative. Centrality and $fh(z)f^{-1}=h(\overline z)$ force $p=q$. Diagonal types decompose into the two eigenspaces of $f$, giving the stated sums. Their group is a quotient of $\mu_2\times\Gm$, proving the remaining implications.

Finally, if a polarization $q$ of weight $n$ is compatible with $f$, and $C=h(i)$ is the Weil operator, then its positive real form is invariant under $f$. In the convention $B(x,y)=(2\pi i)^nq(x,Cy)$, one has
\[
 q(fx,fy)=(-1)^nq(x,y),\qquad fCf^{-1}=C^{-1},\qquad C^{-1}=(-1)^nC.
\]
Substitution gives $B(fx,fy)=B(x,y)$. The sign on the Tate target is necessary for this equality.

\subsection{Real descent and twists of the Betti functor}
The realization statements just used have a precise motivic scope. Let $\mathcal A_{\R}$ be the rational category for absolute Hodge cycles generated by real abelian varieties and Artin motives. A \emph{Hodge pair} means a polarizable rational Hodge structure with a rational involution $f$ exchanging $H^{p,q}$ and $H^{q,p}$; morphisms are Hodge morphisms commuting with $f$.

\begin{proposition}\label{prop:real-fullness}
Real Betti realization is fully faithful from $\mathcal A_{\R}$ to Hodge pairs. Its degree-one image is exactly the pairs of type $\{(1,0),(0,1)\}$, and its whole image is their rigid, additive, idempotent-complete tensor closure together with the sign Artin line. For a tensor generator with Hodge homomorphism $h$, its group is
\[
                 G=\operatorname{MT}(V)\cup f\operatorname{MT}(V).
\]
In particular, $G^0=\operatorname{MT}(V)$ and $\pi_0(G)$ is trivial or the constant group of order two. For any subcategory of real homological motives, the same fullness and group assertions hold under the Hodge conjecture for all products and twists required in that subcategory; this last assertion does not characterize its geometric essential image.
\end{proposition}
\begin{proof}
A Hodge morphism between the indicated tensor constructions over $\C$ is an absolute Hodge class by Deligne's theorem \cite[Theorem~2.11, author p.~19]{DeligneHodgeCycles}. Commuting with $f$ is exactly invariance under conjugation; the definition of absolute Hodge cycles over $\R$ therefore gives the morphism over $\R$. Betti realization is faithful. In the homological category, the stated Hodge conjecture instead supplies an algebraic correspondence $z$ over $\C$; replacing it by $\tfrac12\operatorname{Tr}_{\C/\R}(z)$ descends an invariant class. Projectors commute with this averaging. These arguments prove fullness between geometric objects, without asserting that every Hodge pair is geometric.

For degree one, use the rational Riemann theorem \cite[Theorem~6.20 and Proposition~6.21, p.~65]{DeligneMilne}. On the dual homology space let $J$ be the complex structure and choose a lattice $\Lambda_0$. Then $\Lambda=\Lambda_0+f\Lambda_0$ is $f$-stable. For a rational Riemann form $E_0$, put
\[
                E(x,y)=E_0(x,y)-E_0(fx,fy).
\]
Since $fJ=-Jf$, one has
$E(x,Jx)=E_0(x,Jx)+E_0(fx,Jfx)>0$ for $x\ne0$, and $f^*E=-E$. After multiplication by a positive integer, $E$ is integral on $\Lambda$. Thus $(\Lambda\otimes\R,J)/\Lambda$ is an abelian variety with an antiholomorphic group involution induced by $f$. As a holomorphic map to its conjugate, this involution is algebraic; its square is the identity. Projective descent gives a real abelian variety realizing the pair. Fullness now gives the stated tensor-generated image.

Put $H=\operatorname{MT}(V)$. The connected real torus defining $h$ lies in $H^0_{\R}$, so the minimality of the rational algebraic group $H$ gives $H=H^0$. The identity $fh(z)f^{-1}=h(\bar z)$ shows that $f$ normalizes $H$. Hence $K=H\cup fH$ is an algebraic subgroup with identity component $H$. In every mixed tensor construction $T(V)$, fullness identifies
\[
 T(V)^G=\{t:t\text{ has Hodge type }(0,0),\ ft=t\}=T(V)^K.
\]
These invariants determine the groups here. Indeed polarizable Hodge structures are semisimple; a Hodge retraction onto a subpair becomes an $f$-equivariant retraction after averaging it with $fpf$. Thus the Hodge-pair category is semisimple. Both $H$ and its finite extension $K$ are reductive in characteristic zero. Every representation of $K$ is a direct summand of sums of mixed tensor words in the faithful representation $V$. Equality of the displayed invariants gives equality of all Hom spaces and idempotents on those words. Tensor reconstruction \cite[Theorem~2.11]{DeligneMilne} identifies $K$ with $G$. If the second component exists, multiplication by the rational point $f$ identifies it with $G^0$ as a rational scheme.
\end{proof}

Consequently the kernel and centrality classifications in Section~\ref{sec:real-artin-sign} apply to $\mathcal A_{\R}$ itself. For homological motives they retain exactly the Hodge-conjecture hypothesis in Proposition~\ref{prop:real-fullness}.

Let now $V$ be any neutral rational tensor fibre functor and let $f\in G(\Q)$ be a tensor involution, where $G=\Aut^\otimes(V)$. Define
\[
                  W(X)=V(X)^+\oplus iV(X)^-
                       \subset V(X)\otimes_{\Q}\Q(i).
\]
This is the fixed space of the semilinear involution $f\otimes c$.

\begin{proposition}\label{prop:betti-twist}
Coefficient multiplication makes $W$ a rational tensor fibre functor. Its comparison torsor has cocycle $f^{\epsilon(\sigma)}$, where $\sigma(i)=(-1)^{\epsilon(\sigma)}i$. If $i\notin F$, its $F$-points are described, in the inverse ambient-matrix coordinate, by
\[
       \{b\in G(F(i)):b^{-1}\bar b=f\}.
\]
It is trivial over every field containing $i$. For the sign Artin category it is nontrivial over $\Q$ and $\R$; for the Tate category and the category generated by a real non-CM elliptic $h^1$, it is trivial over $\Q$. For a real CM elliptic $h^1$ it is nontrivial over $\R$.
\end{proposition}
\begin{proof}
The four tensor components multiply with factors $1,i,i,-1$. They give exactly $W(X\otimes Y)$; associativity, unit and symmetry are inherited from the ambient multiplication. Exactness and faithfulness follow after the field extension $\Q(i)/\Q$, where the inclusions identify $W$ with $V$. Conjugating that identification changes it by $f$, giving the stated cocycle. If $a$ is the ambient matrix of a tensor map $V_{F(i)}\to W_{F(i)}\hookrightarrow V_{F(i)}$, descent is $f\bar a=a$. With $b=a^{-1}$ this is $b^{-1}\bar b=f$. This is the usual fibre-functor torsor construction \cite[Theorem~3.2]{DeligneMilne}.

For the Artin line, $G=\mu_2$ and $f=-1$. Every $b\in\mu_2(\Q(i))$ or $\mu_2(\C)$ is fixed by conjugation, so the equation cannot hold. Its torsor is $\operatorname{Spec}\Q(i)$. For the Tate line, $G=\Gm$ and $b=i$ works. On all Tate powers the tensor isomorphism is $e_n\mapsto i^ne_n$; there is no order-two tensor relation on the Tate object. For a non-CM elliptic curve, $G=\GL_2$ and $f$ has rational eigenvalues $1,-1$, so $b=\operatorname{diag}(1,i)$ works in an eigenbasis.

For a CM elliptic curve, the trace-zero line of its imaginary quadratic endomorphism field is a Hodge line of type $(0,0)$ inside $\End(h^1)$. Real Frobenius acts by $-1$ on it. Proposition~\ref{prop:real-fullness} identifies it with the sign Artin object, so the Artin obstruction prevents a real tensor comparison. Equivalently, the CM group has two components and $f$ lies in the nonidentity component; a coboundary over $\C/\R$ has trivial image in that constant component group.
\end{proof}

The objectwise additive map $v_++v_-\mapsto v_++iv_-$ is generally not a tensor map: on two negative vectors its tensor square is minus its value on their tensor product. The valid tensor maps above specify all tensor powers consistently.

There is also a distinction between the torsor and its distinguished point. For real de Rham comparison,
\[
             \bar\alpha=\alpha f,\qquad
             \alpha\in P(\R)\ \Longleftrightarrow\ f=1.
\]
The first equation is the conjugation comparison already proved; the second follows by cancellation. It identifies the real torsor class with the cocycle $f$. If $\Q(-1)$ belongs to the generated category, its comparison coefficient is a nonzero real multiple of $(2\pi i)^{-1}$, so $\R(\alpha)=\C$. The Tate torsor can nevertheless have a real point, as Proposition~\ref{prop:betti-twist} shows; that point is not its distinguished period point.

\subsection{Deforming the entire Hodge realization}

An individual Hodge structure need not be geometric, as above. There is also
a failure for a tensor-compatible Hodge assignment on the entire geometric
Betti diagram. The underlying integral diagram can stay fixed while its Hodge
structures change. We first give the simultaneous lifting fact needed when
the category has uncountably many objects.

Here the integral diagram consists of the geometric cohomology lattices and
integral algebraic correspondence maps; its rational tensor closure is the
motivic category. As in Section~\ref{sec:integral}, no additive integral
functor on a nonzero $\Q$-linear category is meant.

\begin{lemma}\label{lem:prolie-lift}
Let $q:G\to\SL_{2,\R}$ be a faithfully flat homomorphism of affine group
schemes over $\R$, without a finite-type assumption on $G$. Suppose that an
involution $\theta$ of $G$ induces an involution $\theta_0$ of $\SL_2$.
Every $\theta_0$-fixed $Y\in\mathfrak{sl}_2(\R)$ has a $\theta$-fixed lift
$X\in\Lie(G)(\R)$. Its exponentials in the finite algebraic quotients of $G$
are compatible and define a real point of $G$.
\end{lemma}

\begin{proof}
An affine group is an inverse limit of algebraic quotient groups with
faithfully flat transition maps; every map to an algebraic group factors
through one of them \cite[Proposition~2.6 and Corollary~2.7, pp.~19--20]{DeligneMilne}.
Take a cofinal system $G_i$ through which $q$ factors and which is stable under
$\theta$. One obtains the latter condition by including a finite set of
matrix coefficients together with its image under $\theta$.
The maps $G_i\to\SL_2$ are faithfully flat. Their kernels are smooth in
characteristic zero, and the geometric fibres are translates of those
kernels. The fibrewise smoothness criterion makes the maps smooth
\cite[Tags 047N and 01V8]{StacksE}, so their differentials are surjective.
A lift of $Y$ in
$\Lie(G_i)(\R)$ becomes invariant after averaging it with its
$d\theta_i$-image.

Choose coordinates in the finite-dimensional spaces $\Lie(G_i)(\R)$.
Compatibility, invariance, and the prescribed image $Y$ are linear equations
in those coordinates. Every finite subset of the equations is solvable at
a common dominating index. To see that all the equations have a solution,
regard their left sides as vectors in the free vector space on the coordinate
symbols. Finite consistency says that their prescribed right sides define a
well-defined linear functional on their span. Extend that functional to the
whole free space by extending a basis. Its values on the coordinate symbols
are the required solution. This argument applies to any cardinality of the
index set and gives a compatible invariant family $X_i$.

For every $i$, the matrix Lie exponential $\exp(X_i)$ belongs to $G_i(\R)$.
Homomorphisms of real algebraic groups commute with the Lie exponential,
so these elements are compatible and define a point of $G(\R)$. They are
fixed by $\theta$ because the $X_i$ are fixed. The same reasoning applies to
$\exp(tX)$ for every real $t$.
\end{proof}

\begin{proposition}\label{prop:global-hodge-deformation}
On the geometric integral Betti diagram of pure motivated motives over $\Q$
there is a tensor-compatible assignment of polarizable integral Hodge
structures which preserves weights, every specified motivic tensor and
polarization, Tate objects, real Frobenius, and the unchanged finite
realizations, but is induced by no complex place of $\Q$.

The same assertion holds for the entire pure motivated category over $\C$,
with compatibility between conjugate objects and with real Frobenius on
objects descended from $\R$; no embedding $\C\hookrightarrow\C$ induces the
new assignment. Both constructions also apply to the absolute-Hodge abelian
subcategory. Its new degree-one objects are realizable by complex abelian
varieties, and by real abelian varieties when the original real descent datum
is present.
\end{proposition}

\begin{proof}
Let $G=\Aut^\otimes(V)$ for the rational Betti functor of the chosen category,
and let $\varepsilon$ be its character on $\Q(-1)$. Use the rational elliptic
curve
\[
 E_*:\quad y^2=4x^3-4x-1,\qquad j(E_*)=110592/37.
\]
Its nonintegral $j$ excludes CM. The motivated group of $H^1(E_*)$ is
$\GL_2$, by the argument in Section~4; the same conclusion holds in the
absolute-Hodge abelian category. The determinant is the Tate character.
The full tensor subcategory theorem gives a faithfully flat quotient
$G\to\GL_2$ \cite[Proposition~2.21(a), p.~26]{DeligneMilne}, and therefore
\[
 K:=\ker\varepsilon=G\mathbin{\times}_{\GL_2}\SL_2
 \longrightarrow\SL_2
\]
is faithfully flat.

Over $\Q$, real Frobenius is a rational tensor involution $f$ of $V$ and
acts on $K$ by conjugation. Over $\C$, write $c^*M$ for the conjugate motive
and let $j_M:V(M)\to V(c^*M)$ be the rational tensor map induced by
conjugating complex points. It has $j_{c^*M}j_M=1$, and
\[
 \theta(g)_M=j_M^{-1}g_{c^*M}j_M
\]
defines an involution of $G$ preserving $K$. On the elliptic quotient this
is again conjugation by its real Frobenius $f_E$. For an object with real
descent, the equation $\theta(g)=g$ is exactly commutation with that object's
real Frobenius, by naturality with the descent isomorphism.

Choose a rational $f_E$-eigenbasis in which $f_E=\operatorname{diag}(1,-1)$.
The holomorphic line has the form $\C(1,iy)$ with $y>0$, after choosing the
sign of the second basis vector. Put
\[
 b=\sqrt[3]{2},\qquad a=\sqrt{y/b},\qquad
 Y=\operatorname{diag}(\log a,-\log a).
\]
The matrix $Y$ is fixed by the induced involution. Lemma~\ref{lem:prolie-lift}
gives an invariant lift in $\Lie K$ whose exponential is a real tensor
automorphism $g\in K(\R)$ with
$g_E=\operatorname{diag}(a,a^{-1})$.

Leave all rational Betti spaces and geometric integral lattices unchanged
and put
\[
 F'^pV(M)_\C=g_M F^pV(M)_\C.
\]
Reality of $g_M$ gives the conjugation rule for the new Hodge decomposition.
Tensor compatibility and naturality preserve every motivic morphism and
tensor; the motivic weight projectors remain fixed. Tate objects remain
fixed because $g\in K$. For a specified motivic polarization
$q:M\otimes M\to\Q(-n)$ one has $q(gx,gy)=q(x,y)$. The new Weil operator
is $gCg^{-1}$, so each polarization inequality is the original inequality
evaluated on $g^{-1}x$. It is therefore still positive. The invariant-lift
condition retains the asserted conjugation and real-Frobenius structures.
All original integral morphisms remain Hodge morphisms by the same
naturality. The finite realizations and their comparisons with the unchanged
Betti lattices have not been altered.

The new elliptic line has ratio $ib$. This number is algebraic and is not
imaginary quadratic: its square is $-\sqrt[3]{4}\notin\Q$, whereas a purely
imaginary quadratic number has rational square. Passing from the rational
eigenbasis to an integral lattice basis is a fractional-linear change over
$\Q$, which preserves the generated number field and hence the nonquadratic
property. Choose the integral orientation giving a ratio in the upper
half-plane. Schneider's theorem \cite[\S1, result IIc, p.~5]{Schneider}
then makes the new elliptic $j$-invariant transcendental. The retained
positive Riemann form ensures that the resulting torus is an elliptic curve.
An integral Hodge isomorphism to the genuine realization of the labelled
$E_*$ would preserve $j$, contradicting its rational value. Every embedding
of either $\Q$ or $\C$ into $\C$ fixes this rational equation. Thus neither
category's assignment is induced by such an embedding.

For degree-one objects in the abelian category, the Riemann theorem realizes
the new polarized integral Hodge structures by abelian varieties. When a real
involution is present, the real-descent argument of
Proposition~\ref{prop:real-fullness} realizes it over $\R$. This is an
objectwise geometric realization, compatible with the Hodge morphisms; it
does not identify the new objects with the original labelled algebraic
models or supply a complex place.
\end{proof}

One may also retain the algebraic de Rham functor together with a
\emph{chosen} tensor comparison. If the original comparison is
$\alpha_M:V(M)_\C\to D(M)_\C$, set
\[
 \alpha'_M=\alpha_Mg_M^{-1}.
\]
Then $(\alpha'_M)^{-1}F_D=F'_M$. This comparison preserves the Tate and real
structures by the identities already proved. It is not claimed to be the
distinguished comparison obtained by integration. Under the forgetful map to
the plain integral diagram the example has its original, unchanged class.
Consequently it gives the stated failure for the enriched Hodge target and
does not classify the image of plain integral comparison classes or of any
target with further unspecified conditions.

\section{Integral comparisons and places}\label{sec:integral}
This section proves the two assertions about integral comparisons and identifies the exact distinction between finite compatibility and a simultaneous lattice comparison.

An integral Betti functor here means the diagram of smooth projective varieties, integral cohomology modulo torsion, and integral algebraic correspondences. We also use its covariant restriction $H_1(-,\Z)$ on abelian varieties with homomorphisms. It cannot literally mean an additive functor from a nonzero $\Q$-linear category to free $\Z$-modules: the invertible morphism $n\cdot1$ would have to act as multiplication by $n$, which is not invertible for $n>1$. Duality of free lattices transfers the obstructions below between $H_1$ and $H^1$.

\subsection{One path, two indistinguishable places}
Put $F=\Qb(t)$ and $k=\overline F$. Choose an embedding $v:k\hookrightarrow\C$ fixing $\Qb$ and sending $t$ to $e$. Use the polygonal path
\[
                e\longrightarrow e+ie\longrightarrow e+1+ie
                   \longrightarrow e+1
\]
in the complex $t$-line. It avoids $\Qb$: on each vertical segment its real part is $e$ or $e+1$, and on the horizontal segment its imaginary part is $e$. Real and imaginary parts of algebraic numbers are algebraic.

Every element of $k$ is an algebraic function of $t$ with finite branch and pole set in $\Qb$. The branch specified by $v$ at $e$ continues uniquely along the path. Its endpoint is finite. Define $w$ to be the endpoint-value map. Analytic continuation preserves sums, products and algebraic relations. A nonzero algebraic function has zeros only over algebraic values of $t$, so it does not become zero at $e+1$. Thus $w$ is an embedding fixing $\Qb$ and sending $t$ to $e+1$. The two embeddings are distinct and not conjugate, because both fix $i$.

For a finite collection of varieties and correspondences over $k$, descend to a finite extension $F_0/F$. Normalize an open subset of the $t$-line in $F_0$, then delete its branch points and the images of the finitely many bad-reduction loci. These excluded points are all algebraic. The displayed path lifts from the sheet specified by $v$ to that specified by $w$. Smooth proper transport yields integral cohomology isomorphisms compatible with all the correspondences, cup products and K\"unneth maps in the collection. Enlarging the collection or $F_0$ yields the same maps, by uniqueness of path lifting and functoriality of cohomology local systems. The finite comparisons therefore give one natural tensor comparison on the whole diagram. It extends to specified projector lattices, real scalar extensions, and finite-coefficient comparisons. The latter compatibility is the Riemann-existence identification between topological and etale paths. This proves the first part of assertion (4) in Theorem~\ref{thm:other}.

\subsection{Finite comparisons with no global integral comparison}
The opposite phenomenon occurs for two embeddings agreeing on $F$. Consider the Legendre curve
\[
                     E/F:\quad y^2=x(x-1)(x-t).
\]
The local integral monodromy around $0$ has matrix
\[
                    T_0=\begin{pmatrix}1&2\\0&1\end{pmatrix}
\]
in a suitable oriented homology basis \cite[Example 5.2, pp.~16--17]{LocalSystems}. Riemann existence identifies the torsion monodromy with algebraic Galois action on the torsion fields; invariance under algebraically closed characteristic-zero base extension is given by \cite[Theorem 1.1]{Landesman}. The image on the full Tate module contains the closure of the powers of $T_0$, namely
\[
                  u(2a)=\begin{pmatrix}1&2a\\0&1\end{pmatrix}
                         \quad(a\in\widehat\Z).
\]
Choose the particular $a$ with compatible residues
\[
            a\bmod n!=s_n:=\sum_{m=1}^{n-1}m!\pmod{n!}.
\]
They are compatible since $s_{n+1}-s_n=n!$. For $n\ge4$, one has $0<s_n<n!/2$, and $s_n\to\infty$. The least residue of a fixed nonnegative integer is eventually constant, while that of a fixed negative integer eventually exceeds $n!/2$. Thus $a\notin\Z$.

Choose an automorphism of the torsion field acting as $u(2a)$ and extend it to $\sigma\in\Gal(k/F)$. Set $w=v\circ\sigma$. The embeddings agree on $F$, but differ on the torsion field and are not conjugate. For every cyclic subgroup $H\subset E[n]$, the quotient isogeny $f_H:E\to E/H$ is defined over $k$. If $\Lambda=H_1(E_v,\Z)$, analytic uniformization identifies its homology map with a lattice inclusion
\[
                 \Lambda\subset\Lambda_H\subset n^{-1}\Lambda,
                    \qquad\Lambda_H/\Lambda=H_v.
\]
For $w$ the corresponding subgroup is $u(2a)H_v$.

A natural integral comparison would act on $\Lambda$ by $A\in\GL_2(\Z)$. Naturality with each $f_H$ forces $A\Lambda_H=\Lambda_{H_w}$: the rational extension of its map on $H_1(E/H)$ is already determined on the finite-index sublattice $\Lambda$. Reducing modulo every $n$, the matrix $u(2a)^{-1}A$ stabilizes every cyclic subgroup of $(\Z/n)^2$. Stabilizing the summands generated by $e_1,e_2,e_1+e_2$ forces a scalar unit. The scalars are compatible as $n$ varies, so
\[
                         A=u(2a)\lambda I,\qquad
                            \lambda\in\widehat\Z^\times.
\]
The two diagonal entries of $A$ are the same integer $\lambda$. Its determinant is $\lambda^2=\pm1$, hence $\lambda=\pm1$. Then its upper-right entry makes $2a\in\Z$. An element of $\widehat\Z$ whose double is integral is integral, because that integer is even in $\Z_2$ and multiplication by two is injective on $\widehat\Z$. This contradicts $a\notin\Z$.

Nevertheless each finite subdiagram admits a comparison. It descends to a finite extension of $F$, the function field of a connected smooth curve over $\Qb$. Spread the finite data over a suitable open set. The two embeddings give complex points of that open curve; a path joins them and yields the finite comparison by smooth proper transport. This proves the second part of assertion (4).

The failure of an inverse limit can already be seen on paths. Put $k'=\bigcup_n\Qb(t^{1/n!})$ and $u_n=t^{1/n!}$. The assignments
\[
 v(u_n)=\exp(1/n!),\qquad
 w(u_n)=\exp(1/n!+2\pi i s_n/n!)
\]
are compatible embeddings. A compatible system of paths between their values would project to a loop with integer winding number congruent to $s_n$ modulo every $n!$. No such integer exists. Each individual path set is nonempty. Replacing winding numbers by $\widehat\Z$ admits the compatible element $a$, but completion has changed the coefficient object and no longer guarantees preservation of an integral lattice.

\subsection{An ideal-class obstruction over an algebraic field}
Let $K_0=\Q(\sqrt{-5})$, $\mathcal O=\Z[\sqrt{-5}]$ and $I=(2,1+\sqrt{-5})$. The ideal is invertible, has norm two, and is nonprincipal: a generator $a+b\sqrt{-5}$ would satisfy $a^2+5b^2=2$, which has no integral solution. Moreover $I^2=(2)$, by multiplying the generators and comparing norms.

The curves $\C/\mathcal O$ and $\C/I$, with their labelled $\mathcal O$-actions, have algebraic models. The fundamental theorem of complex multiplication, in the ideal-class form of \cite[Theorem 3.13]{Milne}, gives a Galois automorphism fixing $K_0$ carrying one to the other. The reflex field is $K_0$ and the reflex norm is the identity. Either inverse-ideal convention gives the same class because $[I]^2=1$.

For the two corresponding embeddings, integral Betti homology of the same labelled algebraic CM object has $\mathcal O$-modules $\mathcal O$ and $I$. An integral natural comparison must commute with its integral endomorphisms, so would be an $\mathcal O$-linear isomorphism. The image of $1$ would generate $I$, a contradiction. Rationally both become the one-dimensional $K_0$-vector space $K_0$, and at every prime the invertible ideal is principal. Thus rational and all locally integral comparisons exist. Taking $\Z$-duals gives the same obstruction in cohomology, since the nontrivial ideal class is its own inverse. This proves assertion (5) and refutes an unrestricted integral comparison as in cote 18, PDF page 55.

\subsection{What labelled Hodge filtrations recover}\label{subsec:enrichment}
The preceding integral examples forget the Hodge filtrations. Retaining those filtrations changes what a comparison can identify.

\begin{proposition}\label{prop:enriched-embeddings}
Let $\sigma,\tau:k\hookrightarrow\C$ be field embeddings. Suppose their labelled integral Betti diagrams have a natural isomorphism preserving the Hodge filtrations for every elliptic curve over $k$. Then $\sigma=\tau$. If simultaneous exchange of Hodge labels is allowed, the conclusion is equality up to complex conjugation. For a finitely generated field, a finite list of elliptic curves suffices.
\end{proposition}
\begin{proof}
An integral Hodge isomorphism on $H^1$ dualizes to an isomorphism of integral first homology preserving its complex structure. The complex elliptic torus, and hence its $j$-invariant, is recovered from this lattice. For each $a\in k$, choose an elliptic curve over $k$ with $j=a$: for $a\ne0,1728$, the equation with $g_2=3a/(a-1728)$ and $g_3=a/(a-1728)$ does so; the two exceptional values have their usual equations. The labelled comparison gives $\sigma(a)=\tau(a)$ for every $a$. For finitely generated $k$, do this only for a finite set of field generators. Exchanging all Hodge labels replaces one complex structure by its conjugate and gives the second assertion.
\end{proof}

In contrast, over $k=\overline{\Qb(t)}$ no fixed finite geometric presentation detects every embedding. Its equations, correspondence cycles and projectors descend to a finite extension $F/\Qb(t)$. There is a nontrivial automorphism of $k/F$: the one-variable function field $F$ has a discrete valuation, and an element of odd valuation is not a square, so $F$ is not algebraically closed. Extend a nontrivial automorphism of a finite normal extension to the algebraic closure. Two embeddings differing by this automorphism give the same finite presentation. The complete labelled Hodge diagram distinguishes them by Proposition~\ref{prop:enriched-embeddings}.

Even degree-zero objects can distinguish real embeddings. If a number field has real embeddings $\sigma_1,\ldots,\sigma_r$, choose a primitive element $\theta$ and order its distinct real values. For each gap choose a rational $q_j$ between the adjacent values and consider the quadratic algebra $k(\sqrt{\theta-q_j})$. At a real embedding where $\theta-q_j>0$, conjugation fixes its two geometric points and has trace two on $H^0$; where $\theta-q_j<0$, it interchanges them and has trace zero. The resulting $r-1$ sign tests distinguish all real embeddings. The construction is vacuous, as it should be, for $r\le1$.

\subsection{A congruence model for the compatibility obstruction}
A finite-stage comparison need not produce a common rational point even in an elementary matrix model. Choose $a\in\widehat\Z\setminus\Q$, compatible integral residues $a_n\bmod n!$, and put
\[
 U(a_n)=\begin{pmatrix}1&a_n\\0&1\end{pmatrix},\qquad
 C_n=U(a_n)\ker\bigl(\GL_2(\widehat\Z)\to\GL_2(\Z/n!\Z)\bigr).
\]
These are nested nonempty compact cosets, each containing an integral matrix. Their intersection is the single matrix $U(a)$, which is not rational. Indeed congruence modulo every $n!$ determines each profinite matrix entry. This proves the exact limitation of finite congruence tests, independently of any motivic reconstruction assertion.

\subsection{Relative frames and the three finite coefficient rings}
The rings
\[
 \widehat\Z=\prod_\ell\Z_\ell,\qquad
 \mathbf A_f=\widehat\Z\otimes_\Z\Q=\prod_\ell'\Q_\ell,
 \qquad R_{\rm all}=\prod_\ell\Q_\ell
\]
serve different purposes. The restricted product uses $\Z_\ell$ and makes $\widehat\Z$ compact and open. Its inclusion in the full product is continuous and dense, but is proper, since $(1/\ell)_\ell\notin\mathbf A_f$. It does not induce the restricted topology: a product neighbourhood restricts only finitely many coordinates and cannot be contained in $\widehat\Z$. In particular, $\prod_\ell\GL_m(\Q_\ell)$ has no compact open subgroup for $m>0$. Any open identity neighbourhood projects onto an entire noncompact factor outside finitely many primes, whereas the image of a compact set is compact. In contrast, $\GL_m(\widehat\Z)$ is compact and open in $\GL_m(\mathbf A_f)$. For an affine group $G/\Q$, an integral model away from finitely many primes defines the corresponding restricted product $G(\mathbf A_f)$; changing that model alters only finitely many integral subgroups and leaves this topology unchanged.

For an abelian variety and an embedding $v$ of its algebraically closed ground field into $\C$, finite-coefficient comparison gives
\[
 L_v=H^1(A_v(\C),\Z),\qquad
 L_v\otimes\Z_\ell\simeq H^1_{\rm et}(A,\Z_\ell)=T_\ell.
\]
The input is SGA~4, Expos\'e XVI, Th\'eor\`eme~4.1, original p.~233 \cite{SGAComparison}, applied to $\Z/\ell^a$ and then to the inverse limit. The analytic torus has free integral cohomology, so its finite coefficient groups have surjective transition maps and no inverse-limit obstruction. Comparison preserves products and correspondences. The resulting map $L_v\to\prod_\ell T_\ell$ is injective, with dense image and completion isomorphic to the target: after choosing a basis this is the Chinese remainder theorem for $\Z^m$. Thus the integral lattice is not discrete in the finite adelic realization. A discrete diagonal lattice requires the archimedean component as well.

\begin{proposition}\label{prop:relative-frames}
Let $V,V_1,\ldots,V_n$ be tensor fibre functors on a neutral rational Tannakian category with a tensor generator. Put $G=\Aut^\otimes(V)$, $g=\dim G$ and $P_i=\Isom^\otimes(V,V_i)$. Then the diagonal right quotient is represented by
\[
 Q=(P_1\times\cdots\times P_n)/G
   \simeq\prod_{i=2}^n\Isom^\otimes(V_1,V_i),
 \qquad \dim Q=(n-1)g.
\]
For $q_\ell\in Q(\Q_\ell)$ one has
$\trdeg_\Q\Q(q_\ell)\le(n-1)g$. For a point over $R_{\rm all}$, the joint invariant supplied by evaluation is instead
\[
 \dim\im\bigl(\Q[Q]\longrightarrow R_{\rm all}\bigr)\le(n-1)g;
\]
the image need not be a domain.
\end{proposition}
\begin{proof}
The map $(p_i)\mapsto(p_ip_1^{-1})_{i\ge2}$ is invariant under the diagonal action. Conversely, choose $p_1$ fpqc-locally and put $p_i=q_ip_1$. Changing that choice changes exactly the common right factor. These constructions are mutually inverse as fpqc sheaves and prove representability by the displayed product of affine torsors. Each factor has dimension $g$. Evaluation at a local point has prime kernel; the fraction field of its image is $\Q(q_\ell)$, giving the local bound. Evaluation in a product ring has an arbitrary ideal as kernel, and a quotient of a finitely generated algebra cannot have larger Krull dimension.

For the domain distinction take $Q=\Gm$, $q_2=1$, and $q_\ell=2$ at the other primes. The coordinate algebra is $\Q[x,x^{-1}]/((x-1)(x-2))\simeq\Q\times\Q$. Conversely $q_\ell=\ell+1$ is locally rational at every prime, but no nonzero Laurent polynomial vanishes at all these distinct rational values. Its joint coordinate algebra is $\Q[x,x^{-1}]$, of dimension one. Thus neither joint genericity nor its dimension can be read off from the individual transcendence degrees.
\end{proof}

When $G$ is connected, equality in the local bound is precisely genericity over $\Q$ in the geometrically integral relative torsor. Coincident places give the identity frame and degree zero. Forgetting tensor constraints replaces $g$ by the upper bound $m^2$ for rank $m$; it does not prove equality for a prescribed tuple of places.

Retaining filtrations changes the relevant fibres. For two two-dimensional factors with equal determinants, the group
\[
 \Delta_2=\{(a,b)\in\GL_2^2:\det a=\det b\}
\]
has dimension seven. The stabilizer of a pair of lines is a pair of Borels with equal determinants, of dimension $3+3-1=5$. Fixing the determinant leaves dimension four. A semilinear change of embedding need not lie in either linear fibre. For example choose algebraically independent complex numbers $t_1,t_2$ and embeddings of $\Qb(t)$ sending $t$ to them. The Legendre curves have
\[
 j(t)=256\frac{(1-t+t^2)^3}{t^2(1-t)^2}.
\]
Their $j$-values are algebraically independent: each $t_i$ is algebraic over its nonconstant rational value $j(t_i)$. An isogeny would impose a modular-polynomial relation, so the two curves are not isogenous. Any linear commensurability of their Betti lattices preserving the Hodge line would be a rational Hodge isomorphism and hence an isogeny by the degree-one Riemann theorem. This is impossible. One may append the same non-CM curve $y^2=x^3-x+1$ on both sides to obtain two-factor examples. Section~4 computed its nonintegral $j=-6912/23$; this algebraic value cannot be isogenous to either transcendental $j(t_i)$. The field embeddings still give semilinear comparison maps.

\subsection{Adelic genericity and its elliptic applications}
\begin{proposition}\label{prop:adelic-generic}
Let $G/\Q$ be connected affine of finite type. Every coset of a compact open subgroup of $G(\mathbf A_f)$ has a set of full Haar measure whose points are generic over $\Q$ in every $\ell$-component. The same holds for a fixed finite tuple, with joint genericity in the corresponding power of $G$.
\end{proposition}
\begin{proof}
A nonzero regular function on $G_{\Q_\ell}$ has a zero set of local Haar measure zero. To see this in analytic coordinates, its germ is a nonzero convergent power series: a regular local ring injects into its completion, and a regular function on the integral variety cannot vanish on an analytic open set. In one variable a nonzero convergent series has isolated zeros on sufficiently small discs. In higher dimension expand in the last variable. The locus where all coefficient series vanish is contained in the zero set of one nonzero coefficient and has measure zero by induction; elsewhere apply the one-variable assertion and Fubini. A countable analytic cover proves the claim.

The projection of a compact open adelic subgroup $U$ is compact open locally, and normalized Haar measure pushes forward to its normalized Haar measure. For a coset $Ua$, the set defined by $f((ua)_\ell)=0$ therefore has measure zero for each prime and each nonzero $f\in\Q[G]$. There are countably many such pairs. Removing their union gives the assertion. Apply precisely this argument to $G^r$ and $U^r$ for joint genericity; separate genericity of the $r$ entries would not suffice. Local trivialization gives the identical proof for a torsor.
\end{proof}

Suppose the continuous adelic Galois image $\Gamma$ of a chosen integral realization is open. It is compact, hence closed. Changing $n-1$ places independently multiplies their relative frames by independent elements of $\Gamma$. Apply Proposition~\ref{prop:adelic-generic} to the product of these cosets and lift the chosen elements to Galois automorphisms. Such automorphisms extend to an ambient algebraically closed field by choosing a transcendence basis. Hence suitably changed places attain the bound of Proposition~\ref{prop:relative-frames} simultaneously at every prime.

For one non-CM elliptic curve over a number field, the required openness is Serre's theorem \cite[\S4.4, Th\'eor\`eme~3 and $3'$, p.~299]{SerreAdelic}; the resulting degree is $4(n-1)$. For geometrically pairwise non-isogenous non-CM elliptic curves $E_1,\ldots,E_r$, the rational group and its integral form are
\[
 \Delta_r=\{(g_i)\in\GL_2^r:\det g_i=\det g_j\},
 \qquad \dim\Delta_r=3r+1.
\]
The group description is the known Hodge-group calculation for elliptic products \cite[Proposition~5.14, pp.~20--21]{KSV}: the Hodge group is $\SL_2^r$, and adjoining the common weight gives $\Delta_r$. Its Hodge tensors are generated by divisors; the identities and polarizations here are defined over the ground number field, so they give the same motivated group. Lombardo's adelic theorem gives finite index in $\Delta_r(\widehat\Z)$, including the prime two \cite[Theorem~1.1 and Remark~1.2, pp.~1--2]{LombardoProducts}. A compact finite-index subgroup is open: its finitely many cosets are closed, so each coset has closed complement. The attained local degrees are therefore
\[
                         (n-1)(3r+1).
\]
These are applications of the cited arithmetic theorems, not new open-image theorems. Serre's original two-curve statement has the stronger hypothesis that the two representation systems never become isomorphic over a finite extension \cite[\S6.1, Th\'eor\`eme~6, p.~324]{SerreAdelic}. Faltings's isogeny theorem supplies that implication from geometric non-isogeny \cite[\S5, Satz~4 and Korollar~2, pp.~360--361]{FaltingsIsogeny}. The conclusion concerns suitably changed places, rather than every tuple.

\subsection{Approximation, commensurability and openness}
Write $\overline{G(\Q)}=G(\mathbf A_f)$ for rational approximation in the finite adeles. For a connected affine group in characteristic zero, the known strong-approximation criterion says that this holds exactly when its reductive quotient is semisimple and simply connected and every almost-simple rational factor has noncompact real points. The unipotent radical contributes no obstruction. The reduction and necessary conditions are Kneser's Satz~1 and Hilfss\"atze~2.3--2.5 \cite[pp.~191--194]{KneserApprox}; the sufficiency for simply connected simple groups is Platonov's Theorem~A \cite[p.~1140]{PlatonovApprox}, with its addendum \cite[\S\S2 and 5]{PlatonovAddendum}. We use that exact criterion as a known theorem.

The cases relevant here can also be seen directly. The additive group has approximation by the Chinese remainder theorem. For $\SL_2$, the closure of $\SL_2(\Q)$ contains both elementary unipotent groups with parameters in $\mathbf A_f$, since $\Q$ is dense in $\mathbf A_f$. They generate $\SL_2(\mathbf A_f)$. Indeed for a matrix with first column $(a,c)$ choose an adelic $t$ so that $a+tc$ is invertible at every prime and a unit almost everywhere; at integral primes the unimodular column permits such an integral $t$. Elementary elimination then reduces to a diagonal matrix, and $\operatorname{diag}(u,u^{-1})$ is a product of elementary matrices for every adelic unit $u$. Thus $\SL_2$ has approximation.

In contrast, $\Q^\times\cap\widehat\Z^\times=\{1,-1\}$, so the nonempty compact open coset of units congruent to $2$ modulo $5$ contains no rational point. The determinant map, with its section $a\mapsto\operatorname{diag}(a,1)$, transfers this failure to $\GL_2$ and to $\Delta_r$ by using the same diagonal matrix in every factor. Openness of the elliptic adelic image therefore does not require approximation in its full motivic group.

\begin{proposition}\label{prop:all-frames}
Let $H\subset G(\mathbf A_f)$ be a compact subgroup. If $G(\Q)$ is dense, then
\[
 G(\mathbf A_f)=G(\Q)H\quad\Longleftrightarrow\quad H\text{ is open}.
\]
In particular the condition concerns all adelic frames. It cannot be inferred from a statement about an unspecified smaller collection of frames arising from places.
\end{proposition}
\begin{proof}
If $H$ is open, density makes every coset $xH$ meet $G(\Q)$. Conversely $G(\Q)$ is countable, so the displayed equality covers the locally compact Baire space by countably many closed translates of $H$. One has nonempty interior, and therefore $H$ has nonempty interior. Translating an interior point to the identity shows that the subgroup $H$ is open.
\end{proof}

Local genericity has no such converse. In $\SL_2(\mathbf A_f)$ let $I_\ell$ be the inverse image of the upper triangular subgroup of $\SL_2(\mathbf F_\ell)$ and put $H=\prod_\ell I_\ell$. Each $I_\ell$ is open, of index $\ell+1$, but $H$ is not open in the restricted product: an identity neighbourhood contains the full integral group at almost every prime. The measure proof of Proposition~\ref{prop:adelic-generic} still applies to product Haar measure on $H$ and on its finite powers, giving simultaneous local genericity. This is an algebraic-group example, not a claimed motivic image.

A genuine motivic example is the square of the cyclotomic character, the finite image of $\Q(-2)$ over $\Q$:
\[
                  H=(\widehat\Z^\times)^2.
\]
It is closed as a compact image and has open projections at every prime. At each odd prime the squares have index two, so $H$ has infinite index and is not adelically open. The quotient maps modulo any finite set of odd primes already give arbitrarily large powers of two as index bounds. Nevertheless product Haar measure again supplies local genericity in every component. This separates genericity, rational approximation and global openness.

\subsection{Finite generation and uncountable fields}

The coherent path in Section~10.1 is a special case of a general finite-type
comparison. Finite generation gives enough common open sets to choose one
path for the whole diagram.

\begin{proposition}\label{prop:finite-generated-comparison}
Let $k\subset\C$ have characteristic zero and let $k_0$ be its algebraic part
over $\Q$. If $v,w:k\hookrightarrow\C$ agree on $k_0$, every finite
smooth-projective algebraic-correspondence diagram has an integral Betti
comparison, compatible with its cup products, K\"unneth maps and duality.
If $k/k_0$ is finitely generated, one comparison works simultaneously on the
entire diagram. It extends to the induced rational tensor constructions and
to the real spans of the integral realizations.
\end{proposition}

\begin{proof}
First suppose that $k/k_0$ is finitely generated. The constants $k_0$ are
relatively algebraically closed in $k$; in characteristic zero this makes
the extension regular, hence geometrically integral
\cite[Tag 037P and Section~10.45]{StacksE}. Choose a smooth geometrically integral model $S/k_0$
with function field $k$. Both embeddings define $k_0$-generic complex points
of $S$, hence points in every nonempty $k_0$-open subset.

The field $k_0$ is countable, so the proper closed $k_0$-subvarieties $Z_j$
of $S$ form a countable family. We show that the complement of their complex
points is path connected. In a complex coordinate ball of dimension $d$,
a proper analytic set has real Hausdorff dimension at most $2d-2$. Cover
its smooth strata by countably many bounded coordinate pieces. For a fixed
point $x$ outside the set, endpoints whose segment from $x$ meets one of
those pieces lie in cones parametrized by
\[
 (z,t)\longmapsto x+t(z-x),\qquad t\geq1.
\]
On bounded parameter charts these maps are locally Lipschitz and their
domains have dimension at most $2d-1$. Their images have $2d$-dimensional
measure zero. A countable union of such bad endpoint sets is still null.
Given two points outside all $Z_j$ in the ball, choose an intermediate point
outside the bad cones from either endpoint. The two joining segments avoid
every $Z_j$. A finite chain of overlapping coordinate balls along an initial
path in the connected manifold $S(\C)$, with intermediate overlap points
outside the same countable union, gives a path $\gamma$ joining the original
points and avoiding every $Z_j$. In dimension zero the embeddings agree.

Each smooth projective variety over $k$ spreads to a smooth projective family
over a nonempty open of $S$: spread its finite homogeneous equations and
then descend smoothness \cite[Tags 01ZM and 0C0C]{StacksE}.
Each finite list of correspondences and tensor operations spreads after
further shrinking. Relative cycle classes give
maps of integral cohomology local systems: one can resolve the finitely many
cycle supports by the characteristic-zero resolution theorem and shrink so that their resolutions
are smooth and proper over the base, then use their proper pushforwards.
The generic cycle class is unchanged by this resolution. Thus this assertion
also follows from functoriality of smooth proper transport.

Every excluded set is among the $Z_j$, so the single path $\gamma$ lies in
the good open for each object and each finite list of arrows. Ehresmann's
theorem gives the required local systems
\cite[Theorem~12.1.1, printed p.~97, PDF p.~98]{Arapura}; path transport
is their functorial transport
\cite[Theorem~12.3.1, printed p.~101, PDF p.~102]{Arapura}.
It preserves integral cohomology, cup products, product maps and the complex
orientation. Hence it also preserves the duality and correspondence maps.
Two choices of spreading agree on a common smaller open, which still
contains $\gamma$, so the resulting maps are independent of these choices.
They define the simultaneous comparison. Passing to free quotients, rational
projector images and real scalar extensions preserves it.

For a finite diagram over arbitrary $k$, choose a finitely generated
subfield $l/\Q$ containing its finitely many defining coefficients. Its
algebraic part is contained in $k_0$, so the restricted embeddings agree
there. Apply the result just proved to $l$ and the descended diagram.
This proves the finite assertion without selecting an arbitrary inverse
system of comparisons.
\end{proof}

\begin{corollary}\label{cor:uncountable-comparisons}
There are distinct complex places of an uncountable field inducing
isomorphic entire integral Betti correspondence diagrams. There are also
two places of an uncountable field, agreeing on $\Qb$, for which no such
comparison exists. Both examples can be realized on the field $\C$ itself.
\end{corollary}

\begin{proof}
Choose a transcendence basis $B$ of $\R/\Q$, an element $b\in B$, and the
real closure $L$ of $\Q(B\setminus\{b\})$ inside $\R$. The field $L$ is
uncountable and $b$ is transcendental over it. Its algebraic closure is
$L(i)$. Let $k$ be an algebraic closure of $L(t)$ and choose an embedding
$v:k\to\C$ sending $t$ to $b$ and agreeing with the fixed inclusion on
$L(i)$. Since $\C$ is algebraic over $L(b)$, this identifies $k$ with $\C$.

The polygon
\[
 b\longrightarrow b+ib\longrightarrow b+1+ib\longrightarrow b+1
\]
avoids $L(i)=L+iL$: the real coordinate on either vertical edge is outside
$L$, and the imaginary coordinate on the horizontal edge is $b\notin L$.
Every algebraic function over $L(t)$ has its finite branch and pole set in
$L(i)$. Continue the branches selected by $v$ along this polygon. Their
finite endpoint values preserve all algebraic operations, and a nonzero
algebraic function cannot vanish above the transcendental endpoint $b+1$.
This gives an embedding $w:k\to\C$ fixing $L(i)$ with $w(t)=b+1$.
It is distinct from $v$ and not conjugate to it, since both fix $i$.

A finite collection of geometric data descends to a finite extension of
$L(t)$ and spreads over a curve finite above an open of the $t$-line. The
branch points and images of bad loci lie in $L(i)$, so the same polygon lifts
to the good locus, from its $v$-sheet to its $w$-sheet. Smooth proper
transport gives the comparisons. Enlarging the field or the collection
preserves them by unique path lifting and functoriality. They therefore
define one integral comparison on the entire uncountable diagram.
Here the common excluded set, rather than its cardinality, makes the
simultaneous construction possible.

For the negative example use the pair on
$\overline{\Qb(t)}$ in Section~10.2. Their difference is an automorphism
over $\Qb(t)$. Extend it to $\C$ by choosing a transcendence basis and then
extending through an algebraic closure. A full integral comparison for the
two resulting embeddings of $\C$ would restrict to the forbidden comparison
on the original abelian and isogeny diagram. Both embeddings still fix
$\Qb$.
\end{proof}

The coefficient field matters separately from this integral obstruction.
Wibmer proves that every torsor under an affine group scheme over an
algebraically closed field is trivial, without a finite-type or cardinality
restriction \cite[Theorems~1.1--1.2, p.~447]{Wibmer}. Thus the tensor-Isom
torsor of two rational fibre functors has a point after extension to $\C$,
using its standard affine torsor representation
\cite[Theorem~3.2, p.~31]{DeligneMilne}. This known theorem concerns the complexified
torsor; it does not supply an integral point or remove the obstruction in
Section~10.2.

\section{Number fields: finite tests and a simultaneous CM comparison}
\label{sec:number-fields}
For number fields, finite labelled diagrams can distinguish many places even
after Hodge filtrations have been forgotten. The class of objects in the test
matters: we shall also construct two real places whose entire CM-abelian and
Artin diagrams admit a simultaneous integral comparison. Throughout this
section a complex place is an embedding into $\C$ modulo complex conjugation.
The integral correspondence diagram has exactly the convention of
Section~\ref{sec:integral}; in particular its objects include smooth projective
schemes that need not be geometrically connected.

\subsection{A labelled elliptic curve and descent of finite tests}
\begin{proposition}\label{prop:ring-class-test}
Let $\mathcal O$ be an order in an imaginary quadratic field $K$, and let
$j$ be the invariant of an elliptic curve with endomorphism ring $\mathcal O$.
Put $H=K(j)$, and choose an elliptic curve $E/H$ with invariant $j$.
The integral Betti diagram of $E$ and one labelled nonrational endomorphism
distinguishes the complex places of $H$.
\end{proposition}
\begin{proof}
Every geometric endomorphism of $E$ is defined over $H$. Indeed the differential
at the origin injects $\End(E_{\Qb})$ into $\Qb$: after an embedding into $\C$,
a homomorphism of elliptic tori is determined by its complex linear derivative.
The image is the specified order in $K$. If a Galois element fixes $H$, it
fixes that derivative; it consequently fixes the endomorphism itself.

Choose $a\in\mathcal O\setminus\Q$. At an embedding $v:H\hookrightarrow\C$,
let $A_v$ be its labelled operator on $H_1(E_v,\R)$. The analytic complex
structure is recovered from the operator by
\[
 J_v=\frac{A_v-\operatorname{Re}(v(a))\,1}
                 {\operatorname{Im}(v(a))}.
\]
Here the identification of $\mathcal O$ with its differential action has
been fixed once and for all. Suppose two embeddings $v,w$ give an integral
comparison commuting with $a$. Replacing $w$ by its complex conjugate if
necessary, arrange that $v|_K=w|_K$. Complex conjugation itself supplies
the corresponding integral comparison of the labelled tori. The displayed
formula now says that the original lattice map is complex linear. It gives
an isomorphism of elliptic tori, so $v(j)=w(j)$. Since $H=K(j)$, the embeddings
are equal. This proves the assertion modulo conjugation.
\end{proof}

The arithmetic description of $H$ is the classical ring class theorem:
it is the ring class field of $\mathcal O$
\cite[Theorem~11.1, p.~220]{CoxClassFields}. This named input identifies the
fields covered by the proposition; the reconstruction of $J_v$ proves the
separation assertion directly. For a concrete example take
\[
 K=\Q(\sqrt{-5}),\qquad H=\Q(i,\sqrt5),\qquad
 j_\pm=632000\pm282880\sqrt5.
\]
These are the CM invariants of discriminant $-20$; the positive-sign value
and this Hilbert class field are explicitly recorded in
\cite[\S5, p.~11, before (19)]{CampagnaPengo}. The equation
\[
 y^2=x^3-\frac{3j_+}{j_+-1728}x
              -\frac{2j_+}{j_+-1728}
\]
is nonsingular and has invariant $j_+$. The two embeddings of $H$ fixing
$K$ send $j_+$ to $j_+$ and $j_-$. The labelled action of $\sqrt{-5}$
therefore already separates them. This is the explicit field version of
the ideal-class obstruction in Section~\ref{sec:integral}.

\begin{proposition}\label{prop:finite-test-descent}
Let $L/k$ be a finite extension of number fields. A finite integral Betti
diagram over $L$ separating its complex places gives a finite such diagram
over $k$ separating its complex places. More generally it is enough that
comparisons between embeddings of $L$ force their restrictions to $k$ to
agree up to conjugation. The same statements hold over $\R$.
The resulting diagram is contained in the tensor constructions of one
finite labelled direct sum.
\end{proposition}
\begin{proof}
For $X/L$ write $X_{[k]}$ for the same scheme regarded as a $k$-scheme by
composition of its structural map. Finite separability preserves smoothness
and projectivity. With $A_L=h^0(\operatorname{Spec}L)$, one has
\[
 H^*(X_{[k],v},\Z)/\mathrm{tors}
   =\bigoplus_{u|_k=v}H^*(X_u,\Z)/\mathrm{tors}.
\]
The structure map and cup product make this an $A_L$-module. The primitive
idempotent $e_u$ of $H^0(\operatorname{Spec}L_v,\Z)$ projects to the
$u$-summand. Include $A_L$, its unit and multiplication, the $X_{[k]}$,
their module structures, and the original labelled cycles in the descended
diagram. A cycle on $X\times_LY$ is a cycle of the same codimension on the
open and closed diagonal component of $X_{[k]}\times_kY_{[k]}$.
The diagonal class of $\operatorname{Spec}L$ gives
\[
                         \sum_u e_u\otimes e_u,
\]
which selects these relative products.

A comparison at $v,w$ is a unital algebra isomorphism on $A_L$, hence
permutes its primitive idempotents. If $e_u$ goes to $e_{u'}$, naturality
for all module structures restricts the comparison to one comparison of
the original diagram at $u,u'$. The diagonal class ensures tensor
compatibility on these blocks. Their restrictions are $v,w$, so the
assumed separating property gives $w=v$ or $\overline v$.
Over $\R$ the identical argument uses the primitive idempotents of a
product of copies of $\R$.

Only finitely many objects and maps were added. Their direct sum, with
the specified summand inclusions and projections, and the finitely many
Tate lines involved, is one generator for this labelled diagram. No
integral K\"unneth problem is hidden here: the integral K\"unneth sequence
for finite CW complexes splits as abelian groups, its Tor terms are
torsion, and its canonical map becomes an isomorphism after quotienting
both sides by torsion. Grading and the usual signs are retained.
\end{proof}

\subsection{The exact transporter and the remaining universal question}
Let $\omega_{\Q,\Z}$ denote the usual integral Betti diagram over $\Q$.
For a number field $k$ write
\[
 A_k=h^0(\operatorname{Spec}k),\qquad
 \omega_{\Q,\Z}(A_k)=\Z^{\operatorname{Hom}(k,\C)}.
\]
Its primitive idempotents are $e_v$, indexed by embeddings $v$.

\begin{proposition}\label{prop:exact-transporter}
For embeddings $v,w:k\hookrightarrow\C$, restriction to objects obtained
by base change from $\Q$ gives a bijection
\[
 \Isom^\otimes(\omega_{v,\Z},\omega_{w,\Z})
 \simeq
 \{g\in\Aut^\otimes(\omega_{\Q,\Z}):g_{A_k}(e_v)=e_w\}.
\]
The action of $g$ on all Artin objects is induced by a unique
$\sigma\in\Gal(\Qb/\Q)$; for elements in this transporter it satisfies
$\sigma v=w$.
\end{proposition}
\begin{proof}
If $\eta$ is a comparison over $k$, its component on $M_k$ defines an
automorphism $g_M$ of the Betti realization of every $M/\Q$.
The diagonal section of $\operatorname{Spec}k\times_\Q\operatorname{Spec}k$
over $k$ gives evaluation at $v$ and $w$, respectively. Naturality gives
$\operatorname{ev}_w g_{A_k}=\operatorname{ev}_v$. Since $g_{A_k}$ permutes
primitive idempotents, this says $g_{A_k}(e_v)=e_w$.

Conversely, for such a $g$ apply it to $X_{[\Q]}$ for each $X/k$.
Commutation with the $A_k$-module action sends the $v$-block to the
$w$-block. Restrict $g$ to those blocks. The diagonal component and its
cycle, as in Proposition~\ref{prop:finite-test-descent}, recover every
$k$-correspondence and tensor product; units, duality and Tate maps
restrict as well. Thus these restrictions form $\eta$.
These constructions are inverse: $X$ is the diagonal open and closed
summand of $X_{[\Q]}\times_\Q k$, whose inclusion and projection recover
the initial $\eta_X$. In the other direction the underlying $\Q$-object
of $M_k$ is $M\otimes A_k$. Restricting $g_M\otimes g_{A_k}$ from its
$e_v$-block to its $e_w$-block recovers exactly $g_M$.

For a finite Galois extension $L/\Q$, naturality with every deck
transformation makes the permutation of $\operatorname{Hom}(L,\C)$
commute with its regular right Galois action. It is consequently a left
translation. These translations agree on subextensions, and their
inverse limit is the unique $\sigma$. Its action on $e_v$ is the stated one.
\end{proof}

There is an exact unresolved formulation of universal separation:
\[
 \operatorname{im}\bigl(\Aut^\otimes(\omega_{\Q,\Z})
                  \longrightarrow\Gal(\Qb/\Q)\bigr)=\{1,c\}.
\]
This equality is equivalent to separation of the complex places of every
number field by its entire integral diagram. The identity and ordinary
complex conjugation occur. Any other Galois element differs from both
on some finite Galois subextension and gives a nontrivial transporter
there; conversely a comparison between distinct complex places gives
such an element by Proposition~\ref{prop:exact-transporter}. We prove
neither the displayed equality nor its negation.

\subsection{The generalized-dihedral construction and its limitation}
\begin{proposition}\label{prop:dihedral-test}
Let $L/\Q$ be finite Galois and contain an imaginary quadratic field $K$.
Assume $B=\Gal(L/K)$ is abelian and actual complex conjugation satisfies
$cbc^{-1}=b^{-1}$ for all $b\in B$. Then every subfield of $L$ has a finite
integral diagram separating its complex places. This includes every cubic
field of negative discriminant. In contrast, a totally real Galois
subfield of an elliptic ring class field has elementary abelian
$2$-group as its Galois group.
\end{proposition}
\begin{proof}
The group is $B\rtimes\langle c\rangle$, with $c^2=1$. Its transfer to $B$
is trivial: on $b\in B$ it is $bcbc^{-1}=1$, and on $c$ it is $c^2=1$.
Artin reciprocity identifies transfer with extension of ideals from
$\Q$ to $K$ \cite[Chapter~V, \S3, Proposition~3.18 and the accompanying
diagram, p.~162]{MilneCFT}. Choose an integer $f>0$ such that $f\mathcal O_K$
is a defining modulus for $L/K$, enlarging it to contain all ramified
primes. The Artin kernel contains ray principal ideals and all extended
rational ideals prime to $f$.

Suppose $\alpha\equiv a\pmod{f\mathcal O_K}$ with $a\in\Z$ prime to $f$.
Locally at $f$, the ratio $\alpha/a$ is congruent to $1$. Thus $(\alpha)$
is the product of the ray principal ideal $(\alpha/a)$ and the extended
rational ideal $(a)$, and both have trivial Artin image. By the order
class-group description
\cite[Proposition~7.22, pp.~145--146]{CoxClassFields}, the Artin map
factors through the class group of $\Z+f\mathcal O_K$. Hence $L$ is
contained in its ring class field. This is the generalized-dihedral
criterion of \cite[Theorem~9.18, pp.~191--192]{CoxClassFields}; the transfer
calculation above verifies the direction used here. Theorem~11.1 of
that reference and Propositions~\ref{prop:ring-class-test}--\ref{prop:finite-test-descent}
now give the test.

For a cubic field of negative discriminant its normal closure has Galois
group $S_3$. The quadratic resolvent is imaginary, the subgroup $A_3$ is
cyclic, and complex conjugation is a transposition, which acts on $A_3$
by inversion. This proves the cubic assertion.

A ring class field itself has this generalized-dihedral group
\cite[Lemma~9.3, pp.~180--181]{CoxClassFields}. In a totally real Galois
quotient, the image of $c$ is trivial. The inversion relation then says
that every image of $b$ has square one, so the quotient is elementary
abelian of exponent two. The normal closure of any totally real subfield
inside a ring class field is still totally real and is still contained
in that field. Consequently a totally real field with nonabelian
dihedral normal closure cannot arise from this ring class construction.
This restriction concerns that construction alone.
\end{proof}

\subsection{A polarized CM test for all totally real fields}
\begin{theorem}\label{thm:real-field-CM-test}
Every totally real number field has a finite labelled algebraic
correspondence diagram separating its embeddings, even with real
coefficients. Every CM field has a finite such test for its complex places.
In particular every cubic number field has a finite integral separating test.
\end{theorem}
\begin{proof}
We first explain what a labelled CM action and polarization recover.
Let $A$ have dimension $n$ and an action by a CM field $E$ of degree $2n$.
Write $F=E^+$ and $V=H_1(A,\R)$. The space $V$ has rank one over
$E\otimes\R=\prod_{\tau:F\hookrightarrow\R}\C$.
Let $Q$ be a polarization with Rosati involution complex conjugation
on $E$. The idempotents of $F\otimes\R$ are self-adjoint, so the real
two-dimensional factors $V_\tau$ are mutually orthogonal and each
restriction of $Q$ is nondegenerate. On each factor choose multiplication
$I_\tau$ by $i$ in one of its two complex identifications. The analytic
complex structure commutes with $E$ and has square $-1$, so
\[
 J|_{V_\tau}=s_\tau I_\tau,\qquad s_\tau\in\{1,-1\}.
\]
The homological positivity convention is $Q(x,Jx)>0$. Therefore the
sign of $Q(x,I_\tau x)$ is $s_\tau$ for every nonzero $x\in V_\tau$.

A real comparison commuting pointwise with the labelled action of $E$
preserves these factors and their $I_\tau$. The one labelled Tate-valued
polarization gives a similitude
\[
 Q'(Tx,Ty)=\lambda Q(x,y),\qquad\lambda\in\R^\times,
\]
with the same scalar on every factor. Hence
$s'_\tau=\operatorname{sign}(\lambda)s_\tau$ for all $\tau$.
The two CM types must be equal or everywhere conjugate. Integrally the
Tate scalar is $\pm1$. In cohomological notation the polarization is its
algebraic Chern class $\Z(-1)\to\bigwedge^2H^1(A,\Z)$, and the homology
comparison is the inverse dual of the cohomology comparison. Dualizing
only inverts the common Tate scalar. Thus this argument uses algebraic
labels in the stated diagram and does not assume that the Hodge
filtration is an additional label.

Now let $F\subset\R$ be totally real of degree $n\ge3$ and put $E=F(i)$.
Choose the type $\Phi$ that sends $i$ to $i$ at the given embedding of
$F$ and to $-i$ at every other embedding. Weak approximation gives
$b\in\mathcal O_F$ positive at the distinguished embedding and negative
at all the others. On $R=\mathcal O_F[i]$ put
\[
 A_\Phi(\C)=\C^\Phi/\Phi(R),\qquad
 Q(x,y)=\operatorname{Tr}_{E/\Q}(ibx\overline y).
\]
This form is integral, alternating and nondegenerate. If
$\phi(ib)=i d_\phi$ for $\phi\in\Phi$, then every $d_\phi>0$ and
\[
 Q(z,w)=2\sum_{\phi\in\Phi}d_\phi
              \operatorname{Im}(\overline {z_\phi}w_\phi),\qquad
 Q(z,Jz)=2\sum_{\phi\in\Phi}d_\phi|z_\phi|^2>0.
\]
The Riemann bilinear criterion makes this a polarized abelian variety,
with the labelled $R$-action and its required Rosati involution
\cite[Example~2.11, p.~26, and Proposition~3.17, p.~31]{MilneCMNotes}.
CM algebraicity and full faithfulness under algebraically closed base
extension give a model over $\Qb$, together with these endomorphisms and
polarization \cite[Proposition~7.9 and Corollary~7.10, pp.~54--55]{MilneCMNotes}.
Choose a finite field $M$ containing their coefficients and the normal
closure of $E$. A $\Z$-basis of $R$ provides finitely many actual
endomorphism labels. The polarization is represented by its algebraic
map to the dual and the Poincar\'e pairing.

At an embedding $u:M\hookrightarrow\C$ the tangent representation
has type $u\Phi$: applying $u$ to the endomorphism matrices applies it
to every eigencharacter. The type traces satisfy, for $a\in F$,
\[
 t_\Phi(i)=(2-n)i,\qquad
 t_\Phi(ia)=i\bigl(2a-\operatorname{Tr}_{F/\Q}(a)\bigr).
\]
They generate $F(i)$ because $n\ne2$. Equal or conjugate types at
$u,u'$ have equal or conjugate type traces. Thus $u'|_F=u|_F$.
Proposition~\ref{prop:finite-test-descent}, in its restriction-of-place
form and with real coefficients, descends this finite diagram to $F$.

For $F=\Q$ there is only one embedding. If $F$ is real quadratic,
choose another real quadratic field $F_0\ne F$. The totally real
biquadratic field $FF_0$ has degree four, so its test just constructed
descends to $F$. Finally restriction from a CM field to its maximal
totally real subfield bijects complex places with real embeddings.
Base change of that subfield's test therefore gives the CM assertion.
The cubic assertion combines this theorem with
Proposition~\ref{prop:dihedral-test}.
\end{proof}

Here are explicit integral data for the totally real cubic
$F=\Q(\theta)$, where $\theta$ is the largest root of $x^3-4x+1$.
Its three roots lie in $(-3,-2),(0,1),(1,2)$, respectively, and its
discriminant is $229$. In the preceding construction use
$R=\Z[\theta,i]$ and $b=\theta-1$. In the ordered basis
$1,\theta,\theta^2,i,i\theta,i\theta^2$, the matrices are
\[
 Q=\begin{pmatrix}0&2T\\-2T&0\end{pmatrix},\qquad
 T=\begin{pmatrix}-3&8&-11\\8&-11&35\\-11&35&-52\end{pmatrix},
\]
\[
 [\theta]=\operatorname{diag}(C,C),\quad
 C=\begin{pmatrix}0&0&-1\\1&0&4\\0&1&0\end{pmatrix},\qquad
 [i]=\begin{pmatrix}0&-I_3\\I_3&0\end{pmatrix}.
\]
To derive $T$, reduce powers using $\theta^3=4\theta-1$ and use
$\operatorname{Tr}(1)=3$, $\operatorname{Tr}(\theta)=0$,
$\operatorname{Tr}(\theta^2)=8$. Then
$T_{ab}=\operatorname{Tr}((\theta-1)\theta^{a+b})$ for $0\le a,b\le2$.
These calculations give $\det T=458$ and $\det Q=3664^2$,
$[\theta]^tQ=Q[\theta]$ and $[i]^tQ=-Q[i]$.
The signs of $b$ are $(-1,-1,1)$; the CM type chooses the corresponding
signs for $i$. Its traces are $t_\Phi(i)=-i$ and
$t_\Phi(i\theta)=2i\theta$. These specify the polarized torus and the
finite action labels. The algebraic model is supplied by the CM
algebraicity theorem above; projective equations for it are not needed.

\subsection{Two quartic places with the same entire CM diagram}
\begin{theorem}\label{thm:quartic-CM-comparison}
Let $k=\Q(\alpha)$, where $\alpha$ is a real root of $x^4-x-1$.
The two real embeddings of $k$ are distinct complex places but have
isomorphic integral Betti functors on the entire following diagram:
CM abelian varieties over $k$, finite etale $k$-schemes, their
cohomology, products, Tate twists, tensor constructions and duals,
algebraic correspondences, and compatible direct summands with their
assigned intersection lattices. The comparison on CM abelian generators
comes from actual algebraic isomorphisms. It preserves labelled
polarizations and ordinary Hodge structures. It does not intertwine
the geometric real-Frobenius involutions.
\end{theorem}
\begin{proof}
In the fixed $\Qb\subset\C$ let $\Q^{\mathrm{cm}}$ be the compositum
of all CM fields and let $F_{\mathrm{cm}}$ be the compositum of their
maximal abelian extensions. Both are normal over $\Q$.
The extension $F_{\mathrm{cm}}/\Q^{\mathrm{cm}}$ is abelian, since its
Galois group injects into a product of abelian Galois groups. Also
$\Q^{\mathrm{ab}}\subset F_{\mathrm{cm}}$: for any imaginary quadratic
$K$, the extension $K\Q^{\mathrm{ab}}/K$ is abelian.

Fix $\sigma\in\Gal(\Qb/F_{\mathrm{cm}})$ and a CM abelian variety
$A/\Qb$. Choose a CM algebra $E$ acting on $A$ and its type $\Phi$.
Its reflex field $E^*$ is a CM field even when $E$ is a product
\cite[Proposition~1.18, p.~14]{MilneCMNotes}. The fundamental theorem
of complex multiplication states that, for any finite idele $s$ of
$E^*$ with Artin image $\sigma|_{E^{*\mathrm{ab}}}$, there is a unique
$E$-linear quasi-isogeny $a_A:A\to\sigma A$ such that
\[
                   a_A(N_\Phi(s)x)=\sigma x\quad(x\in V_f A).
\]
This is precisely \cite[Theorem~9.10, p.~78]{MilneCMNotes}; the quoted
``isogeny'' there is defined as an invertible rational homomorphism
on p.~23. Since $\sigma$ fixes $E^{*\mathrm{ab}}$, take $s=1$.
Thus $a_A$ acts as $\sigma$ on every rational Tate module.

We check integrality using this one global quasi-isogeny. Write
$a_A=f/m$ with $m>0$ and $f:A\to\sigma A$ an actual homomorphism.
Galois transport preserves every integral Tate lattice, so
\[
                  f(T_\ell A)\subset mT_\ell(\sigma A)
                         \quad\hbox{for every prime }\ell.
\]
If $\ell^e\Vert m$, reduction modulo $\ell^e$ shows that $f$ kills
$A[\ell^e]$. Hence it kills $A[m]$. The quotient property of
$[m]:A\to A$ gives $f=a_0\circ[m]$ with an actual homomorphism $a_0$
\cite[Chapter~I, Remark~8.12, p.~39]{MilneAV}. Equality in
$\operatorname{Hom}(A,\sigma A)\otimes\Q$ gives $a_0=a_A$. The inverse quasi-isogeny
has inverse integral Tate action, so the same argument makes its inverse
an actual morphism. Therefore $a_A$ is an algebraic isomorphism.

A homomorphism of abelian varieties is determined by one rational Tate
realization: over $\C$, zero Tate action gives zero integral homology
action and hence zero linear map of tori. This faithfulness proves that
$a_A$ is independent of the chosen CM algebra and that for every
homomorphism $f:A\to B$,
\[
                    (\sigma f)a_A=a_Bf.
\]
It also proves compatibility with products and duals. Since $\sigma$
fixes $\Q^{\mathrm{ab}}$, its cyclotomic character is one. The Weil
pairing makes the Tate comparison the identity and every labelled
polarization compatible. Equivalently the multiplier in
\cite[Remark~9.11(c), p.~78]{MilneCMNotes} is
$\chi_{\mathrm{cyc}}(\sigma)/N_{E^*/\Q}(s)=1$.
These algebraic isomorphisms preserve ordinary Hodge structures.

For completeness, compatibility with arbitrary correspondences is more
than compatibility with homomorphisms. The cohomology comparison is
$(a_A^{-1})^*$. Integral cohomology of an abelian variety is the exterior
algebra of $H^1$, so these maps give its exterior and product
comparisons. Finite-coefficient Betti--etale comparison
\cite[Expos\'e~XVI, Theorem~4.1]{SGAComparison} identifies their
$\Z_\ell$-completions with Galois transport. Algebraic cycle classes
and their correspondence operators are Galois compatible. Naturality
for any algebraic correspondence therefore holds over $\Q_\ell$,
hence over $\Q$ by injectivity, and then on the free integral groups
because the maps are integral. This treats all labelled cycles,
without assuming they are generated by endomorphisms.

On an Artin object, send the basis indexed by a geometric point $x$
to the basis indexed by $\sigma x$. On products use these permutations
of components and the isomorphisms $a_A$. The same etale argument gives
naturality for all correspondences involving Artin factors. Tensor
products and duals use the corresponding tensor maps and dual inverses.
For a compatible projector $p$, the relation $Tp=(\sigma p)T$ restricts
to its image and carries its assigned intersection lattice to the
conjugate lattice. We have constructed one simultaneous comparison on
the entire stated diagram.

We now choose $\sigma$ to exchange the two real roots. Let $L$ be the
splitting field of $x^4-x-1$. Its discriminant is $-283$. Modulo $2$
the polynomial has no root and is not divisible by the sole irreducible
quadratic $x^2+x+1$, so it is irreducible. Modulo $7$ it factors as
\[
       (x-3)(x^3+3x^2+2x+5),
\]
and the cubic has no root. These squarefree factorizations give a
$4$-cycle and a $3$-cycle in the transitive Galois group. A transitive
subgroup of $S_4$ with those elements is $S_4$: its order is divisible
by $12$, and the sole index-two subgroup $A_4$ has no $4$-cycle.
There are exactly two real roots, one in $(-1,0)$ and one in $(1,2)$;
for instance $4x^3-1$ has one real zero, the unique minimum of the
quartic, and the indicated sign changes give both roots. Thus actual
complex conjugation is a transposition. Label it $(12)$, with real
roots $3,4$.

Complex conjugation is central on $\Q^{\mathrm{cm}}$. Indeed for every
embedding $r$ of a CM field $E$, total reality of $E^+$ gives
$c r=r c_E$, where $c_E$ is its intrinsic involution. These identities
on all CM fields imply centrality on their compositum. Consequently
the Galois group of $L\cap\Q^{\mathrm{cm}}$ is a quotient of $S_4$ in
which $(12)$ is central. The normal subgroup generated by
$[g,(12)]$ is $A_4$: all these commutators are even, and taking a
transposition sharing one letter produces a $3$-cycle, whose conjugates
generate $A_4$. Thus this quotient has order at most two. It contains
the quadratic discriminant field $K=\Q(\sqrt{-283})$, itself a CM field.
We have proved
\[
              L\cap\Q^{\mathrm{cm}}=K.
\]
Restriction from $\Gal(\Qb/\Q^{\mathrm{cm}})$ onto $\Gal(L/K)=A_4$
is therefore surjective. Its commutator subgroup maps onto
$[A_4,A_4]=V_4$, while every such commutator fixes $F_{\mathrm{cm}}$
because $F_{\mathrm{cm}}/\Q^{\mathrm{cm}}$ is abelian. Choose a finite
product of these commutators restricting to $(12)(34)$. This supplies
$\sigma\in\Gal(\Qb/F_{\mathrm{cm}})$ exchanging the two real roots.
This argument needs no description of $L\cap F_{\mathrm{cm}}$.

If $v$ sends $\alpha$ to root $3$, then $w=\sigma v$ sends it to root
$4$. Apply the comparison already constructed to every object after
$v$. It compares precisely its $v$- and $w$-realizations. Finally the
Artin object $\operatorname{Spec}k(\sqrt\alpha)$ has geometric
real-Frobenius trace zero at the negative root and two at the positive
root. Any comparison intertwining those involutions would preserve
trace. The plain comparison preserves the ordinary Hodge structure
of this Artin object, which has type $(0,0)$, but cannot preserve that
extra involution. This proves the last assertion.
\end{proof}

Theorem~\ref{thm:quartic-CM-comparison} is confined to its specified
CM-abelian and Artin diagram. Its automorphisms have not been extended
to the full integral diagram in Proposition~\ref{prop:exact-transporter}.
It therefore settles neither universal number-field separation nor
the full-diagram comparison for these particular quartic places.

\subsection{A finite fixed-field isogeny skeleton}
\begin{proposition}\label{prop:isogeny-skeleton}
Let $A/k$ be an abelian variety over a number field. In the category
$\mathcal I_k(A)$ of abelian varieties $k$-isogenous to $A$, with
$k$-homomorphisms, there is a finite labelled diagram such that
restriction to it bijects natural integral $H_1$ comparisons at any
two embeddings with comparisons on that finite diagram.
\end{proposition}
\begin{proof}
Finiteness I gives representatives $B_1,\ldots,B_s$ for the finitely
many $k$-isomorphism classes in this fixed $k$-isogeny class. This
precise arithmetic theorem is stated in
\cite[Introduction, p.~1; see also \S7, pp.~10--11]{MasserWustholzIsogenies},
where the known Faltings finiteness input and the isogeny-estimate
approach are distinguished. We use the theorem as a named input.
For each $(i,j)$, a complex embedding injects $\operatorname{Hom}_k(B_i,B_j)$ into
the finite free abelian group of maps between their integral homology
lattices. It is consequently free of finite rank. Choose a $\Z$-basis
of each such Hom group. The finite diagram consists of the $B_i$
and those basis arrows.

Write $F_u(B)=H_1(B_u,\Z)$ and suppose $T_i:F_v(B_i)\to F_w(B_i)$
is a comparison on that diagram. For each $B$ choose a $k$-isomorphism
$b_B:B\to B_i$ and set
\[
             T_B=F_w(b_B)^{-1}T_iF_v(b_B).
\]
Any second choice differs by a $k$-automorphism of $B_i$. Commutation
with the basis arrows implies commutation with every integral linear
combination, hence with that automorphism. The definition is independent
of the choice. For $f:B\to C$, the conjugate $b_Cfb_B^{-1}$ is another
linear combination of the chosen Hom basis; it follows that
$T_CF_v(f)=F_w(f)T_B$. This proves naturality. Conversely naturality
for $b_B$ forces the displayed formula, giving uniqueness.
\end{proof}

\subsection{Geometric cyclic quotients force a projective Tate condition}
The preceding finite skeleton has fixed geometric rank. The full diagram
over $k$ also includes varieties defined over arbitrarily large finite
extensions and then regarded as $k$-schemes. Those objects impose the
following further necessary condition.

\begin{proposition}\label{prop:projective-Tate-necessity}
Let $A/k$ have dimension $g>0$, and let $\eta$ be a tensor comparison
between the full integral Betti diagrams at $v,w:k\hookrightarrow\C$.
Let $\sigma\in\Gal(\Qb/\Q)$ be the Artin element selected by
Proposition~\ref{prop:exact-transporter}, so that $\sigma v=w$.
Put $\Lambda_u=H_1(A_u,\Z)$ and
$T_A=(\eta_A^{-1})^\vee:\Lambda_v\to\Lambda_w$.
For every $n$ define
\[
 \kappa_{u,n}:\Lambda_u/n\Lambda_u\longrightarrow A_u[n](\C),
       \qquad x\longmapsto x/n\bmod\Lambda_u,
\]
and let
$\rho_{A,n}(\sigma)=\kappa_{w,n}^{-1}\sigma\kappa_{v,n}$.
Their compatible inverse limit is $\rho_A(\sigma)$.
There is a unique $\lambda_A\in\widehat\Z^\times$ such that
\[
                  T_A\otimes\widehat\Z
                         =\rho_A(\sigma)\lambda_A.
\]
\end{proposition}
\begin{proof}
Fix $n\ge2$ and a cyclic subgroup $H\subset A_v[n]$. All its points
are defined over a finite extension $L/k$ with an embedding $u$ extending
$v$. The quotient isogeny $q:A_L\to B=A_L/H$ exists over $L$:
finite subgroup quotients descend to their field of definition
\cite[Chapter~IV, Lemma~2.1, p.~135]{MilneAV}.
Regard these as $k$-schemes and include $A_L^{\mathrm{Art}}=
h^0(\operatorname{Spec}L)$ and its module actions. The primitive
idempotent $e_u$ selects the original homology block, and $\eta$ carries
it to $e_{\sigma u}$. On $A_L=A\times_k\operatorname{Spec}L$, tensor
compatibility makes the selected map exactly $T_A$. The graph of $q$
is a correspondence on the diagonal open and closed product component.
Naturality yields an integral square
\[
             T_Bq_*=(q_{\sigma H})_*T_A
\]
with $T_B$ an integral isomorphism between the selected quotient blocks.

Under $\kappa_{v,n}$, the kernel of $q_*\bmod n$ is exactly $H$.
Indeed $q(x/n)$ is zero precisely when $q_*(x)\in nH_1(B_u,\Z)$.
This remains true when the order of $H$ properly divides $n$. The square
therefore says that $T_A\bmod n$ and $\rho_{A,n}(\sigma)$ act in the
same way on every cyclic subgroup of $\Lambda_v/n\Lambda_v$.

An invertible matrix on $(\Z/n)^r$, $r\ge2$, preserving every cyclic
subgroup is scalar. The coordinate cyclic subgroups force the matrix
to be diagonal, its diagonal entries are units, and the subgroup
generated by $e_i+e_j$ forces the $i$-th and $j$-th entries to agree.
Apply this to $\rho_{A,n}(\sigma)^{-1}(T_A\bmod n)$ with $r=2g$.
We obtain a unique $\lambda_n\in(\Z/n)^\times$.
If $n\mid m$, reduction of the scalar at level $m$ is the scalar at
level $n$ by uniqueness. The maps $\kappa_{u,n}$ have the same
compatibility, using multiplication by $m/n$ on $m$-torsion. Thus the
$\lambda_n$ form the asserted unique $\lambda_A$.
\end{proof}

All objects in this argument belong to the full diagram over the fixed
field $k$. The underlying $k$-scheme of $B/L$ has $[L:k]$ geometric
components and degree-one rank $2g[L:k]$, so these quotients are not
confined to $\mathcal I_k(A)$. The proposition imposes a necessary
projective condition on the one Artin element selected by the entire
comparison. It makes no assertion that its scalar is independent of
$A$, or that an arbitrary projective Tate condition is sufficient for
a full tensor comparison.

\section{Cycle, Tate and Lefschetz normalizations}\label{sec:normalizations}
The comparison triangle is sensitive to the direction of comparison and to the Tate convention. At an isolated transverse intersection with a codimension-$p$ subvariety of a smooth $n$-fold, the normal transversal has complex dimension $p$: its tangent space is complementary to a space of dimension $n-p$. We collect the normalization statements needed to keep those choices consistent.

\begin{proposition}\label{prop:cycle-normalization}
Let $X/\C$ be smooth proper and let $z$ be an integral cycle of codimension $p$. With complex orientation and residue-normalized algebraic cycle class,
\[
 \alpha_X(\operatorname{cl}_{\B}z)
          =(2\pi i)^{-p}\operatorname{cl}_{\dR}z.
\]
If $f:X\to Y$ is proper between smooth varieties of dimensions $d,e$, put $p=e-d$. In every degree, including odd degrees and negative $p$,
\[
 \alpha_Y f_*^{\B}=(2\pi i)^{-p}f_*^{\dR}\alpha_X.
\]
A codimension-$r$ correspondence from $X$ to $Y$ sends degree $m$ to degree $m+2(r-d)$ and satisfies
\[
                \alpha_Y z_*^{\B}=(2\pi i)^{d-r}z_*^{\dR}\alpha_X.
\]
It is compatible without an additional factor when $r=d$, and in general after the target is twisted by $(r-d)$.
\end{proposition}
\begin{proof}
For a smooth divisor with local equation $z_1=0$, integration of $dz_1/z_1$ around its positive normal circle is $2\pi i$, whereas the topological Thom class has value one. Products of normal coordinates give $(2\pi i)^p$ in codimension $p$. To extend across a singular cycle support $Z$, let $S$ be its singular locus. Every component of $S$ has codimension at least $p+1$ in $X$. Poincare duality with supports identifies $H_S^{2p}$ and $H_S^{2p+1}$ with Borel--Moore homology in degrees exceeding the real dimension of $S$; both vanish. The support localization sequence thus identifies $H_Z^{2p}(X)$ with $H_{Z\setminus S}^{2p}(X\setminus S)$. Using ordinary algebraic de Rham comparison \cite[Theorems 1 and 1', pp.~95--96]{GrothendieckDeRham}, its version with supports follows from the comparison of the localization triangles. Consequently the two fundamental classes extend uniquely from the smooth locus, and the local factor proves the cycle formula globally. Multiplicities are handled by additivity.

For proper $X$, normalize algebraic trace to give one on the class of a point. The cycle formula gives
\[
 \int_X c_X(x)=(2\pi i)^d\operatorname{Tr}_{X,\dR}(x),\qquad c_X=\alpha_X^{-1}.
\]
For $x\in H^r_{\dR}(X)$ and $y\in H^{2d-r}_{\dR}(Y)$, trace duality and the projection formula yield
\[
 \int_Y c_Y(f_*^{\dR}x)c_Y(y)
   =(2\pi i)^{e-d}\int_X c_X(x)f^*c_Y(y).
\]
Nondegeneracy proves the Gysin formula. If the smooth varieties are not proper, use compactly supported complementary cohomology; properness of $f$ makes the needed pullback preserve compact support. Relative codimensions add under composition. Finally a correspondence is pullback, cup product with its cycle class, and projection from $X\times Y$ to $Y$. The last operation has relative codimension $-d$, so the two factors multiply to $(2\pi i)^{d-r}$. The target Tate line cancels this factor. The diagonal acts identically and the projection formula proves compatibility with composition.
\end{proof}

The position of the rational Betti form and the $k$-de Rham form in their common complex space is precisely the period point. An integral de Rham lattice additionally requires an integral model. Even the standard projective-line model carries the factor $(2\pi i)^{-1}$ on its degree-two generator, as the cycle formula shows.

\subsection{Changing the archimedean identification}
Let $K$ be an abstract field isomorphic to $\C$ and choose an isomorphism $\sigma:\C\to K$. Scalar transport gives the cycle comparison coefficient $\sigma(2\pi i)^{-p}$. Replacing $\sigma$ by $\sigma c$ multiplies it by $(-1)^p$ and preserves the embedded integral Betti lattice: conjugation on complex points is itself an integral cohomology isomorphism, and comparison commutes with conjugation. The sign is degree-dependent; on $\mathbf P^2$ conjugation negates the hyperplane class and fixes its square.

Different identifications can give different lattices. Fix a transcendence basis containing $\pi$, send $\pi$ to $2\pi$, fix $i$ and the remaining basis elements, and extend to an automorphism of $\C$. On $H^2_{\dR}(\mathbf P^1)=Kv$, the two resulting lattices are $(2\pi i)^{-1}\Z v$ and $(4\pi i)^{-1}\Z v$, which are unequal. More generally the images of an oriented integral generator, as $\sigma$ varies, are exactly $av$ with $a$ transcendental over $\Q$. Every such coefficient occurs because a bijection between transcendence bases sending $2\pi i$ to $a^{-1}$ extends to an isomorphism of algebraic closures. Algebraic nonzero coefficients never occur. Primitive elements refer to the two generators of one chosen lattice, not to all nonzero vectors of its ambient line.

For a real smooth proper $X$, put $L=H^r(X(\C),\Z)/\mathrm{torsion}$. Its diagonal image under Betti--etale and de Rham comparison in
\[
 \left(\prod_\ell H^r_{\mathrm{et}}(X_\C,\Z_\ell)/\mathrm{torsion}\right)
 \times H^r_{\dR}(X/\R)\otimes_\R\C
\]
is stable under real Galois action. If $f_\infty$ is the integral involution on $L$, finite-coefficient comparison gives $c_\ell f_\infty=c\,c_\ell$, and
\[
       (1\otimes c)\alpha=\alpha(f_\infty\otimes c).
\]
Thus conjugating the diagonal image of $x$ gives the diagonal image of $f_\infty x$. Stability does not assert that the involution is trivial.

Only two identifications of an abstract topological complex field with $\C$ are continuous: any continuous field automorphism fixes $\Q$, hence $\R$ by density, and sends $i$ to $i$ or $-i$. Thus its two possibilities are the identity and conjugation. The comparison square between the two associated complex varieties uses both the antiholomorphic identification of points and conjugation of scalar coefficients. Each operation exchanges $(p,q)$ and $(q,p)$ once; their composite preserves the labelled summands and the de Rham filtration. Consequently the Hodge splitting transported to a fixed topological complex field is independent of these two choices. This coefficient-conjugating identification is not an additional complex-linear symmetry preserving the labelled filtration.

This also specifies the scalar field in an intersection of realizations. If $v:k\hookrightarrow\C$, put $R_v=v^{-1}(\R)$. A Hodge component intersected with de Rham cohomology is a $k$-space, since the component is a complex vector subspace; intersecting further with the transported Betti real form gives an $R_v$-space. If the latter intersection $W$ is nonzero, its scalar stabilizer inside $k$ is exactly $R_v$: for $0\ne w\in W$, both $w$ and $aw$ are Betti-real only if $v(a)=\overline{v(a)}$. The zero intersection has stabilizer all of $k$ and must be excluded from this assertion.

For $k=\Qb$, the intersection of the fields $R_v$ over all embeddings is the field of totally real algebraic numbers. For any characteristic-zero field $k$ admitting a complex embedding, the same intersection consists precisely of its algebraic totally real elements. Indeed a transcendental element extends to a transcendence basis and can be sent to a nonreal transcendental number; choosing the remaining images algebraically independently and extending through the algebraic extension gives an embedding of $k$. Every embedding of the number field of an algebraic element likewise extends to $k$, so that element lies in the intersection exactly when all its conjugates are real. For a number field this is its largest totally real subfield. These statements concern scalar fields of real intersections, not fields generated by the coefficients of a period matrix.


\subsection{Filtration shifts and the integral Lefschetz distinction}
For increasing weight filtration and decreasing Hodge filtration, the Tate convention is
\[
 W_i(E(n))=W_{i+2n}(E)\otimes\Q(n),\qquad
 F^a(E(n))=F^{a+n}(E)\otimes\C,
\]
so $\gr_i^W(E(n))=(\gr_{i+2n}^WE)(n)$. This follows directly from the weight $-2n$ and type $(-n,-n)$ of the Tate line, as in Deligne's definitions \cite[(2.1.13)--(2.1.14), pp.~26--27]{DeligneHodgeII}. An unshifted graded formula fails already for $E=\Q(0)$ and $n=-1$.

\begin{proposition}\label{prop:lefschetz-integral}
On $V_m=\bigoplus_{j=0}^m\Q v_j$ define, with zero at the endpoints,
\[
 Ev_j=v_{j+1},\quad \Lambda v_j=v_{j-1},\quad
 Fv_j=j(m-j+1)v_{j-1},\quad Hv_j=(2j-m)v_j.
\]
Then $[E,F]=H$, $[H,E]=2E$ and $[H,F]=-2F$. On $\bigoplus_{m=0}^dV_m$,
\[
 \Q[E,\Lambda]=\Q[E,F,H]=\prod_{m=0}^d\operatorname{Mat}_{m+1}(\Q).
\]
On the displayed integral lattice in $V_2$, however,
\[
               \Z[E,F,H]\ne\Z[E,\Lambda]=\operatorname{Mat}_3(\Z).
\]
\end{proposition}
\begin{proof}
Evaluate the commutators on each basis vector, including the two endpoints. Put $B=1-E\Lambda$ and $T=1-\Lambda E$, the projections to the bottom and top of each string. Then
\[
                    B_m=B\Lambda^mTE^mB
\]
projects to the bottom of the length-$m$ string and kills all other strings: $E^m$ kills shorter strings, and $T$ kills its image in longer ones. Hence $E^bB_m\Lambda^a$ are matrix units for the $m$th factor. This proves the first operator-algebra description, over $\Z$ on one string as well.

For the other description, $C=H^2+2H+4FE$ acts on $V_m$ by $m(m+2)$. Interpolation in these distinct scalars projects to individual strings; interpolation in $H$ there projects to individual $v_j$. On each such line, $F$ is a nonzero rational multiple of $\Lambda$ away from the bottom. This recovers $\Lambda$ from $E,F,H$. The matrix-unit description also expresses $F,H$ in terms of $E,\Lambda$.

On $V_2$ both $F$ and $H$ reduce to zero modulo two, so the reduction of $\Z[E,F,H]$ is contained in $\mathbf F_2[E]$. Every such operator preserves the line $\mathbf F_2v_2$, whereas $\Lambda v_2=v_1$. Thus $\Lambda$ does not belong to that integral algebra. No integral equality follows from the rational interpolation.
\end{proof}

Weight scaling and the Lefschetz Cartan scaling also differ as maps, despite having the same conjugation on degree-changing operators. On $H^i$ of a smooth projective $d$-fold they are $w(t)=t^i$ and $\lambda(t)=t^{i-d}$, so $\lambda=t^{-d}w$. Both conjugate an operator of degree $j$ by $t^j$. On $H^*(\mathbf P^1)$, the Tate action $\rho(z)=\operatorname{diag}(1,z)$ satisfies $\rho E\rho^{-1}=zE$ and $\rho F\rho^{-1}=z^{-1}F$, but the nontrivial unipotent $\exp(sE)$ cannot factor through that diagonal action. On elliptic $H^1$, the Lefschetz raising and lowering maps have zero targets, and its Cartan operator is zero. This action is therefore distinct from the nontrivial standard $\SL_2$ representation on the elliptic period space.

\section{Descent and comparison orbits}\label{sec:descent-boundaries}
These examples separate an actual descent datum or cohomological comparison from weaker data that can resemble one.

\subsection{A discontinuous semilinear action}
\begin{proposition}
An abstract automorphism-indexed semilinear action on a finite etale scheme need not be effective algebraic descent data.
\end{proposition}
\begin{proof}
Let $q_1,q_2,\ldots$ be distinct primes and $L=\bigcup_n\Q(\sqrt{q_1},\ldots,\sqrt{q_n})$. Independence of square classes gives $\Gamma=\Gal(L/\Q)=\prod_{n\ge1}\mathbf F_2$. Its finite-support subgroup $D=\bigoplus_{n\ge1}\mathbf F_2$ is dense. The all-ones vector is nonzero in $\Gamma/D$; extend it to a vector-space basis and take its coordinate functional. This gives a nonzero homomorphism $\eta:\Gamma\to\mathbf F_2$ vanishing on $D$, hence discontinuous.

On $L\times L$, let $\sigma$ act on coefficients and additionally swap the factors when $\eta(\sigma)=1$. These operations commute, and the homomorphism identity proves the semilinear action law. If this action came by base change from a $\Q$-scheme, faithful-flat descent of finite etale morphisms would make that scheme finite etale of degree two. Its two geometric points are defined over a finite extension, so their permutation action has open kernel. The specified permutation action has the dense proper kernel $\ker\eta$, a contradiction. This example concerns abstract actions, not an fpqc cocycle satisfying the hypotheses of effective descent.
\end{proof}

\subsection{A hypersurface model does not preserve ordinary cohomology}
\begin{proposition}\label{prop:hypersurface-normalization}
Every smooth integral projective $d$-fold $X$ over a characteristic-zero field admits a finite birational morphism to an integral hypersurface $Y\subset\mathbf P^{d+1}$, identifying $X$ with the normalization of $Y$. Ordinary cohomology need not be preserved.
\end{proposition}
\begin{proof}
Choose a very ample line bundle $L$. Over the infinite ground field, choose $d+1$ sections $s_0,\ldots,s_d$ with no common zero: successively avoid the generic points of the previous intersections, and choose the last section nonvanishing at their finite residual intersection. The map $g=[s_0:\cdots:s_d]$ has $g^*\mathcal O(1)=L$. A positive-dimensional geometric fibre would contain a projective curve, on which $L$ would be both ample and trivial. Thus $g$ is quasi-finite, hence finite by properness \cite[Tag 02LS]{StacksFinite}. Its image is all of $\mathbf P^d$, so $F=k(s_1/s_0,\ldots,s_d/s_0)\subset k(X)$ is finite separable.

Complete the sections to an embedding system. Their ratios generate $k(X)$ over $F$. A $k$-linear combination $r$ of them is a primitive element: each pair of distinct $F$-embeddings of $k(X)$ excludes a proper $k$-linear subspace of coefficients, and finitely many such subspaces cannot cover $k^N$. If $r=t/s_0$, the map $[s_0:\cdots:s_d:t]$ is finite by the same ampleness argument and birational onto its image. That image has codimension one in $\mathbf P^{d+1}$; its homogeneous prime ideal has height one in a polynomial UFD, hence is principal. It is therefore an integral hypersurface. On an affine open, its inverse-image ring is finite, normal and has the same fraction field. It is the integral closure, since every element integral over the smaller ring is integral over it and hence belongs to it. This proves the normalization assertion.

For failure of cohomological invariance, take the nodal cubic
\[
 Y:\ y^2z=x^2(x+z),\qquad
 \nu([u:v])=[v(u^2-v^2):u(u^2-v^2):v^3].
\]
The formula is everywhere defined and birational, with the two points $[1:1]$, $[-1:1]$ over the node $P$. The normalization sequence on the analytic spaces is
\[
 0\to\Q_Y\to\nu_*\Q_{\mathbf P^1}\to\Q_P\to0,
\]
where the last map is the difference of the two branch values. A finite proper map has no higher direct images here, by proper base change on its finite fibres. The global-constant map to $\Q_P$ is zero, so its cohomology sequence gives $H^1(Y,\Q)=\Q$ and $H^1(\mathbf P^1,\Q)=0$. The pullback kills precisely this displayed one-dimensional class.
\end{proof}

\subsection{The actual comparison cocycle and its orbit closure}
Let $F\subset\C$, let $D$ be a de Rham fibre functor over $F$, and let $B$ be a rational Betti fibre functor in a neutral category with a tensor generator. In this subsection write comparison in the direction
\[
 p:D_\C\longrightarrow B_\C,\qquad
 P=\Isom^\otimes(D,B_F),\qquad G_D=\Aut^\otimes(D).
\]
Thus $P$ is a right $G_D$-torsor over $F$ \cite[Theorem~3.2, p.~31]{DeligneMilne}. The two cocycle conventions are
\[
 \varphi(\sigma)=(\sigma p)^{-1}p,\qquad
 \psi(\sigma)=p^{-1}\sigma p=\varphi(\sigma)^{-1}.
\]
Typed multiplication gives
\[
 \varphi(\tau\sigma)=\tau(\varphi(\sigma))\varphi(\tau),
 \qquad
 \psi(\tau\sigma)=\psi(\tau)\tau(\psi(\sigma)).
\]
Consequently $\sigma\mapsto(\psi(\sigma),\sigma)$ is a homomorphism into the semidirect product with multiplication $(a,\tau)(b,\sigma)=(a\tau(b),\tau\sigma)$. Its first projection need not be a homomorphism. The relation is semilinear and cannot be replaced by an ordinary representation.

These formulas come from actual fpqc descent. Over $R=\C\otimes_F\C$, let $p_1,p_2$ be the two pullbacks of $p$ and set $z_{12}=p_1^{-1}p_2$. Over the triple tensor product,
$z_{12}z_{23}=z_{13}$. The trivializations $h\mapsto p_i h$ differ by this transition map. Evaluation $a\otimes b\mapsto a\sigma(b)$ sends $z_{12}$ to $\psi(\sigma)$. Thus the automorphism cocycle is the restriction of the descent datum of the already defined torsor; the ineffective abstract action above is not being used as a descent theorem. If a finite extension $F'/F$ supplies $p_0\in P(F')$ and $p=p_0g$, then for $\sigma$ fixing $F'$ one has $\varphi(\sigma)=(\sigma g)^{-1}g$. Such an extension exists because a nonempty affine finite-type scheme over a field has a closed point of finite residue degree. The extension is for trivialization and does not change the field used in the next genericity assertion.

\begin{proposition}\label{prop:orbit-closure}
Let $p\in P(\C)$ and let $Z$ be its Zariski locus over $F$. For the complex algebraic isomorphism
\[
 d_p:P_\C\longrightarrow G_{D,\C},\qquad x\longmapsto x^{-1}p,
\]
one has
\[
 \overline{\im\varphi}^{\rm Zar}=d_p(Z_\C),\qquad
 \dim\overline{\im\varphi}^{\rm Zar}=\trdeg_F F(p).
\]
In particular the actual cocycle image is dense in the full complex group exactly when $p$ is generic in $P$ over the same field $F$. This statement does not replace the actual image by the subgroup it generates.
\end{proposition}
\begin{proof}
Embed $P$ as an affine $F$-scheme. The ideal $J$ over $\C$ vanishing on the orbit $\Aut(\C/F)p$ is stable under coefficient automorphisms. In each finite-dimensional space of polynomials of bounded degree, its unique reduced row-echelon basis is fixed by every such automorphism. All its coefficients therefore lie in $F$. Here the fixed field of $\Aut(\C/F)$ is exactly $F$: an algebraic element outside $F$ has a distinct conjugate in characteristic zero, and a transcendental element can be moved after extending it to a transcendence basis; the resulting field embeddings extend to $\C$. It follows that $J$ is generated by $J\cap F[x]$, which is precisely the ideal of $Z$. Hence the orbit closure is $Z_\C$. Applying $d_p$ proves the formula. The coordinate domain of $Z$ has fraction field $F(p)$, so its dimension is the displayed transcendence degree. Neither argument needs geometric connectedness; all geometric components of $Z_\C$ remain in the formula.
\end{proof}

Already for the trivial $\Gm$-torsor over $\Q$ and $p=2+\sqrt3$, the actual image is $\{1,p^2\}$ because the other conjugate is $p^{-1}$. This finite set is not dense, whereas its generated subgroup is dense: $p^2$ has infinite order and a nonzero one-variable polynomial has finitely many roots. It also directly disproves an ordinary-homomorphism interpretation. Choose an automorphism $\sigma$ carrying $p$ to $p^{-1}$; then $\sigma^2p=p$, so $\varphi(\sigma^2)=1$, while $\varphi(\sigma)^2=p^4\ne1$. No order-two extension of the quadratic automorphism to $\C$ is needed.

Change of place has a similarly typed formula. Given a rational tensor isomorphism $a:B_v\to B_w$, put $g=p_w^{-1}ap_v$ on $D_\C$. Then
\[
 \sigma(g)\varphi_v(\sigma)g^{-1}
 = (\sigma p_w)^{-1}a(\sigma p_v)
   (\sigma p_v)^{-1}p_vp_v^{-1}a^{-1}p_w
 =\varphi_w(\sigma).
\]
It carries the associated rational structures, and carries their integral lattices exactly when $a$ is integral on those lattices. At a finite place, where the action on a frame is conjugation by $\rho_\ell(\sigma)$, this specializes to
$\rho_\ell g_\ell\rho_\ell^{-1}\rho_\ell g_\ell^{-1}=\rho_\ell$. The identity does not supply an integral comparison whose existence is contradicted in Section~\ref{sec:integral}.

\subsection{A motivic difference between the two closures}
The elementary distinction in Proposition~\ref{prop:orbit-closure} occurs for genuine period comparisons, even over a finitely generated field where $\pi$ remains transcendental.

\begin{proposition}\label{prop:cm-orbit-counterexample}
There is a CM elliptic degree-one motive over a field $F$ of transcendence degree one for which the actual comparison-cocycle image has a proper Zariski closure, while the subgroup generated by that image is dense in its two-dimensional motivic torus. One can add an elliptic summand so that transcendence degree one is the minimum field-of-definition degree in the effective degree-one category; the actual closure is still strictly smaller than the generated subgroup's closure. In both cases $\pi$ is transcendental over $F$.
\end{proposition}
\begin{proof}
Take $E_0:y^2=x^3-x$. Over a finite number field $K_1$ one may split its CM torus and choose a point of the comparison torsor: the torus splits after adjoining its CM field, and a closed point of a nonempty finite-type torsor has finite residue degree. In split eigenbases write its distinguished comparison point as $(u,v)\in\Gm^2(\C)$. Up to nonzero algebraic constants, $u$ is a holomorphic period or its inverse and $uv=(2\pi i)^{\epsilon}$ with $\epsilon=1$ or $-1$, according to direction. Proposition~\ref{prop:cm-h} gives algebraic independence of that holomorphic period and $\pi$. Hence $u,v$ are algebraically independent over $\Qb$.

Set $s=u+v$ and $F=K_1(s)$. The field $F$ has transcendence degree one, and $\pi$ is transcendental over it. Indeed if $uv$ were algebraic over $K_1(u+v)$, both $u$ and $v$ would be algebraic over that field by the quadratic $X^2-sX+uv$, contradicting their independence. A nonzero algebraic power or inverse does not change this argument for $\pi$.

The $F$-locus of $(u,v)$ in $\Gm^2$ is exactly $U+V=s$. Its coordinate domain is $F[U,U^{-1},(s-U)^{-1}]$, and $u$ is transcendental over $F$, proving that there are no further relations. For the left cocycle $\psi$ the actual closure is therefore the curve
\[
                       C:\quad ux+vy=s
                         \quad\hbox{in }(\Gm^2)_\C.
\]
It is proper and contains $(1,1)$. It generates the whole torus. To prove this, a proper algebraic subgroup containing the irreducible curve through the identity would impose a nontrivial character relation $x^ay^b=1$. On the parameter line $y=(s-ux)/v$, the valuations at $x=0$ and $x=s/u$ force $a=0$ and $b=0$, respectively. This is a contradiction. Equivalently, the character lattice of any proper subtorus has a nonzero annihilating character. Inversion gives the same strict actual-versus-generated conclusion for $\varphi$.

This CM motive itself descends to a number field. To impose the minimum-degree condition, use the elliptic curve
\[
 E_s:\quad y^2=x^3-\frac{3s}{s-1728}x-\frac{2s}{s-1728},
 \qquad M=h^1(E_0)\oplus h^1(E_s).
\]
Substitution in $j=1728\,4A^3/(4A^3+27B^2)$ gives $j(E_s)=s$; the discriminant is nonzero since $s$ is transcendental. Restriction to the CM summand gives a surjective group map $\rho:G_M\to T$, where $T$ is its CM torus. This is tensor reconstruction for the full subcategory generated by that summand \cite[Proposition~2.21(a), p.~26]{DeligneMilne}. Let $Y$ be the actual left-cocycle closure and $J$ the algebraic subgroup generated by the image. Then $\rho(Y)\subset C$, while $\rho(J)=T$: the image of an algebraic-group homomorphism is closed and contains the subgroup generated by the CM cocycle. Consequently $Y\ne J$. This does not assert $J=G_M$.

It remains to justify the minimum in the stated category. Effective rational degree-one motives of abelian varieties are abelian varieties up to isogeny \cite[Proposition~6.21, p.~65]{DeligneMilne}. If $E_0\times E_s$ descended up to isogeny to a number field, its elliptic isogeny factor corresponding to $E_s$ would be defined over $\Qb$. Here homomorphisms of abelian varieties do not increase under extension of algebraically closed fields: spread a homomorphism over an integral finite-type base; the rigidity lemma makes the family constant because it fixes the origin, and specialization to a closed point recovers it over the smaller algebraically closed field. Apply this to the rational idempotent defining that factor, clear its denominator, and take the image of the resulting algebraic homomorphism. The kernel of an isogeny from the resulting elliptic curve over $\Qb$ is a finite subgroup of its torsion, so it and its quotient are defined over $\Qb$. Its quotient could not have the transcendental $j$-invariant $s$. Thus descent to transcendence degree zero is impossible. The displayed model over $F$ proves the minimum is one.
\end{proof}

A nonzero Tate quotient gives a necessary condition for genericity: restriction of the torsor to $\Q(-m)$, $m\ne0$, is dominant, so a generic comparison point has a transcendental Tate coordinate $(2\pi i)^{\pm m}$. In particular $\pi$ must be transcendental over $F$. Proposition~\ref{prop:cm-orbit-counterexample} proves that this condition is insufficient even with the stated minimum-degree condition. Conversely minimum degree alone does not imply the condition: the same elliptic equation with $s=\pi$ has minimum field degree one over $\Q(\pi)$, but its determinant comparison is algebraic over $\Q(\pi,i)$, excluding full genericity. These examples concern transcendental fields of definition and do not refute a period conjecture over $\Qb$.

\subsection{The complex topology and the total image}
If $F\subset\C$ is countable, the orbit formula has an exact complex-topological version:
\[
 \overline{\Aut(\C/F)p}^{\rm Eucl}=Z_\C(\C),\qquad
 \overline{\im\varphi}^{\rm Eucl}=d_p(Z_\C)(\C).
\]
Indeed the orbit is exactly the set of $F$-generic points of $Z$. Their coordinate fields are isomorphic over $F$, and this isomorphism extends to $\C$ by algebraic closures and transcendence bases. There are countably many proper $F$-subvarieties of $Z$. Each has empty interior on every geometric irreducible component: the components are conjugate, so an $F$-closed subset containing one would contain all and equal $Z$. On a complex variety a proper closed subset of a component has empty interior, as follows on its smooth locus from the identity theorem in analytic charts. Baire's theorem on those locally compact smooth loci shows that the complement of this countable union is dense, and the smooth loci are themselves dense. Applying $d_p$ proves the second formula. No field automorphism is presumed continuous for the usual complex topology.

For a number-field motive write $\Gamma_f$ for the compact finite-adelic Galois image. If its comparison point is generic over $\Qb$, the closure of the combined finite and complex image is
\[
                           \Gamma_f\times G(\C).
\]
In fact $\Aut(\C/\Qb)$ acts trivially on the finite realization. In the fibre over any attained finite value, compose a lift with these automorphisms; the preceding orbit formula makes its complex image dense. The product is closed and already contains the total image, proving both inclusions.

For $\Q(-2)$ over $\Q$, the comparison coordinate $(2\pi i)^{\pm2}$ is transcendental, and the finite image is the closed non-open subgroup $(\widehat\Z^\times)^2$ described above. Its total closure is therefore
\[
             (\widehat\Z^\times)^2\times\C^\times.
\]
It is not open in $\widehat\Z^\times\times\C^\times$. Thus universal openness of the total-image closure fails even over the minimum field $\Q$ and with a generic Tate comparison point.

For a non-CM elliptic $h^1$ over a number field, the exact remaining assertion is different: the total-image closure is an open subgroup if and only if the period point is fully generic. In the forward direction project this closed subgroup to $\GL_2(\C)$. Projection along the compact finite factor is closed, and the projection of an open subgroup is open. Since $\GL_2(\C)$ is connected, its image is the whole group. The complex orbit formula then forces the period locus over the number field to be the full torsor; finite base extension preserves the transcendence degree, so the point is generic over $\Qb$. Conversely genericity gives the product closure and Serre's theorem makes the finite factor open. This equivalence leaves the non-CM elliptic period conjecture as an explicit unresolved arithmetic input.

\section{Questions left by the dimension formula}
The exact image calculation does not decide the period conjecture for the joint motive. For $v=u$, proving that the comparison point is generic remains the arithmetic obstruction to unconditional equality in the table. Even in rank two, a dominant map of algebraic parameter spaces does not turn genericity of one image into genericity of the joint point.

For a motive isomorphic to its conjugate, $v=0$, and genericity for $M$ alone would give every entry of the corrected table. For arbitrary $v$, genericity for $M$ alone gives only the lower bounds
\[
              \trdeg_k k(\beta)\ge u-v,\qquad
              \trdeg_k k(\gamma)\ge b,
\]
by applying the relative upper bounds to $\trdeg_k K\ge g$. Whether other arithmetic input forces larger degrees is a separate question.

The restricted-Lie formula turns a proposed algebraic elliptic splitting parameter $a$ into explicit ordinary congruences $K_p+\overline aH_p=0$. No argument here proves that algebraicity of the complex Hodge splitting imposes these congruences. In particular the finite-field identity is not a proof of an elliptic transcendence converse. The ordinary Hodge--Tate example shows why one tempting replacement of that missing implication is false.

The non-geometric Hodge pair gives an exclusion, not a characterization of the geometric essential image. The integral examples concern the stated correspondence diagrams. They do not classify integral comparisons for every field. Proposition~\ref{prop:enriched-embeddings} settles the separate equality test when all labelled integral elliptic Hodge structures are retained. Their obstruction is precise: finite and profinite comparisons need not preserve one simultaneous integral lattice. A general recovery assertion must specify the source category, coefficient ring and equivalence relation on places. Fixed-object genericity or rational commensurability alone does not recover a place from the entire labelled diagram.

\section*{Method and records}
The mathematics and text were produced by GPT-6 Astra in ChatGPT Pro mode. Checks of the main dimension theorem include a separate derivation by the same model from its statement, and the computational checks use exact arithmetic. The record does not identify a more specific effort setting. No human has checked the mathematics or text. The authors dedicate this text and their accompanying programs to the public domain under CC0 1.0. The designated location for programs and detailed records is \expandafter\url\expandafter{\workbenchurl}; the permanent catalogue address of the records is \expandafter\url\expandafter{\recordurl}.

\begin{thebibliography}{99}
\bibitem{Grothendieck} A. Grothendieck, \emph{Motifs et th\'eorie de Hodge}, Montpellier manuscripts, cote 18, handwritten p.~10 (PDF p.~11), and PDF pp.~13 and 55; and cote 15, PDF pp.~221--223. Universit\'e de Montpellier, \url{https://grothendieck.umontpellier.fr/archives-grothendieck/}.
\bibitem{Andre} Y. Andr\'e, \emph{Pour une th\'eorie inconditionnelle des motifs}, Publ. Math. IH\'ES \textbf{83} (1996), 5--49. \url{https://www.numdam.org/article/PMIHES_1996__83__5_0.pdf}.
\bibitem{BrownDupont} F. Brown and C. Dupont, \emph{Single-valued integration and double copy}, J. reine angew. Math. \textbf{775} (2021), 145--196. Locators use arXiv:1810.07682v3. \url{https://arxiv.org/abs/1810.07682}.
\bibitem{BostCharles} J.-B. Bost and F. Charles, \emph{Some remarks concerning the Grothendieck period conjecture}, J. reine angew. Math. \textbf{714} (2016), 175--208. Locators use the author's 34-page advance-publication PDF. \url{https://www.math.ens.psl.eu/~charles/crelle-periods.pdf}.
\bibitem{KontsevichZagier} M. Kontsevich and D. Zagier, \emph{Periods}, in Mathematics Unlimited---2001 and Beyond, Springer, 2001, 771--808, \S4. \url{https://www.ihes.fr/~maxim/TEXTS/Periods.pdf}.
\bibitem{HuberMuellerStach} A. Huber and S. M\"uller-Stach, \emph{Periods and Nori Motives}, Springer, 2017; author-annotated version dated 1 July 2025. \url{https://home.mathematik.uni-freiburg.de/arithgeom/preprints/buch-errata/buch-errata.pdf}.
\bibitem{BakkerTsimerman} B. Bakker and J. Tsimerman, \emph{Functional transcendence of periods and the geometric Andr\'e--Grothendieck period conjecture}, Forum Math. Sigma \textbf{13} (2025), e97. \url{https://doi.org/10.1017/fms.2025.10036}.
\bibitem{Ayoub} J. Ayoub, \emph{Anabelian presentation of the motivic Galois group in characteristic zero}, Forum Math. Sigma \textbf{14} (2026), e139. Author version, Theorem 4.4.2. \url{https://user.math.uzh.ch/ayoub/PDF-Files/Anabel.pdf}.
\bibitem{Hua} M. Hua, \emph{Grothendieck Montpellier. \'Enonc\'es des lectures modernis\'ees d\'emontr\'es}, version 2, 5 November 2026. \url{https://hexagonmath.org/2611.00125v2}.
\bibitem{Howe} S. Howe, \emph{Transcendence of the Hodge--Tate filtration}, J. Th\'eor. Nombres Bordeaux \textbf{30} (2018), 671--680. \url{https://arxiv.org/abs/1610.05242}.
\bibitem{Ogus} A. Ogus, \emph{$F$-crystals, Griffiths transversality, and the Hodge decomposition}, Ast\'erisque \textbf{221} (1994), ii+183 pp. \url{https://www.numdam.org/item/AST_1994__221__1_0.pdf}.
\bibitem{Totaro} B. Totaro, \emph{Hodge structures of type $(n,0,\ldots,0,n)$}, Int. Math. Res. Not. (2015). Locators use the author version. \url{https://www.math.ucla.edu/~totaro/papers/public_html/structure.pdf}.
\bibitem{Yoshinaga} M. Yoshinaga, \emph{Periods and elementary real numbers}, arXiv:0805.0349, 2008, Theorem 18, p.~11. \url{https://arxiv.org/abs/0805.0349}.
\bibitem{Schneider} Th. Schneider, \emph{Arithmetische Untersuchungen elliptischer Integrale}, Math. Ann. \textbf{113} (1937), 1--13. \url{https://doi.org/10.1007/BF01571618}.
\bibitem{MilneMF} J. S. Milne, \emph{Modular Functions and Modular Forms}, version 1.31, 22 March 2017. \url{https://www.jmilne.org/math/CourseNotes/MF.pdf}.
\bibitem{Milne} J. S. Milne, \emph{The fundamental theorem of complex multiplication}, arXiv:0705.3446, 2007. \url{https://arxiv.org/abs/0705.3446}. Erratum to the proof of Theorem 2.1: \url{https://jmilne.org/math/articles/2007c.html}.
\bibitem{LocalSystems} Y. Zhang, lecture notes by M. Shi, \emph{Local systems in algebraic geometry}, in Ohio State University conference notes, 2024, Example 5.2, pp.~16--17. \url{https://people.math.osu.edu/cueto.5/RTG/rtg24/Local_Systems_in_Algebraic_Geometry_2024_Notes-Final.pdf}.
\bibitem{Landesman} A. Landesman, \emph{Invariance of the tame fundamental group under base change between algebraically closed fields}, arXiv:2005.09690, version 3, 2024, Theorem 1.1. \url{https://arxiv.org/abs/2005.09690}.
\bibitem{GrothendieckDeRham} A. Grothendieck, \emph{On the de Rham cohomology of algebraic varieties}, Publ. Math. IH\'ES \textbf{29} (1966), 95--103. \url{https://www.numdam.org/article/PMIHES_1966__29__95_0.pdf}.
\bibitem{GrothendieckTateLetter} A. Grothendieck, \emph{Letter to John Tate}, May 1966, original thirty-page typescript, \S\S3.1--3.2, pp.~29--30. \url{https://csg.igrothendieck.org/wp-content/uploads/2024/10/L9u.pdf}.
\bibitem{GrothendieckNice} A. Grothendieck, \emph{Groupes de Barsotti--Tate et cristaux}, Actes du Congr\`es International des Math\'ematiciens (Nice, 1970), vol.~1, Gauthier-Villars, 1971, 431--436. \url{https://webusers.imj-prg.fr/~leila.schneps/grothendieckcircle/GroupesB-T.pdf}.
\bibitem{Katz} N. M. Katz, \emph{$p$-adic interpolation of real analytic Eisenstein series}, Ann. of Math. \textbf{104} (1976), 459--571. \url{https://web.math.princeton.edu/~nmk/old/padicinterp.pdf}.
\bibitem{Nesterenko} Yu. V. Nesterenko, \emph{Modular functions and transcendence questions}, Sbornik: Mathematics \textbf{187} (1996), 1319--1348. Theorem 1, p.~1320. \url{https://doi.org/10.1070/SM1996v187n09ABEH000158}.
\bibitem{SebbarBesrour} A. Sebbar and K. Besrour, \emph{Zeros of Katz's $p$-adic Eisenstein series $E_2$}, arXiv:2611.14502v1, 10 November 2026. Lemma 8.1, pp.~25--26. \url{https://arxiv.org/abs/2611.14502}.
\bibitem{DeligneHodgeII} P. Deligne, \emph{Th\'eorie de Hodge, II}, Publ. Math. IH\'ES \textbf{40} (1971), 5--57. \url{https://www.numdam.org/article/PMIHES_1971__40__5_0.pdf}.
\bibitem{StacksFinite} The Stacks Project Authors, \emph{The Stacks Project}, Lemma 37.44.1, Tag 02LS. \url{https://stacks.math.columbia.edu/tag/02LS}.
\bibitem{TateLocal} J. Tate, \emph{$p$-divisible groups}, in \emph{Proceedings of a Conference on Local Fields}, Springer, 1967, 158--183. \url{https://www.math.purdue.edu/~tongliu/teaching/598/p-divisible.pdf}.
\bibitem{KisinLocal} M. Kisin, \emph{Crystalline representations and $F$-crystals}, in \emph{Algebraic Geometry and Number Theory}, Birkh\"auser, 2006, 459--496. Locators use the author's 34-page version. \url{https://people.math.harvard.edu/~kisin/dvifiles/wam.dvi}.
\bibitem{KimLocal} W. Kim, \emph{The classification of $p$-divisible groups over $2$-adic discrete valuation rings}, Math. Res. Lett. \textbf{19} (2012), 121--141. Locators use arXiv:1007.1904v3. \url{https://arxiv.org/abs/1007.1904}.
\bibitem{ColmezFontaineLocal} P. Colmez and J.-M. Fontaine, \emph{Construction des repr\'esentations $p$-adiques semi-stables}, Invent. Math. \textbf{140} (2000), 1--43. Locators use the 37-page author version. \url{https://webusers.imj-prg.fr/~pierre.colmez/CF.pdf}.
\bibitem{ScholzeLocal} P. Scholze, \emph{$p$-adic Hodge theory for rigid-analytic varieties}, Forum Math. Pi \textbf{1} (2013), e1. Locators use the author version dated 3 November 2012. \url{https://people.mpim-bonn.mpg.de/scholze/pAdicHodgeTheory.pdf}.
\bibitem{ScholzeErratumLocal} P. Scholze, \emph{Erratum to ``$p$-adic Hodge theory for rigid-analytic varieties''}, 2016, three pages. \url{https://www.math.uni-bonn.de/people/scholze/pAdicHodgeErratum.pdf}.
\bibitem{TsujiLocal} T. Tsuji, \emph{Semi-stable conjecture of Fontaine--Jannsen: a survey}, Ast\'erisque \textbf{279} (2002), 323--370. \url{https://www.numdam.org/item/AST_2002__279__323_0.pdf}.
\bibitem{TsujiICMLocal} T. Tsuji, \emph{$p$-adic Hodge theory in the semi-stable reduction case}, Doc. Math., Extra Volume ICM 1998, II, 207--216. \url{https://ems.press/content/book-chapter-files/27150?nt=1}.
\bibitem{BMSLocal} B. Bhatt, M. Morrow and P. Scholze, \emph{Integral $p$-adic Hodge theory}, Publ. Math. IH\'ES \textbf{128} (2018), 219--397. Locators use arXiv:1602.03148v3, 15 January 2019. \url{https://arxiv.org/abs/1602.03148}.
\bibitem{PrismaticCrystalsLocal} B. Bhatt and P. Scholze, \emph{Prismatic $F$-crystals and crystalline Galois representations}, Camb. J. Math. \textbf{11} (2023), 507--562. Locators use the 38-page author version. \url{https://people.mpim-bonn.mpg.de/scholze/PrismaticCrystals.pdf}.
\bibitem{KedlayaLocal} K. S. Kedlaya, \emph{Finiteness of rigid cohomology with coefficients}, Duke Math. J. \textbf{134} (2006), 15--97. Locators use arXiv:math/0208027v6. \url{https://arxiv.org/abs/math/0208027}.
\bibitem{DegliseLocal} F. D\'eglise, \emph{Proper cohomological descent in rigid cohomology}, Rennes workshop notes, June 2005. \url{https://deglise.perso.math.cnrs.fr/docs/2005/proper.pdf}.
\bibitem{BerthelotLocal} P. Berthelot, \emph{Finitude et puret\'e cohomologique en cohomologie rigide}, Invent. Math. \textbf{128} (1997), 329--377. Locators use the author version. \url{https://wstein.org/people/berthelo/publis/Finitude.pdf}.
\bibitem{PrismsLocal} B. Bhatt and P. Scholze, \emph{Prisms and prismatic cohomology}, Ann. of Math. \textbf{196} (2022), 1135--1275. Locators use the 125-page author version. \url{https://www.math.uni-bonn.de/people/scholze/prisms.pdf}.
\bibitem{PstragowskiLocal} P. Pstr\k{a}gowski, \emph{$p$-adic analytic groups}, Harvard Math 291Y, Fall 2023, author lecture notes. \url{https://www.kurims.kyoto-u.ac.jp/~piotr/291y_notes.pdf}.
\bibitem{SerreLocal} J.-P. Serre, \emph{Groupes alg\'ebriques associ\'es aux modules de Hodge--Tate}, Ast\'erisque \textbf{65} (1979), 155--188. \url{https://www.numdam.org/item/AST_1979__65__155_0.pdf}.

\bibitem{DeligneHodgeCycles} P. Deligne, notes by J. S. Milne, \emph{Hodge cycles on abelian varieties}, in \emph{Hodge Cycles, Motives, and Shimura Varieties}, Lecture Notes in Math. 900, Springer, 1982, 9--100. Locators use the revised typescript of 1 October 2018. \url{https://www.jmilne.org/math/Documents/Deligne82.pdf}.
\bibitem{DeligneMilne} P. Deligne and J. S. Milne, \emph{Tannakian Categories}, in the same volume, 101--228. Locators use the corrected typescript of 11 February 2022. \url{https://www.jmilne.org/math/xnotes/tc2022.pdf}.
\bibitem{SGAComparison} M. Artin, \emph{Th\'eor\`eme de changement de base par un morphisme lisse, et applications}, SGA 4, tome 3, Expos\'e XVI, \S4, Th\'eor\`eme 4.1, original p.~233; retyped p.~149, PDF p.~156. \url{https://www.cmls.polytechnique.fr/perso/laszlo/sga4/SGA4-3/sga43.pdf}.
\bibitem{SerreAdelic} J.-P. Serre, \emph{Propri\'et\'es galoisiennes des points d'ordre fini des courbes elliptiques}, Invent. Math. \textbf{15} (1972), 259--331. \url{https://www.college-de-france.fr/media/jean-pierre-serre/UPL5874918517843398173_Serre_proprie_te_s_galoisiennes_des_courbes_elliptiques.pdf}.
\bibitem{FaltingsIsogeny} G. Faltings, \emph{Endlichkeitss\"atze f\"ur abelsche Variet\"aten \"uber Zahlk\"orpern}, Invent. Math. \textbf{73} (1983), 349--366. \url{https://math.uchicago.edu/~drinfeld/Deligne's_conjecture_Manin_conf/Faltings_argument/Faltings.pdf}.
\bibitem{LombardoProducts} D. Lombardo, \emph{An explicit open image theorem for products of elliptic curves}, author version, Theorem 1.1 and Remark 1.2, pp.~1--2. \url{https://web.dm.unipi.it/lombardo/files/ProductsEC.pdf}.
\bibitem{KSV} T. Kreutz, M. Shen and C. Vial, \emph{De Rham--Betti classes with coefficients}, arXiv:2206.08618v3, 21 January 2026. \url{https://arxiv.org/abs/2206.08618v3}.
\bibitem{KneserApprox} M. Kneser, \emph{Starke Approximation in algebraischen Gruppen. I}, J. reine angew. Math. \textbf{218} (1965), 190--203. Digitized volume, PDF pp.~195--199 for printed pp.~190--194. \url{https://gdz.sub.uni-goettingen.de/download/pdf/PPN243919689_0218/PPN243919689_0218.pdf}.
\bibitem{PlatonovApprox} V. P. Platonov, \emph{The problem of strong approximation and the Kneser--Tits conjecture for algebraic groups}, Math. USSR-Izv. \textbf{3} (1969), 1139--1147. \url{https://www.mathnet.ru/eng/im2220}.
\bibitem{PlatonovAddendum} V. P. Platonov, \emph{Addendum to the paper ``The problem of strong approximation and the Kneser--Tits conjecture for algebraic groups''}, Math. USSR-Izv. \textbf{4} (1970), 784--786. \url{https://www.mathnet.ru/eng/im2446}.
\bibitem{Masser} D. W. Masser, \emph{Elliptic Functions and Transcendence}, Lecture Notes in Math. 437, Springer, 1975. Appendix One, pp.~113--121. \url{https://link.springer.com/content/pdf/bbm:978-3-540-37410-7/1}.
\bibitem{FonsecaMatthes} T. J. Fonseca and N. Matthes, \emph{Towards algebraic iterated integrals for elliptic curves via the universal vectorial extension}, RIMS K\^oky\^uroku \textbf{2160} (2020), 114--125. Locators use arXiv:2009.10433v1. \url{https://arxiv.org/abs/2009.10433v1}.
\bibitem{BostKunnemann} J.-B. Bost and K. K\"unnemann, \emph{Hermitian vector bundles and extension groups on arithmetic schemes II. The arithmetic Atiyah extension}, Ast\'erisque \textbf{327} (2009), 361--424. \url{https://www.numdam.org/article/AST_2009__327__361_0.pdf}.
\bibitem{BostLeaves} J.-B. Bost, \emph{Algebraic leaves of algebraic foliations over number fields}, Publ. Math. IH\'ES \textbf{93} (2001), 161--221. \url{https://www.numdam.org/article/PMIHES_2001__93__161_0.pdf}.
\bibitem{WustholzAnalytic} G. W\"ustholz, \emph{Algebraische Punkte auf analytischen Untergruppen algebraischer Gruppen}, Max-Planck-Institut f\"ur Mathematik, Preprint 84-7 (1984). Hauptsatz, printed p.~3, PDF p.~4. \url{https://archive.mpim-bonn.mpg.de/id/eprint/3829/1/preprint_1984_7.pdf}.
\bibitem{Tang} Y. Tang, \emph{Cycles in the de Rham cohomology of abelian varieties over number fields}, Compos. Math. \textbf{154} (2018), 850--882. Locators use the author's 26-page preprint. \url{https://math.berkeley.edu/~ytang/dRT.pdf}.
\bibitem{KatzSerreTate} N. M. Katz, \emph{Serre--Tate local moduli}, in \emph{Algebraic Surfaces (Orsay, 1976--78)}, Lecture Notes in Math. 868, Springer, 1981, 138--202. \url{https://web.math.princeton.edu/~nmk/old/serretatelocmod.pdf}.
\bibitem{LiftingTraces} L. Panny, D. Robert and A. Sferlazza, \emph{Hensel-lifting black-box algorithms and fast trace computation for elliptic-curve endomorphisms}, author preprint dated 28 September 2026. \url{https://www.math.u-bordeaux.fr/~damienrobert/pro/publications/articles/lifting_traces.pdf}.
\bibitem{AGOR} C. Alfes, M. Griffin, K. Ono and L. Rolen, \emph{Weierstrass mock modular forms and elliptic curves}, Res. Number Theory \textbf{1} (2015), article 24. Locators use the 33-page author preprint. \url{https://uva.theopenscholar.com/files/ken-ono/files/1406.0443_8.pdf}.
\bibitem{Candelori} L. Candelori, \emph{Harmonic weak Maass forms of integral weight: a geometric approach}, Math. Ann. \textbf{360} (2014), 489--517. Locators use the 32-page author manuscript. \url{https://bpb-us-e1.wpmucdn.com/s.wayne.edu/dist/4/26/files/2018/09/HWMF-Final-1izyvzl.pdf}.
\bibitem{Guerzhoy} P. Guerzhoy, \emph{On Zagier's adele}, Res. Math. Sci. \textbf{1} (2014), article 7. \url{https://doi.org/10.1186/2197-9847-1-7}.
\bibitem{Coleman} R. F. Coleman, \emph{Classical and overconvergent modular forms}, Invent. Math. \textbf{124} (1996), 215--241. \url{https://wstein.org/people/coleman/papers/classical_and_overconvergent_forms.pdf}.
\bibitem{EmertonCM} M. Emerton, \emph{A $p$-adic variational Hodge conjecture and modular forms with complex multiplication}, undated author manuscript, 10 pp., Theorem (1.3), pp.~3--4. \url{https://www.math.uchicago.edu/~emerton/pdffiles/cm.pdf}.
\bibitem{CandeloriCastella} L. Candelori and F. Castella, \emph{A geometric perspective on $p$-adic properties of mock modular forms}, Res. Math. Sci. \textbf{4} (2017), article 5. \url{https://web.math.ucsb.edu/~castella/Ramanujan-print.pdf}.
\bibitem{BannaiKobayashi} K. Bannai and S. Kobayashi, \emph{Algebraic theta functions and the $p$-adic interpolation of Eisenstein--Kronecker numbers}, Duke Math. J. \textbf{153} (2010), 229--295. Locators use arXiv:math/0610163v4, 11 December 2007. \url{https://arxiv.org/abs/math/0610163v4}.
\bibitem{LamLitt} Y. H. J. Lam and D. Litt, \emph{Algebraicity and integrality of solutions to differential equations}, arXiv:2501.13175v1 (2025), Theorem 4.1.1, p.~21. \url{https://arxiv.org/abs/2501.13175v1}.

\bibitem{HuberWustholz} A. Huber and G. W\"ustholz, \emph{Transcendence and Linear Relations of $1$-Periods}, Cambridge Tracts in Math. 227, Cambridge University Press, 2022. Locators use arXiv:1805.10104v8, 21 April 2022; printed pp.~43--46 and 133, 136--138 correspond to PDF pp.~59--62 and 149, 152--154. \url{https://arxiv.org/abs/1805.10104v8}.
\bibitem{Movasati} H. Movasati, \emph{On elliptic modular foliations}, arXiv:0806.3926v1, 24 June 2008. Theorem 2, p.~3; Theorem 3 and Corollaries 1--2, pp.~13--14. \url{https://arxiv.org/abs/0806.3926v1}.
\bibitem{Herblot} M. Herblot, \emph{Algebraic points on meromorphic curves}, author manuscript dated 18 April 2012, 51 pp. \url{https://webusers.imj-prg.fr/~mathilde.herblot/publi/SL.pdf}.
\bibitem{MochizukiConnections} S. Mochizuki, \emph{Connections and Related Integral Structures on the Universal Extension of an Elliptic Curve}, author manuscript, May 2000, 80 pp. Theorem 1.3, pp.~13--14; \S2.2 and Theorem 2.2, pp.~16--19. \url{https://www.kurims.kyoto-u.ac.jp/~motizuki/Connections%20and%20Related%20Structures%20on%20the%20Universal%20Extension%20of%20an%20Elliptic%20Curve.PDF}.
\bibitem{WaldschmidtSubgroup} M. Waldschmidt, \emph{On the transcendence methods of Gel'fond and Schneider in several variables}, in \emph{New Advances in Transcendence Theory}, Cambridge University Press, 1988, 375--398. Theorem 4.1, printed pp.~385--386, PDF pp.~6--7. \url{https://webusers.imj-prg.fr/~michel.waldschmidt/articles/pdf/Durham1986TMGSSV.pdf}.
\bibitem{BostChambertLoir} J.-B. Bost and A. Chambert-Loir, \emph{Analytic curves in algebraic varieties over number fields}, arXiv:math/0702593v3, 2 October 2008, 55 pp. \url{https://arxiv.org/abs/math/0702593v3}.
\bibitem{AndreDwork} Y. Andr\'e, \emph{Sur la conjecture des $p$-courbures de Grothendieck--Katz et un probl\`eme de Dwork}, in \emph{Geometric Aspects of Dwork Theory}, vol.~I, de Gruyter, 2004, 55--112. \S5.1, pp.~70--71; \S5.3, p.~74; Definition 5.4.1 and Theorem 5.4.3, p.~76. The publisher preview contains these pages at PDF pp.~84--85, 88 and 90. \url{https://doi.org/10.1515/9783110198133.1.55}.

\bibitem{Arapura} D. Arapura, \emph{Complex Algebraic Varieties and their Cohomology}, author lecture notes, 16 July 2003, 143 PDF pp. Theorem~12.1.1, printed p.~97 (PDF p.~98); Theorem~12.3.1, printed p.~101 (PDF p.~102). \url{https://www.math.purdue.edu/~arapura/preprints/sheaves.pdf}.

\bibitem{Wibmer} M. Wibmer, \emph{A remark on torsors under affine group schemes}, Transformation Groups \textbf{30} (2025), 447--454; first published online 26 August 2022. Theorems~1.1--1.2, p.~447. \url{https://doi.org/10.1007/s00031-022-09767-z}.

\bibitem{StacksE} The Stacks Project Authors, \emph{The Stacks Project}, Section~10.45 and Lemmas~10.47.8, 29.35.3, 32.8.9, 32.10.1 and 39.8.2; Tags 05DU, 037P, 01V8, 0C0C, 01ZM and 047N. \url{https://stacks.math.columbia.edu}.

\bibitem{CoxClassFields} D. A. Cox, \emph{Primes of the Form $x^2+ny^2$: Fermat, Class Field Theory, and Complex Multiplication}, Wiley, 1989. Proposition~7.22, pp.~145--146; Lemma~9.3, pp.~180--181; Theorem~9.18, pp.~191--192; Theorem~11.1, p.~220. Scanned PDF pages are printed pages plus twelve. \url{https://www.math.utoronto.ca/~ila/Cox-Primes_of_the_form_x2%2Bny2.pdf}.

\bibitem{CampagnaPengo} F. Campagna and R. Pengo, \emph{How big is the image of the Galois representations attached to CM elliptic curves?}, Contemp. Math. \textbf{779} (2022), 41--56. Locators use arXiv:2201.04046v1, 11 January 2022; \S5, p.~11. \url{https://arxiv.org/abs/2201.04046v1}.

\bibitem{MilneCFT} J. S. Milne, \emph{Class Field Theory}, version 4.03, 6 August 2020. Chapter~V, \S3, Proposition~3.18 and the Artin-transfer diagram, printed p.~162, PDF p.~171. \url{https://www.jmilne.org/math/CourseNotes/CFT.pdf}.

\bibitem{MilneCMNotes} J. S. Milne, \emph{Complex Multiplication}, version 0.10, 14 July 2020. Proposition~1.18, p.~14; quasi-isogeny convention, p.~23; Example~2.11, p.~26; Proposition~3.17, p.~31; Proposition~7.9 and Corollary~7.10, pp.~54--55; Theorem~9.10 and Remark~9.11, p.~78. \url{https://www.jmilne.org/math/CourseNotes/CM.pdf}.

\bibitem{MilneAV} J. S. Milne, \emph{Abelian Varieties}, version 2.0, 16 March 2008. Chapter~I, Remark~8.12, printed p.~39, PDF p.~45; Chapter~IV, Theorem~1.1 and Lemma~2.1, printed pp.~131 and 135, PDF pp.~137 and 141. \url{https://www.jmilne.org/math/CourseNotes/AV.pdf}.

\bibitem{MasserWustholzIsogenies} D. W. Masser and G. W\"ustholz, \emph{Isogeny estimates for abelian varieties, and finiteness theorems}, Max-Planck-Institut f\"ur Mathematik, Preprint 91-28, 1991. Introduction, printed p.~1, PDF p.~3; \S7, printed pp.~10--11, PDF pp.~12--13. \url{https://archive.mpim-bonn.mpg.de/id/eprint/1497/1/preprint_1991_28.pdf}.

\end{thebibliography}
\end{document}
